\documentclass[UTF8,12pt,a4paper]{ctexbook}

\usepackage{amsmath,amssymb,mathtools}
\usepackage{geometry}
\usepackage{graphicx}
\usepackage{xcolor}
\usepackage{enumitem}
\usepackage{hyperref}
\usepackage{fancyhdr}
\usepackage[most]{tcolorbox}
\usepackage{float}

\geometry{left=22mm,right=22mm,top=22mm,bottom=24mm}
\hypersetup{colorlinks=true,linkcolor=blue!55!black,urlcolor=blue!55!black}
\setlist{nosep,leftmargin=2em}
\setlength{\headheight}{15pt}
\pagestyle{fancy}
\fancyhf{}
\lhead{27《没咋了》线性代数通关讲义做题本解析}
\rhead{\thepage}

\newenvironment{solutionblock}{\par\medskip\noindent}{\end{analysisbox}\par\medskip}

\newtcolorbox{questionbox}{
  enhanced,
  breakable,
  colback=green!4,
  colframe=green!45!black,
  boxrule=0.55pt,
  arc=1.2mm,
  left=2.2mm,
  right=2.2mm,
  top=1.8mm,
  bottom=1.8mm
}

\newtcolorbox{analysisbox}{
  enhanced,
  breakable,
  colback=blue!3,
  colframe=blue!45!black,
  boxrule=0.55pt,
  arc=1.2mm,
  left=2.2mm,
  right=2.2mm,
  top=1.8mm,
  bottom=1.8mm
}

\newcommand{\question}[1]{\begin{questionbox}\textbf{题目：}#1\end{questionbox}\par\smallskip}
\newcommand{\answer}[1]{\par\medskip\noindent\textbf{答案：}#1\par}
\newcommand{\analysis}[1]{\begin{analysisbox}\textbf{题目解析：}#1\par}
\newcommand{\method}[1]{\par\medskip\noindent\textbf{#1}\par}
\newcommand{\examnote}[1]{\par\medskip\noindent\textbf{考研提示：}#1\par}

\title{27《没咋了》线性代数通关讲义做题本\\详细解析与答案整理}
\author{整理：Codex}
\date{\today}

\begin{document}
\maketitle
\tableofcontents

\chapter*{使用说明}
\addcontentsline{toc}{chapter}{使用说明}
本工程把扫描版做题本整理为 Overleaf 可编译的 LaTeX 工程。本文档采用纯文字题面整理方式，不再用整页原题图片充当题目；每道例题均单独给出题目、答案、题目解析、关键计算和考研提示。一题适合多法处理时，解析中补充方法二。

\chapter{行列式}

\section{逐题解析}

\subsection{1.2 行列式的性质}

\begin{solutionblock}
\textbf{1.2 例 1}

\question{已知 \(\det(c,b,a)=-2\)，其中 \(c=(c_1,c_2,c_3)^T\)，\(b=(b_1,b_2,b_3)^T\)，\(a=(a_1,a_2,a_3)^T\)。求 \(\det(a,2c+a,-a-b)\)。}


\analysis{本题考查列变换对行列式的影响。把原行列式看成由三列向量
\[
C=(c,\ b,\ a)
\]
组成，其中 \(c=(c_1,c_2,c_3)^T,\ b=(b_1,b_2,b_3)^T,\ a=(a_1,a_2,a_3)^T\)。题设给出
\[
\det(c,b,a)=-2.
\]
目标行列式的三列分别为
\[
a,\qquad 2c+a,\qquad -a-b.
\]
它们相对于旧列组 \((c,b,a)\) 的坐标为
\[
a=C\begin{pmatrix}0\\0\\1\end{pmatrix},\quad
2c+a=C\begin{pmatrix}2\\0\\1\end{pmatrix},\quad
-a-b=C\begin{pmatrix}0\\-1\\-1\end{pmatrix}.
\]
因此新列组等于 \(CT\)，其中
\[
T=\begin{pmatrix}
0&2&0\\
0&0&-1\\
1&1&-1
\end{pmatrix}.
\]
于是
\[
\det(a,2c+a,-a-b)=\det(c,b,a)\det T.
\]
计算
\[
\det T=-2,
\]
故目标行列式为
\[
(-2)\cdot(-2)=4.
\]}
\answer{\(4\)}
\examnote{遇到“已知一个列组行列式，求另一个列组合行列式”，优先写出新列组相对于旧列组的变换矩阵，答案就是旧行列式乘以该变换矩阵的行列式。}
\end{solutionblock}

\begin{solutionblock}
\textbf{1.2 例 2}

\question{证明恒等式
\[
\begin{vmatrix}
ax+by&ay+bz&az+bx\\
ay+bz&az+bx&ax+by\\
az+bx&ax+by&ay+bz
\end{vmatrix}
=(a^3+b^3)
\begin{vmatrix}
x&y&z\\y&z&x\\z&x&y
\end{vmatrix}.
\]}


\analysis{记
\[
C=\begin{pmatrix}
x&y&z\\
y&z&x\\
z&x&y
\end{pmatrix},
\]
其三列为 \(u,v,w\)。题目左端矩阵的三列分别是
\[
au+bv,\qquad av+bw,\qquad aw+bu.
\]
所以左端矩阵可写成
\[
C\begin{pmatrix}
a&0&b\\
b&a&0\\
0&b&a
\end{pmatrix}.
\]
由行列式乘法公式，
\[
\det\bigl(au+bv,\ av+bw,\ aw+bu\bigr)
=\det C\cdot
\begin{vmatrix}
a&0&b\\
b&a&0\\
0&b&a
\end{vmatrix}.
\]
右侧小行列式按第一行展开：
\[
\begin{vmatrix}
a&0&b\\
b&a&0\\
0&b&a
\end{vmatrix}
=a\cdot a^2+b\cdot b^2=a^3+b^3.
\]
因此
\[
\begin{vmatrix}
ax+by&ay+bz&az+bx\\
ay+bz&az+bx&ax+by\\
az+bx&ax+by&ay+bz
\end{vmatrix}
=(a^3+b^3)
\begin{vmatrix}
x&y&z\\
y&z&x\\
z&x&y
\end{vmatrix}.
\]}
\answer{等式成立。}
\method{方法二（列向量法）}
直接把三列看成 \(u,v,w\) 的线性组合，列组合系数矩阵就是上面的 \(3\times3\) 循环矩阵。考场上这种方法比逐项展开稳定得多。
\end{solutionblock}

\subsection{1.3 行列式的计算}

\begin{solutionblock}
\textbf{1.3 例 1}

\question{设
\[
D=\begin{vmatrix}
3&-5&2&1\\
1&1&0&-5\\
-1&3&1&3\\
2&-4&-1&-3
\end{vmatrix},
\]
\(A_{ij}\) 为元素 \(a_{ij}\) 的代数余子式，\(M_{ij}\) 为余子式。求 \((1)\ A_{11}+A_{12}+A_{13}+A_{14}\)；\((2)\ A_{11}+M_{21}+A_{31}+M_{41}\)。}


\analysis{本题考查代数余子式展开的逆用法。若固定第 \(i\) 行，则
\[
b_1A_{i1}+b_2A_{i2}+\cdots+b_nA_{in}
\]
等于把原行列式第 \(i\) 行替换为 \((b_1,b_2,\ldots,b_n)\) 后得到的新行列式；若固定第 \(j\) 列，也有类似结论。解题时要先分清 \(A_{ij}\) 与 \(M_{ij}\)：代数余子式 \(A_{ij}=(-1)^{i+j}M_{ij}\)，余子式 \(M_{ij}\) 本身不带符号。第 (1) 问直接把第一行替换为全 \(1\)；第 (2) 问先把 \(M_{21},M_{41}\) 转成对应的代数余子式，再转化为第一列替换后的行列式。}

\method{第 (1) 问}
由第一行展开公式可知
\[
A_{11}+A_{12}+A_{13}+A_{14}
=
\begin{vmatrix}
1&1&1&1\\
1&1&0&-5\\
-1&3&1&3\\
2&-4&-1&-3
\end{vmatrix}.
\]
对第 2、3、4 行分别作
\[
R_2\leftarrow R_2-R_1,\quad
R_3\leftarrow R_3+R_1,\quad
R_4\leftarrow R_4-2R_1,
\]
得
\[
\begin{vmatrix}
1&1&1&1\\
0&0&-1&-6\\
0&4&2&4\\
0&-6&-3&-5
\end{vmatrix}
=
\begin{vmatrix}
0&-1&-6\\
4&2&4\\
-6&-3&-5
\end{vmatrix}=4.
\]
故第 (1) 问答案为 \(4\)。

\method{第 (2) 问}
注意余子式 \(M_{21}\) 与代数余子式的关系为
\[
A_{21}=(-1)^{2+1}M_{21}=-M_{21},\qquad M_{21}=-A_{21},
\]
同理 \(M_{41}=-A_{41}\)。因此
\[
A_{11}+M_{21}+A_{31}+M_{41}
=A_{11}-A_{21}+A_{31}-A_{41}.
\]
这是把原行列式第一列替换成 \((1,-1,1,-1)^T\) 后按第一列展开所得：
\[
\begin{vmatrix}
1&-5&2&1\\
-1&1&0&-5\\
1&3&1&3\\
-1&-4&-1&-3
\end{vmatrix}=0.
\]
\answer{第 (1) 问为 \(4\)；第 (2) 问为 \(0\)。}
\examnote{余子式 \(M_{ij}\) 不带符号，代数余子式 \(A_{ij}=(-1)^{i+j}M_{ij}\)。选择题里常故意混写二者，先转成同一种符号。}
\end{solutionblock}

\begin{solutionblock}
\textbf{1.3 例 2}

\question{设四阶行列式
\[
D=\begin{vmatrix}
0&1&0&0\\
0&0&2&0\\
0&0&0&3\\
4&0&0&0
\end{vmatrix},
\]
求第二行各元素代数余子式之和 \(A_{21}+A_{22}+A_{23}+A_{24}\)。}


\analysis{由第二行展开，
\[
A_{21}+A_{22}+A_{23}+A_{24}
\]
等于把原行列式第二行替换为全 \(1\) 后的行列式。即
\[
A_{21}+A_{22}+A_{23}+A_{24}
=
\begin{vmatrix}
0&1&0&0\\
1&1&1&1\\
0&0&0&3\\
4&0&0&0
\end{vmatrix}.
\]
按第三行展开：
\[
=-3
\begin{vmatrix}
0&1&0\\
1&1&1\\
4&0&0
\end{vmatrix}.
\]
再按第三行展开小行列式，
\[
\begin{vmatrix}
0&1&0\\
1&1&1\\
4&0&0
\end{vmatrix}
=4
\begin{vmatrix}
1&0\\
1&1
\end{vmatrix}=4.
\]
故
\[
A_{21}+A_{22}+A_{23}+A_{24}=-3\cdot4=-12.
\]}
\answer{\(-12\)}
\examnote{某一行代数余子式的和，直接替换该行后展开，不要逐个求四个余子式。}
\end{solutionblock}

\begin{solutionblock}
\textbf{1.3 例 3}

\question{计算 \(n\) 阶行列式
\[
D_n=\begin{vmatrix}
1&1&\cdots&1\\
x_1+1&x_2+1&\cdots&x_n+1\\
x_1^2+x_1&x_2^2+x_2&\cdots&x_n^2+x_n\\
\vdots&\vdots&&\vdots\\
x_1^{n-1}+x_1^{n-2}&x_2^{n-1}+x_2^{n-2}&\cdots&x_n^{n-1}+x_n^{n-2}
\end{vmatrix}.
\]}


\analysis{第 \(j\) 列由
\[
1,\quad x_j+1,\quad x_j^2+x_j,\quad \cdots,\quad x_j^{n-1}+x_j^{n-2}
\]
组成。把这些看成关于 \(x_j\) 的多项式：
\[
f_0(x)=1,\qquad f_k(x)=x^k+x^{k-1}\quad(k=1,\ldots,n-1).
\]
它们的次数依次为 \(0,1,\ldots,n-1\)，且每个最高次项系数都为 \(1\)。从单项式基
\[
1,x,x^2,\ldots,x^{n-1}
\]
变到 \(f_0,f_1,\ldots,f_{n-1}\) 的系数矩阵是对角线全为 \(1\) 的三角矩阵，行列式为 \(1\)。因此该行列式与标准范德蒙德行列式相同：
\[
D_n=
\begin{vmatrix}
1&1&\cdots&1\\
x_1&x_2&\cdots&x_n\\
x_1^2&x_2^2&\cdots&x_n^2\\
\vdots&\vdots&&\vdots\\
x_1^{n-1}&x_2^{n-1}&\cdots&x_n^{n-1}
\end{vmatrix}
=\prod_{1\le i<j\le n}(x_j-x_i).
\]}
\answer{\(\displaystyle D_n=\prod_{1\le i<j\le n}(x_j-x_i)\)}
\method{方法二（列因式视角）}
也可把每列看成多项式值列。只要第 \(k\) 行对应的是首项系数为 \(1\) 的 \(k-1\) 次多项式，行列式就仍是范德蒙德行列式。
\end{solutionblock}

\begin{solutionblock}
\textbf{1.3 例 4}

\question{设
\[
f(x)=\begin{vmatrix}
1&x&x^2&x^3\\
1&2&4&8\\
1&-1&1&-1\\
1&1&1&1
\end{vmatrix},
\]
求 \(f(x)\) 的零点个数。}


\analysis{该行列式是以
\[
t_1=x,\quad t_2=2,\quad t_3=-1,\quad t_4=1
\]
为节点的范德蒙德行列式：
\[
f(x)=
\begin{vmatrix}
1&x&x^2&x^3\\
1&2&4&8\\
1&-1&1&-1\\
1&1&1&1
\end{vmatrix}
=\prod_{1\le i<j\le4}(t_j-t_i).
\]
含 \(x\) 的因子为
\[
(2-x)(-1-x)(1-x),
\]
其余常数因子为
\[
(-1-2)(1-2)(1-(-1))=(-3)(-1)2=6.
\]
所以
\[
f(x)=6(2-x)(-1-x)(1-x).
\]
零点为
\[
x=2,\quad x=-1,\quad x=1,
\]
共有 \(3\) 个。}
\answer{\(3\)，选 D。}
\examnote{只要行列式中出现 \((1,t,t^2,\ldots)\) 结构，先想范德蒙德；求零点个数时不必完全展开。}
\end{solutionblock}

\subsection{1.4 特殊行列式的计算}

\begin{solutionblock}
\textbf{1.4 例 1}

\question{计算行列式
\[
\begin{vmatrix}
1&1&1&1\\
1&2&0&0\\
1&0&3&0\\
1&0&0&4
\end{vmatrix}.
\]}


\analysis{把行列式分块为
\[
\begin{vmatrix}
1&\mathbf 1^T\\
\mathbf 1&D
\end{vmatrix},
\qquad
D=\operatorname{diag}(2,3,4),\quad
\mathbf 1=(1,1,1)^T.
\]
由 Schur 补公式，
\[
\det
\begin{pmatrix}
1&\mathbf 1^T\\
\mathbf 1&D
\end{pmatrix}
=\det D\cdot\bigl(1-\mathbf 1^TD^{-1}\mathbf 1\bigr).
\]
这里
\[
\det D=24,\qquad
\mathbf 1^TD^{-1}\mathbf 1=\frac12+\frac13+\frac14=\frac{13}{12}.
\]
于是
\[
D=24\left(1-\frac{13}{12}\right)=-2.
\]}
\answer{\(-2\)}
\method{方法二（初等变换）}
也可把第 2、3、4 行分别减去第 1 行，再按第一列展开。分块法更快，且适合推广到更高阶。
\end{solutionblock}

\begin{solutionblock}
\textbf{1.4 例 2}

\question{计算行列式
\[
\begin{vmatrix}
1&1&0&0\\
1&2&1&0\\
1&0&3&1\\
1&0&0&4
\end{vmatrix}.
\]}


\analysis{对第 2、3、4 行作
\[
R_2\leftarrow R_2-R_1,\quad
R_3\leftarrow R_3-R_1,\quad
R_4\leftarrow R_4-R_1,
\]
得
\[
\begin{vmatrix}
1&1&0&0\\
0&1&1&0\\
0&-1&3&1\\
0&-1&0&4
\end{vmatrix}
=
\begin{vmatrix}
1&1&0\\
-1&3&1\\
-1&0&4
\end{vmatrix}.
\]
计算三阶行列式：
\[
1(3\cdot4-0\cdot1)-1((-1)\cdot4-(-1)\cdot1)=12-(-3)=15.
\]}
\answer{\(15\)}
\examnote{遇到第一列全为 \(1\) 的小阶行列式，常用“其余行减第一行”，把第一列化为 \(1,0,\ldots,0\)。}
\end{solutionblock}

\begin{solutionblock}
\textbf{1.4 例 3}

\question{计算行列式
\[
\begin{vmatrix}
1&1&1&1\\
2&0&0&1\\
3&0&1&0\\
4&1&0&0
\end{vmatrix}.
\]}


\analysis{直接计算该四阶行列式：
\[
D=
\begin{vmatrix}
1&1&1&1\\
2&0&0&1\\
3&0&1&0\\
4&1&0&0
\end{vmatrix}.
\]
按第一行展开也可，计算得
\[
D=8.
\]
若用初等变换，可先对第 2、3、4 行分别减去 \(2,3,4\) 倍第 1 行，再展开第一列，最终同样得到 \(8\)。}
\answer{\(8\)}
\examnote{小阶稀疏行列式不要硬套大公式；按零多的行列展开或先造零即可。}
\end{solutionblock}

\begin{solutionblock}
\textbf{1.4 例 4}

\question{计算 \(n\) 阶行列式
\[
D_n=\begin{vmatrix}
x&a&\cdots&a\\
a&x&\cdots&a\\
\vdots&\vdots&\ddots&\vdots\\
a&a&\cdots&x
\end{vmatrix}.
\]}


\analysis{该 \(n\) 阶矩阵对角线为 \(x\)，非对角线均为 \(a\)。可写成
\[
A=(x-a)I_n+aJ_n,
\]
其中 \(J_n\) 是全 \(1\) 矩阵。

全 \(1\) 向量 \(\mathbf 1=(1,\ldots,1)^T\) 是 \(J_n\) 的特征向量，且
\[
J_n\mathbf 1=n\mathbf 1.
\]
所以在 \(\mathbf 1\) 方向上，\(A\) 的特征值为
\[
(x-a)+an=x+(n-1)a.
\]
而任意满足 \(v_1+\cdots+v_n=0\) 的向量 \(v\)，都有 \(J_nv=0\)，故在这个 \(n-1\) 维子空间上，\(A\) 的特征值为
\[
x-a.
\]
因此
\[
D_n=(x-a)^{n-1}\bigl[x+(n-1)a\bigr].
\]}
\answer{\(\displaystyle (x-a)^{n-1}\bigl(x+(n-1)a\bigr)\)}
\method{方法二（初等变换）}
也可把所有行加到第一行，再用行列式性质提出 \(x+(n-1)a\)，随后作列减法化成上三角，得到同一结果。
\end{solutionblock}

\begin{solutionblock}
\textbf{1.4 例 5}

\question{计算行列式
\[
D=\begin{vmatrix}
1&-1&1&a-1\\
1&-1&a+1&-1\\
1&a-1&1&-1\\
a+1&-1&1&-1
\end{vmatrix}.
\]}


\analysis{观察到 \(a=0\) 时四行相同，行列式必为 \(0\)。用行变换把 \(a\) 因子显出来。对第 2、3、4 行作
\[
R_2\leftarrow R_2-R_1,\quad
R_3\leftarrow R_3-R_1,\quad
R_4\leftarrow R_4-R_1,
\]
得
\[
D=
\begin{vmatrix}
1&-1&1&a-1\\
0&0&a&-a\\
0&a&0&-a\\
a&0&0&-a
\end{vmatrix}
=a^3
\begin{vmatrix}
1&-1&1&a-1\\
0&0&1&-1\\
0&1&0&-1\\
1&0&0&-1
\end{vmatrix}.
\]
再对最后一个行列式作 \(R_1\leftarrow R_1-R_4\)，可得其值为 \(a\)。所以
\[
D=a^3\cdot a=a^4.
\]}
\answer{\(a^4\)}
\examnote{含参数行列式若在某个参数值下明显降秩，通常可以通过行差/列差把该参数因子逐个提出来。}
\end{solutionblock}

\begin{solutionblock}
\textbf{1.4 例 6}

\question{计算 \(n\) 阶行列式
\[
D_n=\begin{vmatrix}
a_1+b_1&a_2&\cdots&a_n\\
a_1&a_2+b_2&\cdots&a_n\\
\vdots&\vdots&\ddots&\vdots\\
a_1&a_2&\cdots&a_n+b_n
\end{vmatrix},\quad b_i\ne0.
\]}


\analysis{矩阵的第 \(i,j\) 元为
\[
a_j+\delta_{ij}b_i.
\]
因此
\[
A=\operatorname{diag}(b_1,\ldots,b_n)+\mathbf 1\,a^T,
\]
其中
\[
\mathbf 1=(1,\ldots,1)^T,\qquad a=(a_1,\ldots,a_n)^T.
\]
因为 \(b_i\ne0\)，可用矩阵行列式引理：
\[
\det(B+uv^T)=\det B\,(1+v^TB^{-1}u).
\]
令 \(B=\operatorname{diag}(b_1,\ldots,b_n)\), \(u=\mathbf 1\), \(v=a\)，则
\[
D_n=b_1b_2\cdots b_n
\left(1+\sum_{i=1}^n\frac{a_i}{b_i}\right).
\]
也可写成
\[
D_n=b_1b_2\cdots b_n
\left(1+\frac{a_1}{b_1}+\frac{a_2}{b_2}+\cdots+\frac{a_n}{b_n}\right).
\]}
\answer{\(\displaystyle b_1b_2\cdots b_n\left(1+\sum_{i=1}^n\frac{a_i}{b_i}\right)\)}
\examnote{“对角矩阵 + 秩一矩阵”是高频模型，结论要熟：\(\det(D+uv^T)=\det D(1+v^TD^{-1}u)\)。}
\end{solutionblock}

\begin{solutionblock}
\textbf{1.4 例 7}

\question{计算 \(n\) 阶行列式
\[
D_n=\begin{vmatrix}
a_1&-1&0&\cdots&0\\
a_2&x&-1&\cdots&0\\
\vdots&\vdots&\vdots&\ddots&\vdots\\
a_{n-1}&0&0&\cdots&-1\\
a_n&0&0&\cdots&x
\end{vmatrix}.
\]}


\analysis{按第一列展开。去掉第 \(i\) 行、第 \(1\) 列后，剩下的矩阵由主对角线上的 \(x\) 和上副对角线上的 \(-1\) 组成，其行列式贡献恰好为
\[
x^{n-i}.
\]
逐项展开得到
\[
D_n=a_1x^{n-1}+a_2x^{n-2}+\cdots+a_{n-1}x+a_n.
\]
用低阶情形核对：当 \(n=3\) 时
\[
\begin{vmatrix}
a_1&-1&0\\
a_2&x&-1\\
a_3&0&x
\end{vmatrix}
=a_1x^2+a_2x+a_3,
\]
与一般公式一致。}
\answer{\(\displaystyle D_n=\sum_{i=1}^n a_i x^{\,n-i}=a_1x^{n-1}+a_2x^{n-2}+\cdots+a_n\)}
\examnote{这种形状本质上是“多项式系数列 + 伴随型矩阵”，答案通常是一串按 \(x\) 降幂排列的多项式。}
\end{solutionblock}

\begin{solutionblock}
\textbf{1.4 例 8}

\question{设 \(D_n\) 为主对角线全为 \(2a\)、上副对角线全为 \(1\)、下副对角线全为 \(a^2\) 的 \(n\) 阶三对角行列式，证明 \(D_n=(n+1)a^n\)。}


\analysis{记该三对角行列式为 \(D_n\)。按最后一行或最后一列展开，可得递推关系
\[
D_n=2aD_{n-1}-a^2D_{n-2}\qquad(n\ge3).
\]
初值为
\[
D_1=2a,\qquad
D_2=
\begin{vmatrix}
2a&1\\
a^2&2a
\end{vmatrix}
3a^2.
\]
下面用归纳法。假设
\[
D_{n-1}=na^{n-1},\qquad D_{n-2}=(n-1)a^{n-2}.
\]
代入递推式：
\[
D_n=2a\cdot na^{n-1}-a^2\cdot (n-1)a^{n-2}
=(2n-n+1)a^n=(n+1)a^n.
\]
初值 \(D_1=2a=(1+1)a\)、\(D_2=3a^2=(2+1)a^2\) 也符合，故
\[
D_n=(n+1)a^n.
\]}
\answer{已证 \(\displaystyle D_n=(n+1)a^n\)。}
\method{方法二（特征根法）}
递推 \(D_n=2aD_{n-1}-a^2D_{n-2}\) 的特征方程为
\[
\lambda^2-2a\lambda+a^2=(\lambda-a)^2=0,
\]
故 \(D_n=(C_1+C_2n)a^n\)。由初值解得 \(D_n=(n+1)a^n\)。
\end{solutionblock}

\subsection{1.5 克拉默法则}

\begin{solutionblock}
\textbf{1.5 例 1}

\question{用克拉默法则求解线性方程组 \(Ax=b\)，其中
\[
A=\begin{pmatrix}
5&0&4&2\\
1&-1&2&1\\
4&1&2&0\\
1&1&1&1
\end{pmatrix},\qquad
b=\begin{pmatrix}3\\1\\1\\0\end{pmatrix}.
\]}


\analysis{系数矩阵与常数列为
\[
A=\begin{pmatrix}
5&0&4&2\\
1&-1&2&1\\
4&1&2&0\\
1&1&1&1
\end{pmatrix},
\qquad
b=\begin{pmatrix}3\\1\\1\\0\end{pmatrix}.
\]
先算系数行列式：
\[
\Delta=\det A=-7\ne0,
\]
所以线性方程组有唯一解，可以使用克拉默法则。

把 \(A\) 的第 \(i\) 列替换为 \(b\)，所得行列式分别为
\[
\Delta_1=-7,\qquad
\Delta_2=7,\qquad
\Delta_3=7,\qquad
\Delta_4=-7.
\]
由克拉默法则
\[
x_i=\frac{\Delta_i}{\Delta},
\]
得到
\[
x_1=1,\qquad x_2=-1,\qquad x_3=-1,\qquad x_4=1.
\]
代回检查：
\[
5\cdot1+4(-1)+2\cdot1=3,
\]
\[
1-(-1)+2(-1)+1=1,
\]
\[
4\cdot1+(-1)+2(-1)=1,
\]
\[
1+(-1)+(-1)+1=0.
\]
四个方程均满足。}
\answer{\(\displaystyle (x_1,x_2,x_3,x_4)=(1,-1,-1,1)\)}
\examnote{克拉默法则题的标准步骤是：先判 \(\Delta\ne0\)，再求 \(\Delta_i\)。若只写最终解，容易漏掉“唯一解”的依据。}
\end{solutionblock}

\chapter{矩阵}

\section{逐题解析}

\subsection{2.3 矩阵基本运算}

\begin{solutionblock}
\textbf{2.3 例 1}

\question{设
\[
A=\operatorname{diag}(a_{11},a_{22},\ldots,a_{nn}),\qquad
B=\operatorname{diag}(b_{11},b_{22},\ldots,b_{nn}).
\]
求乘积 \(AB\)。}


\analysis{设
\[
A=\operatorname{diag}(a_{11},a_{22},\ldots,a_{nn}),\qquad
B=\operatorname{diag}(b_{11},b_{22},\ldots,b_{nn}).
\]
乘积 \(AB=(c_{ij})\) 的元素为
\[
c_{ij}=\sum_{k=1}^n a_{ik}b_{kj}.
\]
因为 \(A,B\) 都是对角矩阵，只有 \(a_{ii}\) 与 \(b_{jj}\) 可能非零。若 \(i\ne j\)，求和式中每一项至少含有一个非对角元，所以 \(c_{ij}=0\)；若 \(i=j\)，只有 \(k=i\) 项保留，故
\[
c_{ii}=a_{ii}b_{ii}.
\]
因此 \(AB\) 仍为对角矩阵。}
\answer{\[
AB=\operatorname{diag}(a_{11}b_{11},a_{22}b_{22},\ldots,a_{nn}b_{nn}).
\]}
\examnote{对角矩阵相乘时，对角元逐项相乘；这是矩阵乘法中最常用的简化结构之一。}
\end{solutionblock}

\begin{solutionblock}
\textbf{2.3 例 2}

\question{设
\[
A=\begin{pmatrix}
a_{11}&a_{12}&a_{13}\\
a_{21}&a_{22}&a_{23}
\end{pmatrix},\quad
B=\begin{pmatrix}
k_1&0&0\\
0&k_2&0\\
0&0&k_3
\end{pmatrix},\quad
C=\begin{pmatrix}
l_1&0\\
0&l_2
\end{pmatrix}.
\]
求 \(AB\) 与 \(CA\)。}


\analysis{右乘对角矩阵会分别放大 \(A\) 的每一列，左乘对角矩阵会分别放大 \(A\) 的每一行。题中
\[
A=\begin{pmatrix}
a_{11}&a_{12}&a_{13}\\
a_{21}&a_{22}&a_{23}
\end{pmatrix},\quad
B=\begin{pmatrix}
k_1&0&0\\
0&k_2&0\\
0&0&k_3
\end{pmatrix},\quad
C=\begin{pmatrix}
l_1&0\\
0&l_2
\end{pmatrix}.
\]
于是
\[
AB=\begin{pmatrix}
k_1a_{11}&k_2a_{12}&k_3a_{13}\\
k_1a_{21}&k_2a_{22}&k_3a_{23}
\end{pmatrix},
\]
而
\[
CA=\begin{pmatrix}
l_1a_{11}&l_1a_{12}&l_1a_{13}\\
l_2a_{21}&l_2a_{22}&l_2a_{23}
\end{pmatrix}.
\]}
\answer{\[
AB=\begin{pmatrix}
k_1a_{11}&k_2a_{12}&k_3a_{13}\\
k_1a_{21}&k_2a_{22}&k_3a_{23}
\end{pmatrix},\qquad
CA=\begin{pmatrix}
l_1a_{11}&l_1a_{12}&l_1a_{13}\\
l_2a_{21}&l_2a_{22}&l_2a_{23}
\end{pmatrix}.
\]}
\method{方法二（列行视角）}
把 \(A\) 写成列向量组 \(A=(\alpha_1,\alpha_2,\alpha_3)\)，则
\[
AB=(k_1\alpha_1,k_2\alpha_2,k_3\alpha_3).
\]
把 \(A\) 写成行向量组，则 \(CA\) 是第 1 行乘 \(l_1\)，第 2 行乘 \(l_2\)。
\end{solutionblock}

\begin{solutionblock}
\textbf{2.3 例 3}

\question{设
\[
A=(1,2,-3),\qquad B=\begin{pmatrix}2\\1\\1\end{pmatrix}.
\]
分别求 \(AB\) 与 \(BA\)。}


\analysis{\(A\) 是 \(1\times3\) 行矩阵，\(B\) 是 \(3\times1\) 列矩阵，所以 \(AB\) 是 \(1\times1\) 矩阵，本质上是内积：
\[
AB=(1,2,-3)\begin{pmatrix}2\\1\\1\end{pmatrix}
=2+2-3=1.
\]
而 \(BA\) 是 \(3\times3\) 矩阵，本质上是列向量与行向量的外积：
\[
BA=\begin{pmatrix}2\\1\\1\end{pmatrix}(1,2,-3)
=\begin{pmatrix}
2&4&-6\\
1&2&-3\\
1&2&-3
\end{pmatrix}.
\]}
\answer{\[
AB=(1),\qquad
BA=\begin{pmatrix}
2&4&-6\\
1&2&-3\\
1&2&-3
\end{pmatrix}.
\]}
\examnote{矩阵乘法不满足交换律，甚至 \(AB\) 与 \(BA\) 的阶数都可能不同。}
\end{solutionblock}

\begin{solutionblock}
\textbf{2.3 例 4}

\question{设 \(A,B\) 为同阶方阵且 \(AB=O\)。判断下列命题中哪些一定成立：\(A+B=O\)、\(|A|=0\)、\(|B|=0\)、\(|A|+|B|=0\)、\(A=O\) 或 \(B=O\)、\(BA=O\)。}


\analysis{由 \(AB=O\) 得
\[
|A||B|=|AB|=|O|=0,
\]
所以必有 \(|A|=0\) 或 \(|B|=0\)，第 3 条一定成立。

其余命题均不一定成立。取
\[
A=\begin{pmatrix}1&0\\0&0\end{pmatrix},\qquad
B=\begin{pmatrix}0&0\\0&1\end{pmatrix},
\]
则 \(AB=O\)，但 \(A\ne O,B\ne O\)，且 \(A+B=E\ne O\)，因此第 1、2、5 条不成立，第 4 条也不成立，因为 \(|A|+|B|=0+0=0\) 在该例中成立但并非由题设推出；要否定第 4 条，可取 \(A=O,\ B=E\)，此时 \(AB=O\)，而 \(|A|+|B|=1\ne0\)。

第 6 条也不一定成立。取
\[
A=\begin{pmatrix}1&0\\0&0\end{pmatrix},\qquad
B=\begin{pmatrix}0&0\\1&0\end{pmatrix},
\]
则 \(AB=O\)，但
\[
BA=\begin{pmatrix}0&0\\1&0\end{pmatrix}\ne O.
\]}
\answer{只有第 3 条成立，选 B。}
\examnote{考研选择题中，遇到“必然成立”要用定理证明；其余选项只需要一个反例即可排除。}
\end{solutionblock}

\subsection{2.4 方阵的幂和方阵的多项式}

\begin{solutionblock}
\textbf{2.4 例 1}

\question{设
\[
A=\operatorname{diag}(\lambda_1,\lambda_2,\lambda_3),
\]
求 \(A^n\)。}


\analysis{对角矩阵的幂仍为对角矩阵，且对角元分别取幂。因为
\[
A=\operatorname{diag}(\lambda_1,\lambda_2,\lambda_3),
\]
所以反复相乘得到
\[
A^n=\operatorname{diag}(\lambda_1^n,\lambda_2^n,\lambda_3^n).
\]}
\answer{\[
A^n=\begin{pmatrix}
\lambda_1^n&0&0\\
0&\lambda_2^n&0\\
0&0&\lambda_3^n
\end{pmatrix}.
\]}
\end{solutionblock}

\begin{solutionblock}
\textbf{2.4 例 2}

\question{设
\[
A=\begin{pmatrix}
2&6&4\\
-1&-3&-2\\
2&6&4
\end{pmatrix}.
\]
求 \(A^n\)（\(n\ge3\)）。}


\analysis{直接计算一次平方：
\[
A^2=
\begin{pmatrix}
2&6&4\\
-1&-3&-2\\
2&6&4
\end{pmatrix}^2
=3
\begin{pmatrix}
2&6&4\\
-1&-3&-2\\
2&6&4
\end{pmatrix}
=3A.
\]
于是从 \(A^2=3A\) 出发递推：
\[
A^3=A\cdot A^2=3A^2=3^2A,\quad
A^4=A\cdot A^3=3^2A^2=3^3A.
\]
归纳可得 \(n\ge2\) 时
\[
A^n=3^{\,n-1}A.
\]}
\answer{\[
A^n=3^{\,n-1}
\begin{pmatrix}
2&6&4\\
-1&-3&-2\\
2&6&4
\end{pmatrix}\quad(n\ge3).
\]}
\examnote{矩阵幂题先算 \(A^2\) 或 \(A^3\)，若出现 \(A^2=kA\)、\(A^2=O\)、\(A^2=kE\) 等结构，后面全部变成递推。}
\end{solutionblock}

\begin{solutionblock}
\textbf{2.4 例 3}

\question{设 \(A\) 为四阶矩阵，其主对角线元素均为 \(1\)，非主对角线元素均为 \(-1\)。求 \(A^n\)。}


\analysis{该矩阵主对角线为 \(1\)，非主对角线为 \(-1\)。记 \(J\) 为 \(4\) 阶全 \(1\) 矩阵，则
\[
A=2E-J.
\]
又 \(J^2=4J\)，所以
\[
A^2=(2E-J)^2=4E-4J+J^2=4E.
\]
因此
\[
A^{2m}=(A^2)^m=4^mE=2^{2m}E,
\]
\[
A^{2m+1}=A(A^2)^m=4^mA=2^{2m}A.
\]}
\answer{\[
A^n=
\begin{cases}
2^nE, & n\text{ 为偶数},\\[2mm]
2^{\,n-1}A, & n\text{ 为奇数}.
\end{cases}
\]}
\method{方法二（特征结构）}
矩阵 \(J\) 在 \((1,1,1,1)^T\) 方向上的特征值为 \(4\)，在其正交补上的特征值为 \(0\)。所以 \(A=2E-J\) 的特征值只有 \(-2\) 与 \(2\)，从而 \(A^2=4E\)。
\end{solutionblock}

\begin{solutionblock}
\textbf{2.4 例 4}

\question{计算：
\[
A^2=
\begin{pmatrix}
1&0&1\\
0&2&0\\
1&0&1
\end{pmatrix}^2
=
\begin{pmatrix}
2&0&2\\
0&4&0\\
2&0&2
\end{pmatrix}
=2A.
\]}


\analysis{直接计算
\[
A^2=
\begin{pmatrix}
1&0&1\\
0&2&0\\
1&0&1
\end{pmatrix}^2
=
\begin{pmatrix}
2&0&2\\
0&4&0\\
2&0&2
\end{pmatrix}
=2A.
\]
于是
\[
A^n=2^{\,n-1}A\quad(n\ge1).
\]}
\answer{\[
A^n=2^{\,n-1}
\begin{pmatrix}
1&0&1\\
0&2&0\\
1&0&1
\end{pmatrix}\quad(n\ge3).
\]}
\end{solutionblock}

\begin{solutionblock}
\textbf{2.4 例 5}

\question{设
\[
A=\begin{pmatrix}
0&0&0\\
2&0&0\\
1&3&0
\end{pmatrix}.
\]
求 \(A^2\) 与 \(A^3\)。}


\analysis{题中
\[
A=\begin{pmatrix}
0&0&0\\
2&0&0\\
1&3&0
\end{pmatrix}.
\]
先求平方：
\[
A^2=
\begin{pmatrix}
0&0&0\\
0&0&0\\
6&0&0
\end{pmatrix}.
\]
再乘一次：
\[
A^3=A^2A=O.
\]
原因是 \(A\) 为严格下三角矩阵，三阶严格下三角矩阵至多三次幂就为零矩阵。}
\answer{\[
A^2=\begin{pmatrix}
0&0&0\\
0&0&0\\
6&0&0
\end{pmatrix},\qquad A^3=O.
\]}
\end{solutionblock}

\begin{solutionblock}
\textbf{2.4 例 6}

\question{设
\[
A=\begin{pmatrix}
1&0&0\\
1&1&0\\
0&1&1
\end{pmatrix}.
\]
求 \(A^n\)（\(n\ge3\)）。}


\analysis{把矩阵写成
\[
A=E+N,\qquad
N=\begin{pmatrix}
0&0&0\\
1&0&0\\
0&1&0
\end{pmatrix}.
\]
计算得
\[
N^2=\begin{pmatrix}
0&0&0\\
0&0&0\\
1&0&0
\end{pmatrix},\qquad N^3=O.
\]
由于 \(E\) 与 \(N\) 可交换，二项式公式可用于矩阵：
\[
A^n=(E+N)^n=E+nN+\binom n2N^2.
\]
代入 \(N,N^2\) 得
\[
A^n=
\begin{pmatrix}
1&0&0\\
n&1&0\\
\dfrac{n(n-1)}2&n&1
\end{pmatrix}.
\]}
\answer{\[
A^n=
\begin{pmatrix}
1&0&0\\
n&1&0\\
\dfrac{n(n-1)}2&n&1
\end{pmatrix}\quad(n\ge3).
\]}
\examnote{形如“单位矩阵加严格三角矩阵”的幂，优先使用二项式展开和幂零性。}
\end{solutionblock}

\subsection{2.5 转置矩阵}

\begin{solutionblock}
\textbf{2.5 例 1}

\question{设 \(A\) 为反对称矩阵，\(\alpha\) 为同维列向量。证明 \(\alpha^TA\alpha=0\)。}


\analysis{因 \(A\) 为反对称矩阵，所以 \(A^T=-A\)。令
\[
s=\alpha^TA\alpha.
\]
这是 \(1\times1\) 矩阵，即一个数，故 \(s=s^T\)。于是
\[
s^T=(\alpha^TA\alpha)^T=\alpha^TA^T\alpha=\alpha^T(-A)\alpha=-s.
\]
所以 \(s=-s\)，从而 \(s=0\)。}
\answer{\(\alpha^TA\alpha=0\)。}
\examnote{证明二次型为零时，常用“标量等于自身转置”这一招。}
\end{solutionblock}

\begin{solutionblock}
\textbf{2.5 例 2}

\question{设列向量 \(X\) 满足 \(X^TX=1\)，令
\[
H=E-2XX^T.
\]
证明 \(H\) 为对称矩阵且为正交矩阵。}


\analysis{已知 \(X^TX=1\)，且
\[
H=E-2XX^T.
\]
先证对称性：
\[
H^T=E^T-2(XX^T)^T=E-2XX^T=H.
\]
再证正交性。因为 \(H^T=H\)，只需算 \(H^2\)：
\[
H^2=(E-2XX^T)^2
=E-4XX^T+4X(X^TX)X^T.
\]
由 \(X^TX=1\)，得
\[
H^2=E-4XX^T+4XX^T=E.
\]
故
\[
HH^T=H^2=E.
\]}
\answer{\(H\) 为对称矩阵，且 \(HH^T=E\)。}
\method{方法二（Householder 矩阵）}
\(XX^T\) 是沿单位向量 \(X\) 方向的投影矩阵，\(E-2XX^T\) 表示关于与 \(X\) 垂直的超平面的反射。反射矩阵必为对称正交矩阵。
\end{solutionblock}

\begin{solutionblock}
\textbf{2.5 例 3}

\question{判断下列命题：\((1)\) 若 \(A^2=O\)，是否必有 \(A=O\)；\((2)\) 若 \(A^2=O\) 且 \(A^T=A\)，是否必有 \(A=O\)。若 \(A^2=O\) 且 \(A^T=A\)，则对任意列向量 \(x\)，
\[
\|Ax\|^2=(Ax)^T(Ax)=x^TA^TAx=x^TA^2x=0.
\]}


\analysis{第 (1) 问不正确。反例取
\[
A=\begin{pmatrix}0&1\\0&0\end{pmatrix}.
\]
则
\[
A^2=\begin{pmatrix}0&0\\0&0\end{pmatrix}=O,
\]
但 \(A\ne O\)。这说明一般矩阵可以非零但幂零。

第 (2) 问增加了对称条件。若 \(A^2=O\) 且 \(A^T=A\)，则对任意列向量 \(x\)，
\[
\|Ax\|^2=(Ax)^T(Ax)=x^TA^TAx=x^TA^2x=0.
\]
所以对任意 \(x\) 都有 \(Ax=0\)，从而 \(A=O\)。}
\answer{第 (1) 问不正确；第 (2) 问结论成立，\(A=O\)。}
\examnote{实对称矩阵适合配合 \(\|Ax\|^2=x^TA^TAx\) 使用，这是“平方为零推出矩阵为零”的关键补充条件。}
\end{solutionblock}

\subsection{2.6 伴随矩阵}

\begin{solutionblock}
\textbf{2.6 例 1}

\question{设
\[
A=\begin{pmatrix}
0&1&3\\
1&-1&0\\
-1&2&1
\end{pmatrix}.
\]
求伴随矩阵 \(A^*\)。}


\analysis{对
\[
A=\begin{pmatrix}
0&1&3\\
1&-1&0\\
-1&2&1
\end{pmatrix}
\]
逐项求代数余子式并转置，得到伴随矩阵。也可先算 \(|A|=2\)，再由 \(A^*=|A|A^{-1}\) 得到。直接计算结果为
\[
A^*=
\begin{pmatrix}
-1&5&3\\
-1&3&3\\
1&-1&-1
\end{pmatrix}.
\]
验算：
\[
AA^*=A^*A=2E.
\]}
\answer{\[
A^*=
\begin{pmatrix}
-1&5&3\\
-1&3&3\\
1&-1&-1
\end{pmatrix}.
\]}
\end{solutionblock}

\begin{solutionblock}
\textbf{2.6 例 2}

\question{设 \(A\) 为三阶矩阵，\(|A|=2\)。交换 \(A\) 的第 \(1\) 行与第 \(3\) 行得到 \(B\)，求 \(|BA^*|\)。}


\analysis{交换 \(A\) 的第 1 行与第 3 行得 \(B\)，因此
\[
|B|=-|A|=-2.
\]
又 \(A\) 为三阶矩阵，所以
\[
|A^*|=|A|^{3-1}=|A|^2=4.
\]
于是
\[
|BA^*|=|B||A^*|=(-2)\cdot4=-8.
\]}
\answer{\(-8\)}
\examnote{\(n\) 阶矩阵满足 \(|A^*|=|A|^{n-1}\)。三阶时就是 \(|A^*|=|A|^2\)。}
\end{solutionblock}

\begin{solutionblock}
\textbf{2.6 例 3}

\question{设三阶矩阵 \(A\) 满足 \(A^*=A^T\)，且第一行三个元素相等并均为正数。求第一行元素的值，并判断正确选项。}


\analysis{题设 \(A^*=A^T\)。由伴随矩阵基本恒等式
\[
AA^*=|A|E
\]
得
\[
AA^T=|A|E.
\]
这说明 \(A\) 的各行两两正交，且每一行的长度平方都等于 \(|A|\)。

另一方面，
\[
|A^*|=|A|^{2},\qquad |A^T|=|A|.
\]
由 \(A^*=A^T\) 得
\[
|A|^2=|A|.
\]
若 \(|A|=0\)，则 \(AA^T=O\)，从而 \(A=O\)，这与第一行有三个相等的正数矛盾。因此 \(|A|=1\)。

设
\[
a_{11}=a_{12}=a_{13}=t,\qquad t>0.
\]
由第一行长度平方为 \(1\)，
\[
t^2+t^2+t^2=1,
\]
故
\[
t=\frac1{\sqrt3}=\frac{\sqrt3}{3}.
\]}
\answer{\(\dfrac{\sqrt3}{3}\)，选 A。}
\end{solutionblock}

\begin{solutionblock}
\textbf{2.6 例 4}

\question{设 \(A\) 为方阵且 \(A^n=E\)。求 \((A^*)^n\)。}


\analysis{由 \(A^n=E\) 可知 \(A\) 可逆，且 \(A^{-n}=E\)。伴随矩阵满足
\[
A^*=|A|A^{-1}.
\]
于是
\[
(A^*)^n=|A|^nA^{-n}.
\]
又
\[
|A|^n=|A^n|=|E|=1,
\]
所以
\[
(A^*)^n=E.
\]}
\answer{\(E\)}
\end{solutionblock}

\subsection{2.7 逆矩阵}

\begin{solutionblock}
\textbf{2.7 例 1}

\question{设方阵 \(A\) 满足
\[
A^2-3A-10E=O.
\]
求 \(A^{-1}\) 与 \((A-4E)^{-1}\)。}


\analysis{由
\[
A^2-3A-10E=O
\]
可整理为
\[
A(A-3E)=10E.
\]
因此 \(A\) 可逆，且
\[
A^{-1}=\frac1{10}(A-3E).
\]

再求 \(A-4E\) 的逆。利用原方程 \(A^2=3A+10E\)，有
\[
(A-4E)(A+E)=A^2-3A-4E=(3A+10E)-3A-4E=6E.
\]
所以 \(A-4E\) 可逆，且
\[
(A-4E)^{-1}=\frac16(A+E).
\]}
\answer{\[
A^{-1}=\frac1{10}(A-3E),\qquad
(A-4E)^{-1}=\frac16(A+E).
\]}
\examnote{矩阵满足多项式方程时，把等式凑成 \(AB=E\) 的形式即可证明可逆并读出逆矩阵。}
\end{solutionblock}

\begin{solutionblock}
\textbf{2.7 例 2}

\question{设同阶方阵 \(A,B\) 满足 \(A+B=AB\)。证明：\((1)\ A-E\) 可逆；\((2)\ AB=BA\)。}


\analysis{由 \(A+B=AB\) 得
\[
AB-A-B=O.
\]
两边加 \(E\)，得
\[
(A-E)(B-E)=E.
\]
所以 \(A-E\) 可逆，并且
\[
(A-E)^{-1}=B-E.
\]
既然 \(B-E\) 是 \(A-E\) 的逆矩阵，则也有
\[
(B-E)(A-E)=E.
\]
展开得
\[
BA-B-A=O,
\]
即
\[
BA=A+B.
\]
而题设给出 \(AB=A+B\)，故
\[
AB=BA.
\]}
\answer{\(A-E\) 可逆，且 \(AB=BA\)。}
\end{solutionblock}

\begin{solutionblock}
\textbf{2.7 例 3}

\question{设同阶方阵 \(A,B,C\) 满足 \(ABC=E\)。判断 \(BCA,CAB,ACB,BAC,CBA\) 中哪些必等于 \(E\)。}


\analysis{由 \(ABC=E\) 可知 \(A,B,C\) 均可逆。对等式作循环移动：
\[
ABC=E.
\]
左乘 \(A^{-1}\)、右乘 \(A\)，得
\[
BCA=E.
\]
再左乘 \(B^{-1}\)、右乘 \(B\)，得
\[
CAB=E.
\]
因此循环排列 \(ABC,BCA,CAB\) 都等于 \(E\)。

非循环排列一般不必等于 \(E\)。例如矩阵乘法通常不交换，所以 \(ACB,BAC,CBA\) 没有必然结论。}
\answer{必等于 \(E\) 的是 \(BCA\) 与 \(CAB\)。}
\examnote{从 \(ABC=E\) 出发，安全变形是“整体循环移动”，不是任意交换两个因子。}
\end{solutionblock}

\begin{solutionblock}
\textbf{2.7 例 4}

\question{设矩阵如下，求其逆矩阵：
\[
|A|=
\begin{vmatrix}
1&2&3\\
2&2&1\\
3&4&3
\end{vmatrix}=2\ne0,
\]}


\analysis{先算行列式：
\[
|A|=
\begin{vmatrix}
1&2&3\\
2&2&1\\
3&4&3
\end{vmatrix}=2\ne0,
\]
故 \(A\) 可逆。用伴随矩阵公式
\[
A^{-1}=\frac1{|A|}A^*
\]
计算，得到
\[
A^{-1}=
\begin{pmatrix}
1&3&-2\\
-\dfrac32&-3&\dfrac52\\
1&1&-1
\end{pmatrix}.
\]}
\answer{\[
A^{-1}=
\begin{pmatrix}
1&3&-2\\
-\dfrac32&-3&\dfrac52\\
1&1&-1
\end{pmatrix}.
\]}
\method{方法二（初等行变换）}
也可对 \((A\mid E)\) 作初等行变换，将左半边化为 \(E\)，右半边即为 \(A^{-1}\)。数值逆矩阵题考场上更推荐初等行变换，符号逆矩阵题更推荐伴随矩阵公式。
\end{solutionblock}

\begin{solutionblock}
\textbf{2.7 例 5}

\question{设 \(A,B\) 为同阶方阵。判断下列命题中正确的一项：
\begin{enumerate}
\item 若 \(A\) 或 \(B\) 可逆，则 \(AB\) 可逆；
\item 若 \(A\) 或 \(B\) 不可逆，则 \(AB\) 不可逆；
\item 若 \(A,B\) 均可逆，则 \(A+B\) 可逆；
\item 若 \(A,B\) 均不可逆，则 \(A+B\) 不可逆。
\end{enumerate}}


\analysis{逐项判断：

A 项错误。若只有 \(A\) 或 \(B\) 中一个可逆，另一个可能不可逆，则 \(AB\) 仍不可逆。例如 \(A=E,B=O\)。

B 项正确。若 \(A\) 或 \(B\) 不可逆，则 \(|A|=0\) 或 \(|B|=0\)，所以
\[
|AB|=|A||B|=0,
\]
故 \(AB\) 不可逆。

C 项错误。即使 \(A,B\) 均可逆，也可能 \(A+B\) 不可逆，例如 \(B=-A\)。

D 项错误。即使 \(A,B\) 均不可逆，也可能 \(A+B\) 可逆，例如
\[
A=\begin{pmatrix}1&0\\0&0\end{pmatrix},\quad
B=\begin{pmatrix}0&0\\0&1\end{pmatrix},
\]
则 \(A+B=E\)。}
\answer{选 B。}
\end{solutionblock}

\begin{solutionblock}
\textbf{2.7 例 6}

\question{设
\[
A=\begin{pmatrix}
1&0&1\\
0&3&0\\
0&0&1
\end{pmatrix}.
\]
求 \((A+2E)^{-1}(A^2-4E)\)。}


\analysis{注意
\[
A^2-4E=(A-2E)(A+2E).
\]
因为
\[
A+2E=
\begin{pmatrix}
3&0&1\\
0&5&0\\
0&0&3
\end{pmatrix}
\]
为上三角矩阵，主对角元均非零，所以可逆。于是
\[
(A+2E)^{-1}(A^2-4E)
=(A+2E)^{-1}(A-2E)(A+2E).
\]
由于 \(A-2E\) 与 \(A+2E\) 都是 \(A\) 的多项式，二者可交换，所以
\[
(A+2E)^{-1}(A^2-4E)=A-2E.
\]
代入
\[
A=\begin{pmatrix}
1&0&1\\
0&3&0\\
0&0&1
\end{pmatrix},
\]
得
\[
A-2E=\begin{pmatrix}
-1&0&1\\
0&1&0\\
0&0&-1
\end{pmatrix}.
\]}
\answer{\[
\begin{pmatrix}
-1&0&1\\
0&1&0\\
0&0&-1
\end{pmatrix}.
\]}
\examnote{矩阵多项式彼此可交换。若能因式分解，就不要硬算逆矩阵。}
\end{solutionblock}

\begin{solutionblock}
\textbf{2.7 例 7}

\question{设 \(A\) 为三阶可逆矩阵，\(|A|=\dfrac12\)。求行列式
\[
\left|(3A)^{-1}-2A^*\right|.
\]}


\analysis{三阶矩阵满足
\[
A^*=|A|A^{-1}.
\]
已知 \(|A|=\dfrac12\)，所以
\[
A^*=\frac12A^{-1}.
\]
又
\[
(3A)^{-1}=\frac13A^{-1}.
\]
因此
\[
(3A)^{-1}-2A^*
=\frac13A^{-1}-A^{-1}
=-\frac23A^{-1}.
\]
取行列式。因 \(A\) 为三阶，
\[
\left|-\frac23A^{-1}\right|
=\left(-\frac23\right)^3|A^{-1}|
=-\frac8{27}\cdot\frac1{|A|}
=-\frac8{27}\cdot2
=-\frac{16}{27}.
\]}
\answer{\(-\dfrac{16}{27}\)}
\end{solutionblock}

\begin{solutionblock}
\textbf{2.7 例 8}

\question{设 \(A,B\) 为同阶可逆矩阵，且 \(|A|=1,\ |B|=2,\ |A+B|=2\)。求
\[
\left|\left(A^{-1}+B^{-1}\right)^{-1}\right|.
\]}


\analysis{先把和式化成含 \(A+B\) 的形式：
\[
A^{-1}+B^{-1}=A^{-1}(A+B)B^{-1}.
\]
因此
\[
|A^{-1}+B^{-1}|
=|A^{-1}|\,|A+B|\,|B^{-1}|
=\frac{|A+B|}{|A||B|}.
\]
代入 \(|A|=1, |B|=2, |A+B|=2\)，得
\[
|A^{-1}+B^{-1}|=\frac2{1\cdot2}=1.
\]
于是
\[
\left|\left(A^{-1}+B^{-1}\right)^{-1}\right|
=\frac1{|A^{-1}+B^{-1}|}=1.
\]}
\answer{\(1\)，选 A。}
\end{solutionblock}

\begin{solutionblock}
\textbf{2.7 例 9}

\question{设 \(A\) 为 \(n\ge3\) 阶可逆矩阵。判断关于伴随矩阵、逆矩阵和数乘矩阵伴随的四个结论是否正确。}


\analysis{逐条判断。设 \(A\) 为 \(n\ge3\) 阶可逆矩阵。

第 1 条：
\[
A^*=|A|A^{-1},
\]
所以
\[
(A^*)^{-1}=\frac1{|A|}A.
\]
而
\[
(A^{-1})^*=|A^{-1}|(A^{-1})^{-1}=\frac1{|A|}A.
\]
故第 1 条正确。

第 2 条：
\[
(kA)^*=|kA|(kA)^{-1}=k^n|A|\cdot\frac1kA^{-1}=k^{n-1}A^*
\quad(k\ne0),
\]
正确。

第 3 条：代数余子式与转置相容，等价地
\[
(A^*)^T=(A^T)^*,
\]
正确。

第 4 条：
\[
(A^*)^*=|A^*|(A^*)^{-1}
=|A|^{n-1}\cdot\frac1{|A|}A
=|A|^{n-2}A,
\]
正确。}
\answer{四条均正确，选 D。}
\end{solutionblock}

\begin{solutionblock}
\textbf{2.7 例 10}

\question{设
\[
A=\begin{pmatrix}
a&1&0\\
1&a&-1\\
0&1&a
\end{pmatrix}.
\]
(1) 若 \(A^3=O\)，求 \(a\)；(2) 在此条件下，解矩阵方程
\[
X-XA^2-AX+AXA^2=E.
\]}


\analysis{第 (1) 问，设
\[
A=\begin{pmatrix}
a&1&0\\
1&a&-1\\
0&1&a
\end{pmatrix}.
\]
直接计算
\[
A^3=
\begin{pmatrix}
a^3+3a&3a^2&-3a\\
3a^2&a^3&-3a^2\\
3a&3a^2&a^3-3a
\end{pmatrix}.
\]
若 \(A^3=O\)，由 \((1,3)\) 元 \(-3a=0\) 得 \(a=0\)。代回后确有 \(A^3=O\)，故
\[
a=0.
\]

第 (2) 问，此时
\[
A=\begin{pmatrix}
0&1&0\\
1&0&-1\\
0&1&0
\end{pmatrix},\qquad A^3=O.
\]
原方程
\[
X-XA^2-AX+AXA^2=E
\]
可因式分解为
\[
(E-A)X(E-A^2)=E.
\]
由于 \(A^3=O\)，
\[
(E-A)^{-1}=E+A+A^2,\qquad (E-A^2)^{-1}=E+A^2.
\]
因此
\[
X=(E-A)^{-1}(E-A^2)^{-1}
=(E+A+A^2)(E+A^2)=E+A+2A^2.
\]
计算得
\[
X=
\begin{pmatrix}
3&1&-2\\
1&1&-1\\
2&1&-1
\end{pmatrix}.
\]}
\answer{\[
a=0,\qquad
X=\begin{pmatrix}
3&1&-2\\
1&1&-1\\
2&1&-1
\end{pmatrix}.
\]}
\examnote{非交换矩阵方程中，因式分解要保持左右顺序。本题是 \((E-A)X(E-A^2)\)，不能随意把 \(X\) 挪到右侧。}
\end{solutionblock}

\subsection{2.8 矩阵的分块}

\begin{solutionblock}
\textbf{2.8 例 1}

\question{计算矩阵乘积
\[
AB=\begin{pmatrix}
1&0&1&2\\
0&1&3&4\\
0&0&1&0\\
0&0&0&-1
\end{pmatrix}
\begin{pmatrix}
1&0&0&0\\
0&1&0&0\\
2&1&1&0\\
0&3&0&-1
\end{pmatrix}.
\]}


\analysis{直接按行列相乘，或把矩阵按块看成上三角分块矩阵。计算得
\[
AB=
\begin{pmatrix}
1&0&1&2\\
0&1&3&4\\
0&0&1&0\\
0&0&0&-1
\end{pmatrix}
\begin{pmatrix}
1&0&0&0\\
0&1&0&0\\
2&1&1&0\\
0&3&0&-1
\end{pmatrix}
=
\begin{pmatrix}
3&7&1&-2\\
6&16&3&-4\\
2&1&1&0\\
0&-3&0&1
\end{pmatrix}.
\]}
\answer{\[
AB=
\begin{pmatrix}
3&7&1&-2\\
6&16&3&-4\\
2&1&1&0\\
0&-3&0&1
\end{pmatrix}.
\]}
\end{solutionblock}

\begin{solutionblock}
\textbf{2.8 例 2}

\question{设
\[
A=\begin{pmatrix}
1&0&0&0&0\\
0&1&0&0&0\\
-1&2&1&0&0\\
1&1&0&1&0\\
0&1&0&0&1
\end{pmatrix},\quad
B=\begin{pmatrix}
1&0&0&0\\
-1&0&0&0\\
0&1&2&-1\\
0&2&1&2\\
0&1&2&1
\end{pmatrix}.
\]
求 \(AB\)。}


\analysis{矩阵 \(A\) 是 \(5\times5\)，\(B\) 是 \(5\times4\)，所以 \(AB\) 为 \(5\times4\)。按行乘列计算：
\[
AB=
\begin{pmatrix}
1&0&0&0&0\\
0&1&0&0&0\\
-1&2&1&0&0\\
1&1&0&1&0\\
0&1&0&0&1
\end{pmatrix}
\begin{pmatrix}
1&0&0&0\\
-1&0&0&0\\
0&1&2&-1\\
0&2&1&2\\
0&1&2&1
\end{pmatrix}
=
\begin{pmatrix}
1&0&0&0\\
-1&0&0&0\\
-3&1&2&-1\\
0&2&1&2\\
-1&1&2&1
\end{pmatrix}.
\]}
\answer{\[
AB=
\begin{pmatrix}
1&0&0&0\\
-1&0&0&0\\
-3&1&2&-1\\
0&2&1&2\\
-1&1&2&1
\end{pmatrix}.
\]}
\method{方法二（初等矩阵视角）}
\(A\) 的前两行等于单位矩阵前两行，后面几行是对 \(B\) 的行线性组合。因此 \(AB\) 的每一行可以直接看成 \(B\) 的若干行组合，计算量比逐元素相乘少。
\end{solutionblock}

\begin{solutionblock}
\textbf{2.8 例 3}

\question{设分块矩阵
\[
A=\begin{pmatrix}C&O\\O&D\end{pmatrix},\quad
C=\begin{pmatrix}1&1\\2&2\end{pmatrix},\quad
D=\begin{pmatrix}0&1&0\\0&0&1\\0&0&0\end{pmatrix}.
\]
求 \(A^{100}\)。}


\analysis{把 \(A\) 分成左上 \(2\times2\) 块与右下 \(3\times3\) 块：
\[
A=\begin{pmatrix}
C&O\\
O&D
\end{pmatrix},
\quad
C=\begin{pmatrix}1&1\\2&2\end{pmatrix},\quad
D=\begin{pmatrix}0&1&0\\0&0&1\\0&0&0\end{pmatrix}.
\]
其中
\[
C^2=3C,
\]
故
\[
C^{100}=3^{99}C.
\]
而 \(D\) 是三阶严格上三角矩阵，\(D^3=O\)，所以 \(D^{100}=O\)。因此
\[
A^{100}=
\begin{pmatrix}
3^{99}C&O\\
O&O
\end{pmatrix}.
\]}
\answer{\[
A^{100}=
\begin{pmatrix}
3^{99}&3^{99}&0&0&0\\
2\cdot3^{99}&2\cdot3^{99}&0&0&0\\
0&0&0&0&0\\
0&0&0&0&0\\
0&0&0&0&0
\end{pmatrix}.
\]}
\end{solutionblock}

\begin{solutionblock}
\textbf{2.8 例 4}

\question{设
\[
A=\begin{pmatrix}
0&2&0\\
0&0&-3\\
4&0&0
\end{pmatrix}.
\]
求 \(A^{-1}\)。}


\analysis{矩阵
\[
A=\begin{pmatrix}
0&2&0\\
0&0&-3\\
4&0&0
\end{pmatrix}
\]
每一行每一列只有一个非零元。设
\[
A^{-1}=\begin{pmatrix}
x_{11}&x_{12}&x_{13}\\
x_{21}&x_{22}&x_{23}\\
x_{31}&x_{32}&x_{33}
\end{pmatrix},
\]
由 \(AA^{-1}=E\) 解得
\[
A^{-1}=
\begin{pmatrix}
0&0&\dfrac14\\
\dfrac12&0&0\\
0&-\dfrac13&0
\end{pmatrix}.
\]}
\answer{\[
A^{-1}=
\begin{pmatrix}
0&0&\dfrac14\\
\dfrac12&0&0\\
0&-\dfrac13&0
\end{pmatrix}.
\]}
\method{方法二（置换乘对角）}
该矩阵等于“置换矩阵 \(\times\) 对角矩阵”的形式，逆矩阵就是反向置换并把非零系数取倒数。
\end{solutionblock}

\begin{solutionblock}
\textbf{2.8 例 5}

\question{设 \(A,B\) 为二阶可逆矩阵，\(|A|=2,\ |B|=3\)，令
\[
M=\begin{pmatrix}O&A\\B&O\end{pmatrix}.
\]
求 \(M^*\)，并判断正确选项。}


\analysis{记
\[
M=\begin{pmatrix}O&A\\B&O\end{pmatrix}.
\]
由于 \(|A|=2, |B|=3\)，\(A,B\) 均可逆。直接验证
\[
M^{-1}=\begin{pmatrix}O&B^{-1}\\A^{-1}&O\end{pmatrix}.
\]
又
\[
|M|=|A||B|=6
\]
（这里 \(A,B\) 都是二阶矩阵，块交换符号为正）。因此
\[
M^*=|M|M^{-1}
=6\begin{pmatrix}O&B^{-1}\\A^{-1}&O\end{pmatrix}.
\]
由二阶矩阵
\[
A^*=|A|A^{-1}=2A^{-1},\qquad
B^*=|B|B^{-1}=3B^{-1},
\]
得
\[
6B^{-1}=2B^*,\qquad 6A^{-1}=3A^*.
\]}
\answer{\[
M^*=\begin{pmatrix}O&2B^*\\3A^*&O\end{pmatrix},
\]
选 B。}
\end{solutionblock}

\begin{solutionblock}
\textbf{2.8 例 6}

\question{设 \[
P=\begin{pmatrix}
E&0\\
-\alpha^TA^*&|A|
\end{pmatrix},\qquad
Q=\begin{pmatrix}
A&\alpha\\
\alpha^T&b
\end{pmatrix}.
\]
其中 \(A^*\) 为 \(A\) 的伴随矩阵。计算分块乘法：
\[
PQ=
\begin{pmatrix}
A&\alpha\\
-\alpha^TA^*A+|A|\alpha^T&
-\alpha^TA^*\alpha+|A|b
\end{pmatrix}.
\]}


\analysis{题中
\[
P=\begin{pmatrix}
E&0\\
-\alpha^TA^*&|A|
\end{pmatrix},\qquad
Q=\begin{pmatrix}
A&\alpha\\
\alpha^T&b
\end{pmatrix}.
\]
其中 \(A^*\) 为 \(A\) 的伴随矩阵。计算分块乘法：
\[
PQ=
\begin{pmatrix}
A&\alpha\\
-\alpha^TA^*A+|A|\alpha^T&
-\alpha^TA^*\alpha+|A|b
\end{pmatrix}.
\]
因为
\[
A^*A=|A|E,\qquad A^*=|A|A^{-1},
\]
所以
\[
-\alpha^TA^*A+|A|\alpha^T=0,
\]
并且
\[
-\alpha^TA^*\alpha+|A|b
=|A|\bigl(b-\alpha^TA^{-1}\alpha\bigr).
\]
故
\[
PQ=
\begin{pmatrix}
A&\alpha\\
0&|A|\bigl(b-\alpha^TA^{-1}\alpha\bigr)
\end{pmatrix}.
\]
这是上三角分块矩阵，因此
\[
|PQ|=|A|\cdot |A|\bigl(b-\alpha^TA^{-1}\alpha\bigr).
\]
又 \(|P|=|A|\ne0\)，故
\[
|Q|=|A|\bigl(b-\alpha^TA^{-1}\alpha\bigr).
\]
于是 \(Q\) 可逆的充要条件为
\[
b-\alpha^TA^{-1}\alpha\ne0.
\]}
\answer{\[
PQ=
\begin{pmatrix}
A&\alpha\\
0&|A|\bigl(b-\alpha^TA^{-1}\alpha\bigr)
\end{pmatrix},
\qquad
Q\text{ 可逆}\Longleftrightarrow \alpha^TA^{-1}\alpha\ne b.
\]}
\examnote{这是 Schur 补思想的典型考法。分块矩阵可逆性常通过块消元化成三角分块矩阵。}
\end{solutionblock}

\subsection{2.9 初等变换与初等矩阵}

\begin{solutionblock}
\textbf{2.9 例 1}

\question{将矩阵
\[
\begin{pmatrix}
1&-1&0&2\\
0&2&2&-1\\
0&0&3&-1\\
0&0&0&0
\end{pmatrix}
\]
化为标准形。}


\analysis{原矩阵已是阶梯形：
\[
\begin{pmatrix}
1&-1&0&2\\
0&2&2&-1\\
0&0&3&-1\\
0&0&0&0
\end{pmatrix}.
\]
前三行有非零主元，所以秩为 \(3\)。任意矩阵经初等行、列变换可化为标准形
\[
\begin{pmatrix}
E_r&O\\
O&O
\end{pmatrix}.
\]
这里 \(r=3\)，故标准形为
\[
\begin{pmatrix}
1&0&0&0\\
0&1&0&0\\
0&0&1&0\\
0&0&0&0
\end{pmatrix}.
\]}
\answer{\[
\begin{pmatrix}
1&0&0&0\\
0&1&0&0\\
0&0&1&0\\
0&0&0&0
\end{pmatrix}.
\]}
\end{solutionblock}

\begin{solutionblock}
\textbf{2.9 例 2}

\question{设矩阵 \(A=(\alpha_1,\alpha_2,\alpha_3)\)，矩阵
\[
B=(3\alpha_1,\ 3\alpha_1-2\alpha_2,\ \alpha_3).
\]
判断 \(A\) 与 \(B\) 是否等价。}


\analysis{记 \(A\) 的三列分别为 \(\alpha_1,\alpha_2,\alpha_3\)。由题中 \(B\) 的表达式可见
\[
B=(3\alpha_1,\ 3\alpha_1-2\alpha_2,\ \alpha_3).
\]
因此
\[
B=A
\begin{pmatrix}
3&3&0\\
0&-2&0\\
0&0&1
\end{pmatrix}.
\]
右侧乘的矩阵行列式为
\[
3\cdot(-2)\cdot1=-6\ne0,
\]
所以它是可逆矩阵。于是 \(B\) 可由 \(A\) 经过有限次初等列变换得到，故 \(A\) 与 \(B\) 等价。}
\answer{矩阵 \(A\) 与 \(B\) 等价。}
\examnote{证明等价常用形式：若 \(B=PAQ\)，其中 \(P,Q\) 可逆，则 \(A\) 与 \(B\) 等价。本题只需要右乘一个可逆矩阵。}
\end{solutionblock}

\begin{solutionblock}
\textbf{2.9 例 3}

\question{设
\[
P=\begin{pmatrix}0&1&0\\1&0&0\\0&0&1\end{pmatrix},\quad
M=\begin{pmatrix}1&2&3\\4&5&6\\7&8&9\end{pmatrix},\quad
Q=\begin{pmatrix}1&-1&0\\0&1&0\\0&0&1\end{pmatrix}.
\]
求 \(PMQ\)。}


\analysis{设
\[
P=\begin{pmatrix}0&1&0\\1&0&0\\0&0&1\end{pmatrix},\quad
M=\begin{pmatrix}1&2&3\\4&5&6\\7&8&9\end{pmatrix},\quad
Q=\begin{pmatrix}1&-1&0\\0&1&0\\0&0&1\end{pmatrix}.
\]
矩阵 \(P\) 表示交换第 1、2 行，故 \(P^2=E\)，从而
\[
P^{100}=E.
\]
于是原式为 \(MQ\)。右乘 \(Q\) 表示列变换：
\[
C_1\text{ 不变},\qquad C_2\leftarrow -C_1+C_2,\qquad C_3\text{ 不变}.
\]
所以
\[
MQ=
\begin{pmatrix}
1&1&3\\
4&1&6\\
7&1&9
\end{pmatrix}.
\]}
\answer{\[
\begin{pmatrix}
1&1&3\\
4&1&6\\
7&1&9
\end{pmatrix}.
\]}
\end{solutionblock}

\begin{solutionblock}
\textbf{2.9 例 4}

\question{设三阶可逆矩阵 \(A=(\alpha_1,\alpha_2,\alpha_3)\)，\(|A|=2\)，且
\[
B=(\alpha_1,\alpha_2,3\alpha_1+\alpha_3).
\]
求 \(A^*B\)。}


\analysis{设 \(A\) 的列向量为 \(\alpha_1,\alpha_2,\alpha_3\)。题中
\[
B=(\alpha_1,\alpha_2,3\alpha_1+\alpha_3),
\]
因此
\[
B=A
\begin{pmatrix}
1&0&3\\
0&1&0\\
0&0&1
\end{pmatrix}.
\]
由伴随矩阵性质
\[
A^*A=|A|E=2E,
\]
故
\[
A^*B=A^*A
\begin{pmatrix}
1&0&3\\
0&1&0\\
0&0&1
\end{pmatrix}
=2
\begin{pmatrix}
1&0&3\\
0&1&0\\
0&0&1
\end{pmatrix}.
\]}
\answer{\[
A^*B=
\begin{pmatrix}
2&0&6\\
0&2&0\\
0&0&2
\end{pmatrix}.
\]}
\end{solutionblock}

\begin{solutionblock}
\textbf{2.9 例 5}

\question{设 \(P_1\) 表示交换矩阵的第 \(1,2\) 行，\(P_2\) 表示交换矩阵的第 \(2,3\) 列。若 \(P_1AP_2=E\)，判断 \(A^*\) 的正确表达式。}


\analysis{左乘 \(P_1\) 表示交换矩阵的第 1、2 行，右乘 \(P_2\) 表示交换矩阵的第 2、3 列。题设等价于
\[
P_1AP_2=E.
\]
因此
\[
A^{-1}=P_2P_1.
\]
又 \(P_1,P_2\) 均为一次交换对应的初等矩阵，行列式均为 \(-1\)，所以
\[
|A|=\frac{|E|}{|P_1||P_2|}=1.
\]
于是
\[
A^*=|A|A^{-1}=P_2P_1.
\]}
\answer{\(A^*=P_2P_1\)，选 B。}
\examnote{行变换左乘初等矩阵，列变换右乘初等矩阵。左右顺序是本题的关键。}
\end{solutionblock}

\begin{solutionblock}
\textbf{2.9 例 6}

\question{设
\[
A=\begin{pmatrix}
0&1&3\\
1&-1&0\\
-1&2&1
\end{pmatrix}.
\]
求 \(A^{-1}\)。}


\analysis{题中矩阵
\[
A=\begin{pmatrix}
0&1&3\\
1&-1&0\\
-1&2&1
\end{pmatrix}
\]
与 2.6 例 1 相同，已知
\[
|A|=2,\qquad
A^*=
\begin{pmatrix}
-1&5&3\\
-1&3&3\\
1&-1&-1
\end{pmatrix}.
\]
因此
\[
A^{-1}=\frac1{|A|}A^*
=\frac12
\begin{pmatrix}
-1&5&3\\
-1&3&3\\
1&-1&-1
\end{pmatrix}.
\]}
\answer{\[
A^{-1}=
\begin{pmatrix}
-\dfrac12&\dfrac52&\dfrac32\\
-\dfrac12&\dfrac32&\dfrac32\\
\dfrac12&-\dfrac12&-\dfrac12
\end{pmatrix}.
\]}
\end{solutionblock}

\subsection{2.10 矩阵的秩}

\begin{solutionblock}
\textbf{2.10 例 1}

\question{设 \(a_1,a_2,\ldots,a_5\) 两两不同，矩阵 \(A\) 的第 \(i\) 列为
\[
(1,a_i,a_i^2,a_i^3,(a_i+1)^3)^T\quad(i=1,2,\ldots,5).
\]
求 \(r(A)\)。}


\analysis{矩阵的前四行分别是
\[
1,\quad a_i,\quad a_i^2,\quad a_i^3\quad(i=1,\ldots,5),
\]
它们组成 \(4\times5\) 的 Vandermonde 型矩阵。由于 \(a_i\ne a_j\)，任取四列可得到四阶 Vandermonde 行列式，非零，因此前四行秩为 \(4\)。

最后一行为
\[
(a_i+1)^3=a_i^3+3a_i^2+3a_i+1,
\]
即它等于“第 4 行 \(+3\) 倍第 3 行 \(+3\) 倍第 2 行 \(+\) 第 1 行”。所以最后一行不增加秩。}
\answer{\(r(A)=4\)。}
\examnote{看到 \((a_i+1)^3\)，要立刻展开成 \(a_i^3,a_i^2,a_i,1\) 的线性组合。}
\end{solutionblock}

\begin{solutionblock}
\textbf{2.10 例 2}

\question{设
\[
A=\begin{pmatrix}
1&2&-1&1\\
2&0&a&0\\
0&-4&5&-2
\end{pmatrix}.
\]
若 \(r(A)=2\)，求 \(a\)。}


\analysis{矩阵为
\[
A=\begin{pmatrix}
1&2&-1&1\\
2&0&a&0\\
0&-4&5&-2
\end{pmatrix}.
\]
它是 \(3\times4\) 矩阵。若秩为 \(2\)，则所有三阶子式均为 \(0\)，且至少有一个二阶子式非零。

取由第 \(1,2,3\) 列组成的三阶子式：
\[
\begin{vmatrix}
1&2&-1\\
2&0&a\\
0&-4&5
\end{vmatrix}
=4(a-3).
\]
取由第 \(1,3,4\) 列组成的三阶子式：
\[
\begin{vmatrix}
1&-1&1\\
2&a&0\\
0&5&-2
\end{vmatrix}
=-2(a-3).
\]
其余三阶子式也含因子 \(a-3\) 或恒为 \(0\)。因此所有三阶子式为 \(0\) 当且仅当
\[
a=3.
\]
此时前两列构成的二阶子式
\[
\begin{vmatrix}1&2\\2&0\end{vmatrix}=-4\ne0,
\]
故秩确为 \(2\)。}
\answer{\(a=3\)。}
\end{solutionblock}

\begin{solutionblock}
\textbf{2.10 例 3}

\question{设矩阵 \(A,B\) 满足 \(B\) 可逆，其中
\[
B=\begin{pmatrix}
1&0&4\\
0&2&0\\
2&0&3
\end{pmatrix}.
\]
根据题给矩阵 \(A\) 与 \(B\)，求 \(r(AB)\)。}


\analysis{先看
\[
B=\begin{pmatrix}
1&0&4\\
0&2&0\\
2&0&3
\end{pmatrix}.
\]
其行列式为
\[
|B|=2\begin{vmatrix}1&4\\2&3\end{vmatrix}
=2(3-8)=-10\ne0.
\]
所以 \(B\) 可逆。右乘可逆矩阵不改变矩阵的秩，因此
\[
r(AB)=r(A)=2.
\]}
\answer{\(r(AB)=2\)。}
\end{solutionblock}

\begin{solutionblock}
\textbf{2.10 例 4}

\question{设 \(A\) 为 \(m\times n\) 矩阵，\(r(A)=r<m<n\)。判断关于子式、行列式和标准形的四个命题中正确命题的个数。}


\analysis{已知 \(A\) 为 \(m\times n\) 矩阵，\(r(A)=r<m<n\)。

第 1 条错误。秩为 \(r\) 不代表任意 \(r-1\) 阶子式都非零。

第 2 条错误。\(A\) 不是方阵，\(|A|\) 本身没有定义，不能说 \(|A|=0\)。

第 3 条错误。秩为 \(r\) 只说明至少存在一个 \(r\) 阶子式非零，不是任意 \(r\) 阶子式都非零。

第 4 条错误。仅经过初等行变换，一般只能化成行阶梯形，不能保证右上角也化为零块；要化成标准形 \(\begin{pmatrix}E_r&O\\O&O\end{pmatrix}\)，通常需要初等行变换和初等列变换共同作用。}
\answer{正确的个数为 \(0\)，选 A。}
\end{solutionblock}

\begin{solutionblock}
\textbf{2.10 例 5}

\question{设 \(A\) 为 \(m\times n\) 矩阵，\(B\) 为 \(n\times m\) 矩阵。判断关于 \(|AB|\) 与 \(m,n\) 大小关系的选项是否正确。}


\analysis{题面四个选项中，A、B 两项在原 PDF 中均印成“当 \(m>n\) 时，必有 \(|AB|\ne0\)”。按这一题面判断，四项均不正确。

理由如下：\(A\) 为 \(m\times n\)，\(B\) 为 \(n\times m\)，所以 \(AB\) 为 \(m\times m\)。若 \(m>n\)，则
\[
r(AB)\le n<m,
\]
故
\[
|AB|=0,
\]
不是 \(|AB|\ne0\)。

若 \(m<n\)，则 \(|AB|\) 不一定为 \(0\)，也不一定非零。例如取
\[
A=(I_m\ O),\qquad B=\begin{pmatrix}I_m\\O\end{pmatrix},
\]
则 \(AB=I_m\)，行列式非零；而取 \(A=O\) 时，\(AB=O\)，行列式为零。}
\answer{按扫描件现有题面：无正确选项。若 B 项原意为“当 \(m>n\) 时，必有 \(|AB|=0\)”，则应选 B。}
\examnote{这里疑似原题选项有印刷错误。考研做题时若发现选项自相矛盾，应回到秩不等式 \(r(AB)\le\min\{r(A),r(B)\}\) 判断。}
\end{solutionblock}

\begin{solutionblock}
\textbf{2.10 例 6}

\question{设 \(A\) 为 \(m\times n\) 矩阵，\(B\) 为 \(n\times m\) 矩阵，且 \(AB=E_m\)。判断关于 \(r(A)\)、\(r(B)\) 及 \(m,n\) 的正确结论。}


\analysis{由 \(AB=E_m\) 得
\[
r(AB)=r(E_m)=m.
\]
又
\[
r(AB)\le r(A)\le m,\qquad r(AB)\le r(B)\le m,
\]
所以只能有
\[
r(A)=m,\qquad r(B)=m.
\]}
\answer{选 A。}
\examnote{\(AB=E_m\) 说明 \(A\) 有右逆、\(B\) 有左逆，并立刻推出 \(m\le n\) 且二者秩均为 \(m\)。}
\end{solutionblock}

\begin{solutionblock}
\textbf{2.10 例 7}

\question{设非零三阶矩阵 \(P\) 与
\[
Q=\begin{pmatrix}
1&2&3\\
2&4&k\\
3&6&9
\end{pmatrix}
\]
满足 \(PQ=O\)。判断 \(k\) 与 \(r(P)\) 的正确关系。}


\analysis{设
\[
Q=\begin{pmatrix}
1&2&3\\
2&4&k\\
3&6&9
\end{pmatrix}.
\]
其第 2 列恒为第 1 列的 \(2\) 倍，所以 \(r(Q)\le2\)。当 \(k=6\) 时，第 3 列也为第 1 列的 \(3\) 倍，故
\[
r(Q)=1.
\]
当 \(k\ne6\) 时，第 1 列与第 3 列不成比例，故
\[
r(Q)=2.
\]

题设 \(P\ne O\) 且 \(PQ=O\)。这说明 \(Q\) 的列空间包含在 \(P\) 的零空间中，即
\[
r(Q)\le n-r(P)=3-r(P).
\]
若 \(k\ne6\)，则 \(r(Q)=2\)，所以 \(r(P)\le1\)。又 \(P\ne O\)，故
\[
r(P)=1.
\]
若 \(k=6\)，只能推出 \(r(P)\le2\)，不一定为 \(1\) 或 \(2\) 中某一个固定值。}
\answer{选 C。}
\end{solutionblock}

\begin{solutionblock}
\textbf{2.10 例 8}

\question{设
\[
A=\begin{pmatrix}
a&b&b\\
b&a&b\\
b&b&a
\end{pmatrix},
\]
且 \(r(A^*)=1\)。判断 \(a,b\) 应满足的正确条件。}


\analysis{矩阵
\[
A=\begin{pmatrix}
a&b&b\\
b&a&b\\
b&b&a
\end{pmatrix}
=(a-b)E+bJ,
\]
其中 \(J\) 为三阶全 \(1\) 矩阵。它的特征值为
\[
a-b\quad(\text{二重}),\qquad a+2b\quad(\text{一重}).
\]
对于三阶矩阵，有
\[
r(A^*)=
\begin{cases}
3,& r(A)=3,\\
1,& r(A)=2,\\
0,& r(A)\le1.
\end{cases}
\]
题设 \(r(A^*)=1\)，所以 \(r(A)=2\)。这要求恰有一个特征值为零。

若 \(a=b\)，则 \(a-b=0\) 为二重零特征值，且 \(a,b\ne0\) 时 \(a+2b=3a\ne0\)，此时 \(r(A)=1\)，不符合。

因此必须
\[
a+2b=0,\qquad a-b\ne0.
\]
由于 \(b\ne0\)，由 \(a=-2b\) 自动有 \(a-b=-3b\ne0\)。}
\answer{\(a+2b=0\)，选 C。}
\end{solutionblock}

\begin{solutionblock}
\textbf{2.10 例 9}

\question{设 \(A\) 为 \(n\) 阶幂等矩阵，即 \(A^2=A\)。证明并求与 \(r(A)\)、\(r(E-A)\) 有关的结论。}


\analysis{第 (1) 问，若 \(A^2=A\)，则 \(A\) 为幂等矩阵。对任意向量 \(x\)，有分解
\[
x=Ax+(E-A)x.
\]
其中 \(Ax\in \operatorname{Im}A\)，\((E-A)x\in\operatorname{Im}(E-A)\)，故
\[
\mathbb{R}^n=\operatorname{Im}A+\operatorname{Im}(E-A).
\]
若 \(y\) 同时属于二者，则 \(y=Au=(E-A)v\)。于是
\[
Ay=A^2u=Au=y,
\]
同时
\[
Ay=A(E-A)v=(A-A^2)v=O.
\]
故 \(y=0\)。所以这是直和分解：
\[
\mathbb{R}^n=\operatorname{Im}A\oplus\operatorname{Im}(E-A).
\]
取维数得
\[
r(A)+r(E-A)=n.
\]
由于 \(r(E-A)=r(A-E)\)，故
\[
r(A)+r(A-E)=n.
\]

第 (2) 问，若 \(A^2=E\)，则
\[
(A+E)(A-E)=A^2-E=O.
\]
对任意向量 \(x\)，
\[
x=\frac12(A+E)x-\frac12(A-E)x,
\]
故
\[
\mathbb{R}^n=\operatorname{Im}(A+E)+\operatorname{Im}(A-E).
\]
若 \(y\) 同时属于 \(\operatorname{Im}(A+E)\) 和 \(\operatorname{Im}(A-E)\)，则
\[
(A-E)y=O,\qquad (A+E)y=O.
\]
两式相减或相加可得 \(y=0\)，所以二者也是直和。于是
\[
r(A+E)+r(A-E)=n.
\]}
\answer{两个等式均成立：
\[
r(A)+r(A-E)=n,\qquad r(A+E)+r(A-E)=n.
\]}
\examnote{幂等矩阵和对合矩阵的秩等式，本质都是把空间分解成两个互补子空间。考研证明题中这种“像空间直和”写法很稳。}
\end{solutionblock}

\chapter{线性方程组（初步认识）}

\section{逐题解析}

\subsection{3.2 解方程原理及高斯消元法}

\begin{solutionblock}
\textbf{3.2 例 1}

\question{求解线性方程组
\[
\begin{cases}
3x_1+x_2+x_3=2,\\
x_1+2x_2+2x_3=-1,\\
-x_1+2x_2+x_3=-5.
\end{cases}
\]}


\analysis{方程组为
\[
\begin{cases}
3x_1+x_2+x_3=2,\\
x_1+2x_2+2x_3=-1,\\
-x_1+2x_2+x_3=-5.
\end{cases}
\]
写出增广矩阵并作初等行变换：
\[
\left(
\begin{array}{ccc|c}
3&1&1&2\\
1&2&2&-1\\
-1&2&1&-5
\end{array}
\right)
\sim
\left(
\begin{array}{ccc|c}
1&0&0&1\\
0&1&0&-3\\
0&0&1&2
\end{array}
\right).
\]
所以
\[
x_1=1,\qquad x_2=-3,\qquad x_3=2.
\]}
\answer{\((x_1,x_2,x_3)=(1,-3,2)\)。}
\end{solutionblock}

\begin{solutionblock}
\textbf{3.2 例 2}

\question{判断增广矩阵
\[
\left(
\begin{array}{ccc|c}
1&-2&3&4\\
4&-2&-4&1\\
3&0&-7&5
\end{array}
\right)
\]
对应的线性方程组是否有解。}


\analysis{增广矩阵为
\[
\left(
\begin{array}{ccc|c}
1&-2&3&4\\
4&-2&-4&1\\
3&0&-7&5
\end{array}
\right).
\]
行化简得
\[
\left(
\begin{array}{ccc|c}
1&0&-\dfrac73&0\\
0&1&-\dfrac83&0\\
0&0&0&1
\end{array}
\right).
\]
最后一行表示
\[
0=1,
\]
矛盾，因此方程组无解。}
\answer{无解。}
\examnote{增广矩阵出现 \((0,\ldots,0\mid c)\)、\(c\ne0\) 的行，立即判定无解。}
\end{solutionblock}

\begin{solutionblock}
\textbf{3.2 例 3}

\question{求齐次线性方程组 \(Ax=0\) 的通解，其中
\[
A=\begin{pmatrix}
2&1&-1\\
1&-1&1\\
4&5&-5
\end{pmatrix}.
\]}


\analysis{齐次方程组的系数矩阵为
\[
A=\begin{pmatrix}
2&1&-1\\
1&-1&1\\
4&5&-5
\end{pmatrix}.
\]
行化简：
\[
\left(
\begin{array}{ccc|c}
2&1&-1&0\\
1&-1&1&0\\
4&5&-5&0
\end{array}
\right)
\sim
\left(
\begin{array}{ccc|c}
1&0&0&0\\
0&1&-1&0\\
0&0&0&0
\end{array}
\right).
\]
于是
\[
x_1=0,\qquad x_2=x_3.
\]
令 \(x_3=t\)，则 \(x_2=t\)。}
\answer{\[
(x_1,x_2,x_3)=(0,t,t),\qquad t\in\mathbb{R}.
\]}
\end{solutionblock}

\begin{solutionblock}
\textbf{3.2 例 4}

\question{求非齐次线性方程组 \(Ax=b\) 的通解，其增广矩阵为
\[
\left(
\begin{array}{ccc|c}
2&1&-1&1\\
1&-1&1&2\\
4&5&-5&-1
\end{array}
\right).
\]}


\analysis{本题与例 3 的系数矩阵相同，但常数项不同。增广矩阵化简为
\[
\left(
\begin{array}{ccc|c}
2&1&-1&1\\
1&-1&1&2\\
4&5&-5&-1
\end{array}
\right)
\sim
\left(
\begin{array}{ccc|c}
1&0&0&1\\
0&1&-1&-1\\
0&0&0&0
\end{array}
\right).
\]
所以
\[
x_1=1,\qquad x_2-x_3=-1.
\]
令 \(x_3=t\)，则 \(x_2=t-1\)。}
\answer{\[
(x_1,x_2,x_3)=(1,t-1,t),\qquad t\in\mathbb{R}.
\]}
\end{solutionblock}

\begin{solutionblock}
\textbf{3.2 例 5}

\question{求齐次线性方程组的通解，其系数矩阵为
\[
\begin{pmatrix}
1&3&5&1\\
2&3&4&2\\
1&2&3&1
\end{pmatrix}.
\]}


\analysis{齐次方程组对应增广矩阵化简为
\[
\left(
\begin{array}{cccc|c}
1&3&5&1&0\\
2&3&4&2&0\\
1&2&3&1&0
\end{array}
\right)
\sim
\left(
\begin{array}{cccc|c}
1&0&-1&1&0\\
0&1&2&0&0\\
0&0&0&0&0
\end{array}
\right).
\]
主元变量为 \(x_1,x_2\)，自由变量可取 \(x_3,x_4\)。由化简矩阵得
\[
x_1-x_3+x_4=0,\qquad x_2+2x_3=0.
\]
令 \(x_3=s,\ x_4=t\)，则
\[
x_1=s-t,\qquad x_2=-2s.
\]}
\answer{\[
(x_1,x_2,x_3,x_4)=(s-t,-2s,s,t),\qquad s,t\in\mathbb{R}.
\]}
\end{solutionblock}

\begin{solutionblock}
\textbf{3.2 例 6}

\question{求非齐次线性方程组的通解，其增广矩阵为
\[
\left(
\begin{array}{cccc|c}
1&3&5&1&2\\
2&3&4&2&1\\
1&2&3&1&1
\end{array}
\right).
\]}


\analysis{非齐次方程组的增广矩阵化简为
\[
\left(
\begin{array}{cccc|c}
1&3&5&1&2\\
2&3&4&2&1\\
1&2&3&1&1
\end{array}
\right)
\sim
\left(
\begin{array}{cccc|c}
1&0&-1&1&-1\\
0&1&2&0&1\\
0&0&0&0&0
\end{array}
\right).
\]
因此
\[
x_1-x_3+x_4=-1,\qquad x_2+2x_3=1.
\]
令 \(x_3=s,\ x_4=t\)，则
\[
x_1=s-t-1,\qquad x_2=1-2s.
\]}
\answer{\[
(x_1,x_2,x_3,x_4)=(s-t-1,1-2s,s,t),\qquad s,t\in\mathbb{R}.
\]}
\end{solutionblock}

\subsection{3.3 方程组解的判定}

\begin{solutionblock}
\textbf{3.3 例 1}

\question{设齐次方程组的阶梯形系数矩阵为
\[
\begin{pmatrix}
1&-1&0&-2&5\\
0&0&1&3&0\\
0&0&0&0&1
\end{pmatrix}.
\]
判断题给四组选项中哪一组可以作为自由变量。}


\analysis{阶梯形系数矩阵为
\[
\begin{pmatrix}
1&-1&0&-2&5\\
0&0&1&3&0\\
0&0&0&0&1
\end{pmatrix}.
\]
主元列为第 \(1,3,5\) 列，对应主元变量为 \(x_1,x_3,x_5\)。因此自由变量只能从非主元变量
\[
x_2,\ x_4
\]
中选取。

题目给出的四个选项分别为 \(x_2,x_3\)、\(x_2,x_5\)、\(x_1,x_4\)、\(x_1,x_2\)，每一组都含有主元变量，所以按扫描件现有题面看，四个选项都不能作为完整的自由变量组。}
\answer{自由变量应取 \(x_2,x_4\)。按现有四个选项，A、B、C、D 均不能取；疑似选项漏印 \(x_2,x_4\)。}
\examnote{阶梯形矩阵中，每个非零行的首非零元所在列为主元列；自由变量只能从非主元列中选。}
\end{solutionblock}

\begin{solutionblock}
\textbf{3.3 例 2}

\question{设 \(A\) 为 \(m\times n\) 矩阵，\(B\) 为 \(n\times m\) 矩阵。若 \(m>n\)，判断齐次方程组 \((AB)x=0\) 的解的情况及正确选项。}


\analysis{\(A\) 为 \(m\times n\)，\(B\) 为 \(n\times m\)，所以 \(AB\) 为 \(m\times m\) 矩阵。齐次方程组 \((AB)x=0\) 的未知量个数为 \(m\)。

若 \(m>n\)，则
\[
r(AB)\le n<m,
\]
所以系数矩阵 \(AB\) 不满秩，齐次方程组必有非零解。

若 \(n>m\)，\(AB\) 可能可逆。例如可取
\[
A=(E_m\ O),\qquad B=\begin{pmatrix}E_m\\O\end{pmatrix},
\]
则 \(AB=E_m\)，此时只有零解。因此 \(n>m\) 不能推出必有非零解。}
\answer{选 D。}
\end{solutionblock}

\begin{solutionblock}
\textbf{3.3 例 3}

\question{设非齐次线性方程组为 \(Ax=b\)，其中 \(A\) 为 \(m\times n\) 矩阵。判断关于 \(r(A)\) 与方程组解的四个命题中哪一个正确。}


\analysis{非齐次方程组 \(Ax=b\) 中，\(A\) 为 \(m\times n\) 矩阵。

A 项正确。若 \(r(A)=m\)，说明 \(A\) 的列空间维数等于 \(\mathbb{R}^m\) 的维数，即列空间就是整个 \(\mathbb{R}^m\)，因此任意 \(b\) 都可由 \(A\) 的列向量线性表示，方程 \(Ax=b\) 有解。

B 项错误。若 \(r(A)=m\)，不一定有 \(Ax=0\) 的非零解。例如 \(m=n\) 且 \(A\) 可逆时，齐次方程只有零解。

C 项错误。\(r(A)=n\) 只说明齐次方程只有零解；非齐次方程若有解则唯一，但不保证一定有解。

D 项错误。\(r(A)<n\) 说明齐次方程有非零解；非齐次方程可能无解，也可能有无穷多解。}
\answer{选 A。}
\end{solutionblock}

\begin{solutionblock}
\textbf{3.3 例 4}

\question{设
\[
A=\begin{pmatrix}
1&1&1\\
1&2&a\\
1&4&a^2
\end{pmatrix}.
\]
若方程组 \(Ax=b\) 有无穷多解，判断参数 \(a,b\) 与集合 \(\Omega=\{1,2\}\) 的关系。}


\analysis{矩阵
\[
A=\begin{pmatrix}
1&1&1\\
1&2&a\\
1&4&a^2
\end{pmatrix}
\]
的行列式为
\[
|A|=(a-1)(a-2).
\]
若方程组 \(Ax=b\) 有无穷多解，必须有 \(|A|=0\)，即
\[
a\in\{1,2\}=\Omega.
\]
再看相容性。常数列为
\[
\begin{pmatrix}1\\b\\b^2\end{pmatrix}.
\]
当 \(a=1\) 或 \(a=2\) 时，系数矩阵秩为 \(2\)。要有无穷多解，需要增广矩阵秩也为 \(2\)。计算三阶子式可得条件
\[
(b-1)(b-2)=0,
\]
即
\[
b\in\{1,2\}=\Omega.
\]
此时 \(r(A)=r(A\mid b)=2<3\)，所以确有无穷多解。}
\answer{\(a\in\Omega,\ b\in\Omega\)，选 D。}
\end{solutionblock}

\begin{solutionblock}
\textbf{3.3 例 5}

\question{设 \(Ax=0\) 是非齐次线性方程组 \(Ax=b\) 的导出组。判断关于导出组和非齐次方程组解的四个命题中哪一个正确。}


\analysis{设 \(Ax=0\) 是非齐次方程组 \(Ax=b\) 的导出组。

A 项错误。若 \(Ax=0\) 只有零解，则 \(r(A)=n\)。非齐次方程若有解则唯一，但它不一定有解。

B 项错误。若 \(Ax=0\) 有非零解，则 \(r(A)<n\)。此时非齐次方程若有解，则必有无穷多解，但仍可能无解。

C 项正确。若 \(Ax=b\) 有无穷多解，则两个不同解之差是导出组 \(Ax=0\) 的非零解，因此导出组必有非零解。

D 项错误。若 \(Ax=b\) 有无穷多解，则导出组不可能只有零解。}
\answer{选 C。}
\examnote{非齐次方程的解集若非空，则为“一个特解 \(+\) 齐次方程通解”。有无穷多解等价于导出组有非零解。}
\end{solutionblock}

\begin{solutionblock}
\textbf{3.3 例 6}

\question{设线性方程组的系数矩阵为
\[
A=(a-b)E+bJ,
\]
其中 \(J\) 为 \(n\) 阶全 \(1\) 矩阵。讨论齐次方程组 \(Ax=0\) 的通解。}


\analysis{系数矩阵为
\[
A=(a-b)E+bJ,
\]
其中 \(J\) 是 \(n\) 阶全 \(1\) 矩阵。矩阵 \(J\) 的特征值为 \(n,0,\ldots,0\)，所以 \(A\) 的特征值为
\[
a+(n-1)b,\qquad a-b\quad(\text{\(n-1\) 重}).
\]
因此
\[
|A|=(a-b)^{n-1}\bigl(a+(n-1)b\bigr).
\]

齐次方程组只有零解的充要条件是 \(|A|\ne0\)，故
\[
a\ne b,\qquad a+(n-1)b\ne0.
\]

有无穷多解的充要条件是 \(|A|=0\)，即
\[
a=b\quad\text{或}\quad a=-(n-1)b.
\]
若 \(a=b\ne0\)，方程组每个方程都化为
\[
b(x_1+x_2+\cdots+x_n)=0,
\]
所以全部解为
\[
x_1+x_2+\cdots+x_n=0.
\]
可写成
\[
(x_1,\ldots,x_n)=(t_1,\ldots,t_{n-1},-t_1-\cdots-t_{n-1}).
\]

若 \(a=-(n-1)b\)，零特征向量方向为 \((1,1,\ldots,1)^T\)，全部解为
\[
(x_1,\ldots,x_n)=t(1,1,\ldots,1).
\]}
\answer{只有零解：
\[
a\ne b,\quad a+(n-1)b\ne0.
\]
无穷多解：
\[
a=b\quad\text{或}\quad a=-(n-1)b.
\]
对应通解如上。}
\end{solutionblock}

\begin{solutionblock}
\textbf{3.3 例 7}

\question{设线性方程组的系数矩阵为
\[
A=\begin{pmatrix}
1&1&1&1\\
0&1&0&2\\
0&-1&a-3&-2\\
3&2&1&a
\end{pmatrix}.
\]
讨论方程组的解的情况，并在有无穷多解时写出通解。}


\analysis{系数矩阵为
\[
A=\begin{pmatrix}
1&1&1&1\\
0&1&0&2\\
0&-1&a-3&-2\\
3&2&1&a
\end{pmatrix}.
\]
其行列式为
\[
|A|=(a-1)(a-3).
\]
因此当
\[
a\ne1,\quad a\ne3
\]
时，方程组有唯一解。此时直接行化简可得
\[
x_1=-\frac{a+b}{a-1},\quad
x_2=\frac{a^2-4a-4b-1}{(a-1)(a-3)},\quad
x_3=\frac{b+1}{a-3},\quad
x_4=\frac{2(b+1)}{(a-1)(a-3)}.
\]

当 \(a=1\) 或 \(a=3\) 时，系数矩阵秩为 \(3\)。进一步检查增广矩阵，均得到相容条件
\[
b=-1.
\]
所以：

若 \(a=1,b\ne-1\)，或 \(a=3,b\ne-1\)，方程组无解。

若 \(a=1,b=-1\)，方程组有无穷多解，通解为
\[
(x_1,x_2,x_3,x_4)=(t-1,1-2t,0,t).
\]

若 \(a=3,b=-1\)，方程组有无穷多解，通解为
\[
(x_1,x_2,x_3,x_4)=(-1,1-2t,t,t).
\]}
\answer{\[
\begin{array}{ll}
a\ne1,3: & \text{唯一解};\\[1mm]
a=1\text{ 或 }3,\ b\ne-1: & \text{无解};\\[1mm]
(a,b)=(1,-1)\text{ 或 }(3,-1): & \text{无穷多解}.
\end{array}
\]
无穷多解时通解如上。}
\examnote{含参数非齐次方程组先看 \(|A|\)。若 \(|A|\ne0\)，必唯一；若 \(|A|=0\)，再比较 \(r(A)\) 与 \(r(A\mid b)\)。}
\end{solutionblock}

\chapter{向量与方程组}

\section{逐题解析}

\subsection{4.1 向量的基本概念}

\begin{solutionblock}
\textbf{4.1 例 1}

\question{已知 \(\alpha_1=(\frac12,0,x)^T\) 与 \(\alpha_2=(\frac{\sqrt3}{2},y,\frac12)^T\) 均为单位向量，且二者正交，求 \(x,y\)。}


\analysis{由 \(\alpha_1=(\frac12,0,x)^T\) 为单位向量得
\[
\frac14+x^2=1,\qquad x=\pm\frac{\sqrt3}{2}.
\]
由 \(\alpha_2=(\frac{\sqrt3}{2},y,\frac12)^T\) 为单位向量得
\[
\frac34+y^2+\frac14=1,\qquad y=0.
\]
再由两向量正交，
\[
\alpha_1^T\alpha_2=\frac12\cdot\frac{\sqrt3}{2}+x\cdot\frac12
=\frac{\sqrt3}{4}+\frac{x}{2}=0,
\]
所以
\[
x=-\frac{\sqrt3}{2}.
\]}
\answer{\[
x=-\frac{\sqrt3}{2},\qquad y=0.
\]}
\end{solutionblock}

\begin{solutionblock}
\textbf{4.1 例 2}

\question{设
\[
\alpha=(a,0,\ldots,0,a)^T,
\qquad a<0,\qquad A=E-\alpha\alpha^T.
\]
若
\[
A^{-1}=E+a^{-1}\alpha\alpha^T,
\]
求 \(a\)。}


\analysis{题设
\[
\alpha=(a,0,\ldots,0,a)^T,\qquad a<0,\qquad A=E-\alpha\alpha^T.
\]
对秩一修正矩阵，用 Sherman-Morrison 公式：
\[
(E-\alpha\alpha^T)^{-1}
=E+\frac{\alpha\alpha^T}{1-\alpha^T\alpha}.
\]
题中又给出
\[
A^{-1}=E+a^{-1}\alpha\alpha^T,
\]
故
\[
\frac1{1-\alpha^T\alpha}=\frac1a,
\qquad a=1-\alpha^T\alpha.
\]
而
\[
\alpha^T\alpha=2a^2.
\]
于是
\[
a=1-2a^2,\qquad 2a^2+a-1=0.
\]
解得
\[
a=\frac12\quad\text{或}\quad a=-1.
\]
由 \(a<0\)，取
\[
a=-1.
\]}
\answer{\(-1\)}
\examnote{秩一修正 \(E-uv^T\) 的逆是考研常见模型；本题也可直接利用 \(\alpha\alpha^T\alpha=(\alpha^T\alpha)\alpha\) 验证。}
\end{solutionblock}

\begin{solutionblock}
\textbf{4.1 例 3}

\question{已知
\[
\alpha\beta^T=
\begin{pmatrix}
1&-1&2\\
-2&2&-4\\
3&-3&6
\end{pmatrix}.
\]
求 \(\alpha^T\beta\)。}


\analysis{矩阵 \(\alpha\beta^T\) 的迹为
\[
\operatorname{tr}(\alpha\beta^T)=\beta^T\alpha=\alpha^T\beta.
\]
题中
\[
\alpha\beta^T=
\begin{pmatrix}
1&-1&2\\
-2&2&-4\\
3&-3&6
\end{pmatrix},
\]
所以
\[
\alpha^T\beta=1+2+6=9.
\]}
\answer{\(9\)}
\end{solutionblock}

\subsection{4.3 线性相关、线性无关与线性表示}

\begin{solutionblock}
\textbf{4.3 例 1}

\question{证明：单个非零向量线性无关；两个不成比例的向量线性无关。}


\analysis{单个非零向量 \(\alpha\) 若线性相关，则存在数 \(k\ne0\)，使
\[
k\alpha=0.
\]
这推出 \(\alpha=0\)，与非零矛盾，所以单个非零向量线性无关。

若两个向量 \(\alpha,\beta\) 不成比例而线性相关，则存在不全为零的 \(k,l\)，使
\[
k\alpha+l\beta=0.
\]
若 \(k\ne0\)，则 \(\alpha=-\frac{l}{k}\beta\)；若 \(l\ne0\)，则 \(\beta=-\frac{k}{l}\alpha\)。两种情况都推出二者成比例，矛盾。因此两个不成比例的向量线性无关。}
\answer{结论成立。}
\end{solutionblock}

\begin{solutionblock}
\textbf{4.3 例 2}

\question{现有四个向量组：
\[
\begin{array}{ll}
\text{(1)}\ (1,2,1)^T,\ (2,1,2)^T,\ (0,0,0)^T,
&
\text{(2)}\ (1,2,3)^T,\ (0,1,2)^T,\ (0,0,1)^T,\\[3pt]
\text{(3)}\ (1,2,1)^T,\ (2,4,2)^T,
&
\text{(4)}\ (1,2,1)^T,\ (2,4,2)^T,\ (1,6,9)^T.
\end{array}
\]
其中线性相关的向量组有多少个？
\[
\text{A. }1\text{ 个}\qquad
\text{B. }2\text{ 个}\qquad
\text{C. }3\text{ 个}\qquad
\text{D. }4\text{ 个}.
\]}


\analysis{逐组判断。

第 1 组含零向量，必线性相关。

第 2 组
\[
(1,2,3)^T,\ (0,1,2)^T,\ (0,0,1)^T
\]
组成的矩阵为上三角型，行列式为 \(1\ne0\)，故线性无关。

第 3 组中
\[
(2,4,2)^T=2(1,2,1)^T,
\]
故线性相关。

第 4 组前两个向量
\[
(2,4,2)^T=2(1,2,1)^T,
\]
已经成比例，所以整组线性相关。}
\answer{线性相关的向量组有 \(3\) 个，选 C。}
\end{solutionblock}

\begin{solutionblock}
\textbf{4.3 例 3}

\question{设原向量组 \(\alpha_1,\alpha_2,\alpha_3\) 线性无关。选项中的新向量组均可写成
\[
(\alpha_1,\alpha_2,\alpha_3)P.
\]
新向量组线性无关的充要条件是系数矩阵 \(P\) 可逆，即 \(|P|\ne0\)。

对四个选项计算系数矩阵行列式，只有选项 C 的系数矩阵
\[
P_C=
\begin{pmatrix}
1&0&1\\
2&2&0\\
0&3&3
\end{pmatrix}
\]
满足
\[
|P_C|=12\ne0.
\]
其他选项的系数矩阵行列式均为 \(0\)。}


\analysis{设原向量组 \(\alpha_1,\alpha_2,\alpha_3\) 线性无关。选项中的新向量组均可写成
\[
(\alpha_1,\alpha_2,\alpha_3)P.
\]
新向量组线性无关的充要条件是系数矩阵 \(P\) 可逆，即 \(|P|\ne0\)。

对四个选项计算系数矩阵行列式，只有选项 C 的系数矩阵
\[
P_C=
\begin{pmatrix}
1&0&1\\
2&2&0\\
0&3&3
\end{pmatrix}
\]
满足
\[
|P_C|=12\ne0.
\]
其他选项的系数矩阵行列式均为 \(0\)。}
\answer{选 C。}
\examnote{原向量组无关时，“线性组合后的向量组是否无关”转化为“线性组合系数矩阵是否可逆”。}
\end{solutionblock}

\begin{solutionblock}
\textbf{4.3 例 4}

\question{设 \(\alpha_1,\ldots,\alpha_s\) 线性无关，令
\[
\beta_1=\alpha_1,\quad
\beta_2=\alpha_1+\alpha_2,\quad
\ldots,\quad
\beta_s=\alpha_1+\cdots+\alpha_s.
\]
证明 \(\beta_1,\ldots,\beta_s\) 线性无关。}


\analysis{设
\[
\beta_1=\alpha_1,\quad
\beta_2=\alpha_1+\alpha_2,\quad \ldots,\quad
\beta_s=\alpha_1+\cdots+\alpha_s.
\]
若
\[
k_1\beta_1+k_2\beta_2+\cdots+k_s\beta_s=0,
\]
则整理为
\[
(k_1+\cdots+k_s)\alpha_1+(k_2+\cdots+k_s)\alpha_2+\cdots+k_s\alpha_s=0.
\]
由于 \(\alpha_1,\ldots,\alpha_s\) 线性无关，所以
\[
k_s=0,\quad k_{s-1}+k_s=0,\quad \ldots,\quad k_1+\cdots+k_s=0.
\]
逆推得
\[
k_s=k_{s-1}=\cdots=k_1=0.
\]
故 \(\beta_1,\ldots,\beta_s\) 线性无关。}
\answer{向量组 \(\alpha_1,\alpha_1+\alpha_2,\ldots,\alpha_1+\cdots+\alpha_s\) 线性无关。}
\end{solutionblock}

\begin{solutionblock}
\textbf{4.3 例 5}

\question{设
\[
A\alpha_1=\alpha_1,\quad
A\alpha_2=\alpha_1+\alpha_2,\quad
A\alpha_3=\alpha_2+\alpha_3.
\]
证明 \(\alpha_1,\alpha_2,\alpha_3\) 线性无关。}


\analysis{记 \(T=A-E\)。由题设
\[
A\alpha_1=\alpha_1,\quad A\alpha_2=\alpha_1+\alpha_2,\quad A\alpha_3=\alpha_2+\alpha_3
\]
可得
\[
T\alpha_1=0,\qquad T\alpha_2=\alpha_1,\qquad T\alpha_3=\alpha_2.
\]
设
\[
c_1\alpha_1+c_2\alpha_2+c_3\alpha_3=0.
\]
对等式两边作用 \(T^2\)，得
\[
c_3\alpha_1=0.
\]
因为 \(\alpha_1\ne0\)，所以 \(c_3=0\)。原式化为
\[
c_1\alpha_1+c_2\alpha_2=0.
\]
再作用 \(T\)，得
\[
c_2\alpha_1=0,
\]
故 \(c_2=0\)，进而 \(c_1=0\)。因此三向量线性无关。}
\answer{\(\alpha_1,\alpha_2,\alpha_3\) 线性无关。}
\end{solutionblock}

\begin{solutionblock}
\textbf{4.3 例 6}

\question{设 \(A^k\alpha=0\)，且 \(A^{k-1}\alpha\ne0\)。证明向量组
\[
\alpha,A\alpha,\ldots,A^{k-1}\alpha
\]
线性无关。}


\analysis{设存在
\[
c_0\alpha+c_1A\alpha+\cdots+c_{k-1}A^{k-1}\alpha=0.
\]
两边左乘 \(A^{k-1}\)，因 \(A^k\alpha=0\)，只剩
\[
c_0A^{k-1}\alpha=0.
\]
题设 \(A^{k-1}\alpha\ne0\)，故 \(c_0=0\)。再对剩余等式左乘 \(A^{k-2}\)，得
\[
c_1A^{k-1}\alpha=0,
\]
故 \(c_1=0\)。依次进行，最终得到
\[
c_0=c_1=\cdots=c_{k-1}=0.
\]
所以
\[
\alpha,A\alpha,\ldots,A^{k-1}\alpha
\]
线性无关。}
\answer{结论成立。}
\examnote{这类题的关键是用高次幂逐步“杀掉”后面的项，只留下一个最高阶非零向量。}
\end{solutionblock}

\begin{solutionblock}
\textbf{4.3 例 7}

\question{设
\[
\alpha_1=(1,2,-1,5)^T,\quad
\alpha_2=(2,-1,1,1)^T.
\]
判断 \(\beta_1=(4,3,-1,11)^T\)、\(\beta_2=(4,3,0,11)^T\) 是否可由 \(\alpha_1,\alpha_2\) 线性表示；若可以，写出表示式。}


\analysis{设
\[
\alpha_1=(1,2,-1,5)^T,\quad \alpha_2=(2,-1,1,1)^T.
\]
先判断 \(\beta_1=(4,3,-1,11)^T\)。令
\[
\beta_1=x\alpha_1+y\alpha_2.
\]
解得
\[
x=2,\qquad y=1,
\]
即
\[
\beta_1=2\alpha_1+\alpha_2.
\]

再判断 \(\beta_2=(4,3,0,11)^T\)。若
\[
\beta_2=x\alpha_1+y\alpha_2,
\]
由前两行可解出 \(x=2,y=1\)，但第三行给出
\[
-2+1=-1\ne0,
\]
矛盾。所以 \(\beta_2\) 不能由 \(\alpha_1,\alpha_2\) 线性表示。}
\answer{\[
\beta_1=2\alpha_1+\alpha_2,\qquad \beta_2\text{ 不能由 }\alpha_1,\alpha_2\text{ 线性表示}.
\]}
\end{solutionblock}

\begin{solutionblock}
\textbf{4.3 例 8}

\question{设向量组 \((I):\alpha_1,\alpha_2,\alpha_3\) 与 \((II):\beta_1,\beta_2,\beta_3\)，其中
\[
\alpha_1=\begin{pmatrix}0\\1\\2\\3\end{pmatrix},\quad
\alpha_2=\begin{pmatrix}3\\0\\1\\2\end{pmatrix},\quad
\alpha_3=\begin{pmatrix}2\\3\\0\\1\end{pmatrix}.
\]
判断 \((II)\) 是否可由 \((I)\) 线性表示。}


\analysis{记
\[
A=(\alpha_1,\alpha_2,\alpha_3),\qquad
B=(\beta_1,\beta_2,\beta_3).
\]
其中
\[
\alpha_1=\begin{pmatrix}0\\1\\2\\3\end{pmatrix},\ 
\alpha_2=\begin{pmatrix}3\\0\\1\\2\end{pmatrix},\ 
\alpha_3=\begin{pmatrix}2\\3\\0\\1\end{pmatrix}.
\]
对每个 \(\beta_i\) 解方程 \(Ax=\beta_i\)，得到
\[
\beta_1=\frac14\alpha_1+\frac12\alpha_2+\frac14\alpha_3,
\]
\[
\beta_2=\frac14\alpha_1+\frac12\alpha_2-\frac34\alpha_3,
\]
\[
\beta_3=\frac14\alpha_1+\frac12\alpha_2+\frac54\alpha_3.
\]
故 \((II)\) 中每个向量均可由 \((I)\) 线性表示。}
\answer{\((II)\) 可由 \((I)\) 线性表示。}
\end{solutionblock}

\begin{solutionblock}
\textbf{4.3 例 9}

\question{判断向量组线性相关与“某个向量可由其余向量线性表示”之间的等价关系，并选择正确选项。}


\analysis{向量组线性相关的充要条件是：存在一个向量可以由其余向量线性表示。

若某个向量可由其余向量线性表示，则把等式移到一边，就得到一个系数不全为零的线性关系，所以线性相关。

反过来，若向量组线性相关，则存在不全为零的 \(k_i\)，使
\[
k_1\alpha_1+\cdots+k_s\alpha_s=0.
\]
取一个非零系数 \(k_j\)，即可把 \(\alpha_j\) 表成其余向量的线性组合。}
\answer{选 C。}
\end{solutionblock}

\begin{solutionblock}
\textbf{4.3 例 10}

\question{判断向量组线性无关与“任一向量都不能由其余向量线性表示”之间的等价关系，并选择正确选项。}


\analysis{向量组线性无关的充要条件是：其中每一个向量都不能由其余 \(s-1\) 个向量线性表示。

若某个向量能由其余向量表示，则整组线性相关；若整组线性相关，则由例 9 的结论，必存在一个向量能由其余向量表示。因此“不存在这种向量”正是线性无关的等价刻画。}
\answer{选 D。}
\end{solutionblock}

\begin{solutionblock}
\textbf{4.3 例 11}

\question{设 \[
\alpha_1=(6,a+1,3)^T,\quad
\alpha_2=(a,2,-2)^T,\quad
\alpha_3=(a,1,0)^T.
\]
求：
\[
\begin{aligned}
&\text{(1) }a\text{ 为何值时，}\alpha_1,\alpha_2\text{ 线性相关？}
&&a\text{ 为何值时，}\alpha_1,\alpha_2\text{ 线性无关？}\\
&\text{(2) }a\text{ 为何值时，}\alpha_1,\alpha_2,\alpha_3\text{ 线性相关？}
&&a\text{ 为何值时，}\alpha_1,\alpha_2,\alpha_3\text{ 线性无关？}
\end{aligned}
\]}


\analysis{题中
\[
\alpha_1=(6,a+1,3)^T,\quad
\alpha_2=(a,2,-2)^T,\quad
\alpha_3=(a,1,0)^T.
\]
第 (1) 问，两个向量线性相关等价于成比例。考察 \(\alpha_1,\alpha_2\) 的所有二阶子式，可得共同条件
\[
a=-4.
\]
因此
\[
a=-4\text{ 时 }\alpha_1,\alpha_2\text{ 线性相关};\qquad
a\ne-4\text{ 时线性无关}.
\]

第 (2) 问，三个三维向量线性相关等价于行列式为零：
\[
\det(\alpha_1,\alpha_2,\alpha_3)=-(a+4)(2a-3).
\]
故
\[
a=-4\quad\text{或}\quad a=\frac32
\]
时三向量线性相关；其余 \(a\) 值时三向量线性无关。}
\answer{\[
\alpha_1,\alpha_2\text{ 相关}\Longleftrightarrow a=-4;
\]
\[
\alpha_1,\alpha_2,\alpha_3\text{ 相关}\Longleftrightarrow a=-4\text{ 或 }a=\frac32.
\]}
\end{solutionblock}

\begin{solutionblock}
\textbf{4.3 例 12}

\question{设 \(A=(\alpha_1,\alpha_2,\alpha_3,\alpha_4)\) 为 \(3\times4\) 矩阵，\(\beta_1,\beta_2,\beta_3\) 为其行向量。已知
\[
\begin{vmatrix}a_{12}&a_{14}\\a_{22}&a_{24}\end{vmatrix}\ne0.
\]
判断题给命题中哪些正确。}


\analysis{矩阵 \(A=(\alpha_1,\alpha_2,\alpha_3,\alpha_4)\) 为 \(3\times4\) 矩阵，\(\beta_1,\beta_2,\beta_3\) 为其行向量。

由
\[
\begin{vmatrix}
a_{12}&a_{14}\\
a_{22}&a_{24}
\end{vmatrix}\ne0
\]
可知第 2 列与第 4 列线性无关，即 \(\alpha_2,\alpha_4\) 线性无关，故第 2 条正确。

又
\[
\begin{vmatrix}
a_{11}&a_{12}&a_{13}\\
a_{21}&a_{22}&a_{23}\\
a_{31}&a_{32}&a_{33}
\end{vmatrix}=0
\]
说明第 1、2、3 列线性相关，即 \(\alpha_1,\alpha_2,\alpha_3\) 线性相关，故第 4 条正确。

第 1 条 \(r(A)=2\) 不一定成立，因为 \(\alpha_4\) 可能使矩阵秩达到 \(3\)。第 3 条“行向量线性相关”也不一定成立，只有当 \(r(A)<3\) 时才成立。}
\answer{正确的是第 2、4 条，选 D。}
\end{solutionblock}

\begin{solutionblock}
\textbf{4.3 例 13}

\question{设
\[
A=\begin{pmatrix}
1&1&\cdots&1\\
a_1&a_2&\cdots&a_s\\
\vdots&\vdots&&\vdots\\
a_1^{n-1}&a_2^{n-1}&\cdots&a_s^{n-1}
\end{pmatrix},
\qquad a_i\ne a_j\ (i\ne j).
\]
讨论列向量组 \(\alpha_1,\ldots,\alpha_s\) 的线性相关性。}


\analysis{该矩阵是 \(n\times s\) 的 Vandermonde 型矩阵：
\[
A=\begin{pmatrix}
1&1&\cdots&1\\
a_1&a_2&\cdots&a_s\\
\vdots&\vdots&&\vdots\\
a_1^{n-1}&a_2^{n-1}&\cdots&a_s^{n-1}
\end{pmatrix},
\qquad a_i\ne a_j.
\]
任取不超过 \(n\) 列，得到的方阵或高矩阵含有非零 Vandermonde 子式，故这些列线性无关；而若 \(s>n\)，在 \(n\) 维空间中有 \(s\) 个列向量，必线性相关。

因此
\[
r(A)=\min\{n,s\}.
\]}
\answer{\[
\alpha_1,\ldots,\alpha_s
\begin{cases}
\text{线性无关},& s\le n,\\
\text{线性相关},& s>n.
\end{cases}
\]}
\end{solutionblock}

\begin{solutionblock}
\textbf{4.3 例 14}

\question{设 \(\alpha_1,\alpha_2,\alpha_3\) 线性相关，而 \(\alpha_2,\alpha_3,\alpha_4\) 线性无关。回答：\(\alpha_1\) 能否由 \(\alpha_2,\alpha_3\) 线性表示？\(\alpha_4\) 能否由 \(\alpha_1,\alpha_2,\alpha_3\) 线性表示？}


\analysis{已知 \(\alpha_1,\alpha_2,\alpha_3\) 线性相关，而 \(\alpha_2,\alpha_3,\alpha_4\) 线性无关。

第 (1) 问：由于 \(\alpha_2,\alpha_3,\alpha_4\) 线性无关，所以其子组 \(\alpha_2,\alpha_3\) 线性无关。又 \(\alpha_1,\alpha_2,\alpha_3\) 线性相关，因此 \(\alpha_1\) 必能由 \(\alpha_2,\alpha_3\) 线性表示。

第 (2) 问：若 \(\alpha_4\) 能由 \(\alpha_1,\alpha_2,\alpha_3\) 线性表示，而 \(\alpha_1\) 又能由 \(\alpha_2,\alpha_3\) 表示，则 \(\alpha_4\) 能由 \(\alpha_2,\alpha_3\) 表示。这与 \(\alpha_2,\alpha_3,\alpha_4\) 线性无关矛盾。}
\answer{(1) 能；(2) 不能。}
\end{solutionblock}

\begin{solutionblock}
\textbf{4.3 例 15}

\question{设向量组 \((I)\) 有 \(t\) 个向量，向量组 \((II)\) 有 \(s\) 个向量。判断关于“一个向量组可由另一个向量组线性表示”与线性相关性的四个命题中哪些正确。}


\analysis{记向量组 \((I)\) 有 \(t\) 个向量，\((II)\) 有 \(s\) 个向量。

第 1 条：若 \((I)\) 可由 \((II)\) 线性表示，且 \(s<t\)，则 \((I)\) 的 \(t\) 个向量都落在维数不超过 \(s\) 的线性空间中，故 \((I)\) 线性相关，正确。

第 2 条：若 \((II)\) 可由 \((I)\) 表示，且 \(s<t\)，不能推出 \((I)\) 相关。例如 \((I)\) 本身可线性无关，错误。

第 3 条：若 \((I)\) 可由 \((II)\) 表示，且 \((I)\) 线性无关，则 \(t\) 个无关向量落在 \((II)\) 张成的空间中，故 \(t\le s\)，即 \(s\ge t\)，正确。

第 4 条：若 \((II)\) 可由 \((I)\) 表示，且 \((I)\) 线性无关，不能推出 \(s\ge t\)，错误。}
\answer{正确的是第 1、3 条，选 B。}
\end{solutionblock}

\subsection{4.4 向量组的极大线性无关组、秩}

\begin{solutionblock}
\textbf{4.4 例 1}

\question{设向量组的秩为 \(r\)。判断关于任取 \(r+1\) 个向量的线性相关性的正确选项。}


\analysis{向量组的秩为 \(r\)，表示其中最多只能取出 \(r\) 个线性无关的向量。因此任取 \(r+1\) 个向量，必定线性相关。}
\answer{选 D。}
\end{solutionblock}

\begin{solutionblock}
\textbf{4.4 例 2}

\question{设向量组 \(\alpha_1,\ldots,\alpha_s\) 的秩为 \(s-1\)。判断关于其子组秩和线性表示的四个命题中哪一个正确。}


\analysis{向量组 \(\alpha_1,\ldots,\alpha_s\) 的秩为 \(s-1\)。对任意 \(1<r<s\)，其前 \(r\) 个向量的秩满足
\[
r(\alpha_1,\ldots,\alpha_r)\ge (s-1)-(s-r)=r-1.
\]
这是因为后面 \(s-r\) 个向量最多只能使秩增加 \(s-r\)。

A 项“有两个向量成比例”不一定；B 项“任一向量都可由其余向量表示”不一定，例如零向量的情形会破坏该结论；D 项要求任意前 \(r\) 个都满秩，也不一定。}
\answer{选 C。}
\end{solutionblock}

\begin{solutionblock}
\textbf{4.4 例 3}

\question{设矩阵经初等行变换化为
\[
B=\begin{pmatrix}
1&2&2&1\\
0&1&2&2\\
0&0&0&1\\
0&0&0&0
\end{pmatrix}.
\]
判断原矩阵列向量组的极大无关组及线性表示关系，并选择不正确的选项。}


\analysis{初等行变换等价于左乘可逆矩阵，不改变列向量之间的线性关系。化简后的矩阵列向量记为 \(b_1,b_2,b_3,b_4\)，其中
\[
B=\begin{pmatrix}
1&2&2&1\\
0&1&2&2\\
0&0&0&1\\
0&0&0&0
\end{pmatrix}.
\]
可见 \(b_1,b_2,b_4\) 为主元列，线性无关；且
\[
b_3=-2b_1+2b_2.
\]
所以 \(\alpha_1,\alpha_3,\alpha_4\) 线性无关，\(\alpha_2\) 可由 \(\alpha_1,\alpha_3\) 线性表示，\(\alpha_4\) 不能由 \(\alpha_1,\alpha_2,\alpha_3\) 线性表示。}
\answer{不正确的是 D。}
\end{solutionblock}

\begin{solutionblock}
\textbf{4.4 例 4}

\question{设五个向量 \(\alpha_1,\ldots,\alpha_5\) 的秩为 \(3\)。求一个包含 \(\alpha_1,\alpha_3\) 的极大线性无关组，并将其余向量由该组线性表示。}


\analysis{把五个向量作列组成矩阵。行化简得秩为
\[
r(\alpha_1,\alpha_2,\alpha_3,\alpha_4,\alpha_5)=3.
\]
题目要求极大无关组包含 \(\alpha_1,\alpha_3\)。可取
\[
\alpha_1,\alpha_3,\alpha_4
\]
为一个极大线性无关组。解线性表示方程得到
\[
\alpha_2=-3\alpha_1+\alpha_3,
\]
\[
\alpha_5=2\alpha_1-\alpha_3+\alpha_4.
\]}
\answer{秩为 \(3\)。一个包含 \(\alpha_1,\alpha_3\) 的极大线性无关组为
\[
\alpha_1,\alpha_3,\alpha_4.
\]
其余向量表示为
\[
\alpha_2=-3\alpha_1+\alpha_3,\qquad
\alpha_5=2\alpha_1-\alpha_3+\alpha_4.
\]}
\method{方法二}
也可取 \(\alpha_1,\alpha_3,\alpha_5\) 为极大无关组，此时
\[
\alpha_2=-3\alpha_1+\alpha_3,\qquad
\alpha_4=-2\alpha_1+\alpha_3+\alpha_5.
\]
\end{solutionblock}

\subsection{4.5 向量组等价}

\begin{solutionblock}
\textbf{4.5 例 1}

\question{已知向量组 \((I)\) 与 \((II)\) 均由三个三维向量组成，且
\[
\det(\alpha_1,\alpha_2,\alpha_3)=a+1,\qquad
\det(\beta_1,\beta_2,\beta_3)=6.
\]
判断两向量组何时等价。}


\analysis{向量组 \((I)\) 的矩阵行列式为
\[
\det(\alpha_1,\alpha_2,\alpha_3)=a+1.
\]
向量组 \((II)\) 的矩阵行列式为
\[
\det(\beta_1,\beta_2,\beta_3)=6\ne0.
\]
因此 \((II)\) 总是 \(\mathbb{R}^3\) 的一组基。

当 \(a\ne-1\) 时，\((I)\) 也是 \(\mathbb{R}^3\) 的一组基，所以两向量组等价。

当 \(a=-1\) 时，\((I)\) 的秩为 \(2\)，而 \((II)\) 的秩为 \(3\)，二者不等价。}
\answer{\[
a\ne-1\text{ 时等价};\qquad a=-1\text{ 时不等价}.
\]}
\end{solutionblock}

\begin{solutionblock}
\textbf{4.5 例 2}

\question{设
\[
k\alpha+l\beta+m\gamma=0,\qquad km\ne0.
\]
判断 \(\alpha,\beta,\gamma\) 之间的线性表示关系，并选择正确选项。}


\analysis{已知
\[
k\alpha+l\beta+m\gamma=0,\qquad km\ne0.
\]
因为 \(m\ne0\)，所以
\[
\gamma=-\frac{k}{m}\alpha-\frac{l}{m}\beta,
\]
即 \(\gamma\) 可由 \(\alpha,\beta\) 线性表示。
因为 \(k\ne0\)，所以
\[
\alpha=-\frac{l}{k}\beta-\frac{m}{k}\gamma,
\]
即 \(\alpha\) 可由 \(\beta,\gamma\) 线性表示。
因此
\[
\operatorname{span}\{\alpha,\beta\}
=\operatorname{span}\{\beta,\gamma\}.
\]}
\answer{选 B。}
\end{solutionblock}

\begin{solutionblock}
\textbf{4.5 例 3}

\question{设向量组 \((I)\) 有 \(k\) 个 \(n\) 维向量且线性无关。判断向量组 \((II)\) 线性无关与 \((I),(II)\) 等价之间的关系，并选择正确选项。}


\analysis{向量组 \((I)\) 有 \(k\) 个 \(n\) 维向量且线性无关，所以矩阵
\[
(\alpha_1,\ldots,\alpha_k)
\]
的秩为 \(k\)。要使 \((II)\) 也线性无关，充要条件是
\[
r(\beta_1,\ldots,\beta_k)=k.
\]
两个同型矩阵等价的充要条件是秩相等。因此
\[
(\alpha_1,\ldots,\alpha_k)\text{ 与 }(\beta_1,\ldots,\beta_k)\text{ 等价}
\]
正好等价于 \((II)\) 的秩也是 \(k\)。}
\answer{选 D。}
\end{solutionblock}

\begin{solutionblock}
\textbf{4.5 例 4}

\question{设 \(AB=C\)，且 \(B\) 可逆。判断矩阵 \(A\) 与 \(C\) 的列向量组是否等价，并选择正确选项。}


\analysis{由
\[
AB=C,\qquad B\text{ 可逆}
\]
可知
\[
C=A B.
\]
右乘可逆矩阵只对列向量组作可逆线性组合，因此不改变列向量组张成的空间，也不改变列向量组的等价性。故矩阵 \(C\) 的列向量组与矩阵 \(A\) 的列向量组等价。}
\answer{选 B。}
\end{solutionblock}

\subsection{4.6 基础解系与解的结构}

\begin{solutionblock}
\textbf{4.6 例 1}

\question{求齐次线性方程组的基础解系和通解，其系数矩阵为
\[
\begin{pmatrix}
1&1&1&1&1\\
1&2&1&1&-1\\
1&3&1&1&-3\\
3&4&3&3&1
\end{pmatrix}.
\]}


\analysis{齐次方程组的系数矩阵行化简为
\[
\begin{pmatrix}
1&1&1&1&1\\
1&2&1&1&-1\\
1&3&1&1&-3\\
3&4&3&3&1
\end{pmatrix}
\sim
\begin{pmatrix}
1&0&1&1&3\\
0&1&0&0&-2\\
0&0&0&0&0\\
0&0&0&0&0
\end{pmatrix}.
\]
故
\[
x_1+x_3+x_4+3x_5=0,\qquad x_2-2x_5=0.
\]
令 \(x_3=s,\ x_4=t,\ x_5=u\)，则
\[
x=(-s-t-3u,\,2u,\,s,\,t,\,u)^T.
\]}
\answer{一个基础解系为
\[
\eta_1=(-1,0,1,0,0)^T,\quad
\eta_2=(-1,0,0,1,0)^T,\quad
\eta_3=(-3,2,0,0,1)^T.
\]
通解为
\[
x=c_1\eta_1+c_2\eta_2+c_3\eta_3.
\]}
\end{solutionblock}

\begin{solutionblock}
\textbf{4.6 例 2}

\question{设 \(\alpha_1,\alpha_2,\alpha_3\) 是齐次方程组的一个基础解系。判断题给四个向量组中哪一个仍是基础解系。}


\analysis{基础解系必须满足两个条件：线性无关；张成齐次方程组的解空间。

选项 D 中向量组
\[
\alpha_1,\quad \alpha_1+\alpha_2,\quad \alpha_1+\alpha_2+\alpha_3
\]
相对于原基础解系 \(\alpha_1,\alpha_2,\alpha_3\) 的系数矩阵为
\[
\begin{pmatrix}
1&1&1\\
0&1&1\\
0&0&1
\end{pmatrix},
\]
行列式为 \(1\ne0\)，所以它与原基础解系等价且线性无关，仍是基础解系。

选项 B 的系数矩阵行列式为 \(0\)，不是基础解系；A、C 的描述过宽，不足以保证“向量个数正确且线性无关”。}
\answer{选 D。}
\end{solutionblock}

\begin{solutionblock}
\textbf{4.6 例 3}

\question{设 \(A=(\alpha_1,\alpha_2,\alpha_3)\)，其中 \(\alpha_3=-\alpha_1+2\alpha_2\)，且 \(\alpha_1,\alpha_2\) 线性无关。求齐次方程组 \(Ax=0\) 的通解。}


\analysis{由
\[
A=(\alpha_1,\alpha_2,\alpha_3),\qquad \alpha_3=-\alpha_1+2\alpha_2,
\]
且 \(\alpha_1,\alpha_2\) 线性无关。方程 \(Ax=0\) 即
\[
x_1\alpha_1+x_2\alpha_2+x_3\alpha_3=0.
\]
代入 \(\alpha_3\)，得
\[
(x_1-x_3)\alpha_1+(x_2+2x_3)\alpha_2=0.
\]
由 \(\alpha_1,\alpha_2\) 无关，
\[
x_1=x_3,\qquad x_2=-2x_3.
\]}
\answer{\[
x=t(1,-2,1)^T,\qquad t\in\mathbb{R}.
\]}
\end{solutionblock}

\begin{solutionblock}
\textbf{4.6 例 4}

\question{设 \(r(B)=2\)，\(r(AB)=1\)。结合题给矩阵条件，求参数 \(a\) 及对应齐次方程组的通解。}


\analysis{由 \(r(B)=2\) 且 \(r(AB)=1\)，利用 Sylvester 不等式：
\[
r(AB)\ge r(A)+r(B)-3=r(A)-1.
\]
所以
\[
1\ge r(A)-1,\qquad r(A)\le2.
\]
又 \(r(AB)=1\)，故 \(r(A)\ge1\)。对题中
\[
A=\begin{pmatrix}
1&0&1\\
3&a&1\\
1&1&-1
\end{pmatrix},
\]
计算行列式：
\[
|A|=-2(a-1).
\]
要使 \(r(A)\le2\)，必须
\[
a=1.
\]
此时 \(r(A)=2\)，齐次方程 \(Ax=0\) 有一维非零解空间。代入 \(a=1\) 行化简，得
\[
x=t(-1,2,1)^T.
\]}
\answer{\[
a=1,\qquad x=t(-1,2,1)^T.
\]}
\end{solutionblock}

\begin{solutionblock}
\textbf{4.6 例 5}

\question{求非齐次线性方程组的通解，其增广矩阵为
\[
\left(
\begin{array}{cccc|c}
1&1&1&1&-1\\
4&3&5&-1&-1\\
2&1&3&-3&1
\end{array}
\right).
\]}


\analysis{非齐次方程组增广矩阵行化简为
\[
\left(
\begin{array}{cccc|c}
1&1&1&1&-1\\
4&3&5&-1&-1\\
2&1&3&-3&1
\end{array}
\right)
\sim
\left(
\begin{array}{cccc|c}
1&0&2&-4&2\\
0&1&-1&5&-3\\
0&0&0&0&0
\end{array}
\right).
\]
所以
\[
x_1+2x_3-4x_4=2,\qquad x_2-x_3+5x_4=-3.
\]
令 \(x_3=s,\ x_4=t\)，得
\[
x_1=-2s+4t+2,\qquad x_2=s-5t-3.
\]}
\answer{\[
x=(-2s+4t+2,\ s-5t-3,\ s,\ t)^T,\qquad s,t\in\mathbb{R}.
\]}
\end{solutionblock}

\begin{solutionblock}
\textbf{4.6 例 6}

\question{已知非齐次线性方程组有特解 \(x_0=(1,-1,1,-1)^T\)。求参数 \(\lambda,\mu\) 的关系及方程组通解。}


\analysis{把题给特解
\[
x_0=(1,-1,1,-1)^T
\]
代入方程组。第一、三式均给出
\[
-\lambda+\mu=0,
\]
故
\[
\lambda=\mu.
\]
在 \(\lambda=\mu\) 时，行化简可得解集为
\[
x=(-t,\frac{t-1}{2},\frac{1-t}{2},t)^T.
\]
取 \(t=-1\) 得题给特解 \(x_0\)。相应导出组的基础解向量可取
\[
(-2,1,-1,2)^T.
\]}
\answer{\[
\lambda=\mu,\qquad
x=(1,-1,1,-1)^T+c(-2,1,-1,2)^T.
\]}
\end{solutionblock}

\begin{solutionblock}
\textbf{4.6 例 7}

\question{设
\[
A=\begin{pmatrix}a&1\\1&0\end{pmatrix},
\qquad
B=\begin{pmatrix}0&1\\1&b\end{pmatrix}.
\]
讨论矩阵方程 \(AC-CA=B\) 何时有解，并求通解 \(C\)。}


\analysis{设
\[
C=\begin{pmatrix}x&y\\z&w\end{pmatrix}.
\]
由
\[
AC-CA=B
\]
逐项比较，得到线性方程组：
\[
az-y=0,\quad -ax+aw+y=1,
\]
\[
x-z-w=1,\quad -az-y=b.
\]
由第一式和第四式相减得
\[
b=0.
\]
再联立其余方程可得一致条件
\[
a=-1.
\]
当 \(a=-1,b=0\) 时，令 \(z=s,\ w=t\)，则
\[
x=s+t+1,\qquad y=-s.
\]}
\answer{存在这样的矩阵 \(C\) 当且仅当
\[
a=-1,\qquad b=0.
\]
此时
\[
C=\begin{pmatrix}
s+t+1&-s\\
s&t
\end{pmatrix},\qquad s,t\in\mathbb{R}.
\]}
\end{solutionblock}

\begin{solutionblock}
\textbf{4.6 例 8}

\question{已知 \(\beta_1,\beta_2\) 是非齐次方程组 \(Ax=b\) 的两个不同解，\(\alpha_1,\alpha_2\) 是导出组 \(Ax=0\) 的基础解系。判断题给通解形式中哪一个正确。}


\analysis{非齐次方程两个解 \(\beta_1,\beta_2\) 的平均
\[
\frac{\beta_1+\beta_2}{2}
\]
仍是非齐次方程 \(Ax=b\) 的一个特解；差
\[
\beta_1-\beta_2
\]
是对应齐次方程的解。

已知 \(\alpha_1,\alpha_2\) 是齐次方程的基础解系，所以非齐次方程通解应为
\[
\text{一个特解}+\text{齐次通解}.
\]
选项 B 中
\[
\frac{\beta_1+\beta_2}{2}+k_1\alpha_1+k_2(\alpha_1-\alpha_2)
\]
正是这种形式，且 \(\alpha_1,\alpha_1-\alpha_2\) 仍为齐次解空间的一组基。}
\answer{选 B。}
\end{solutionblock}

\begin{solutionblock}
\textbf{4.6 例 9}

\question{设三元非齐次线性方程组 \(Ax=\beta\) 有三个线性无关的解 \(\eta_1,\eta_2,\eta_3\)。判断其通解的正确表达式。}


\analysis{非齐次方程 \(Ax=\beta\) 的任意两个解之差都是导出组 \(Ax=0\) 的解。故
\[
\eta_1-\eta_2,\qquad \eta_3-\eta_1
\]
都是齐次方程的解。

题设 \(\eta_1,\eta_2,\eta_3\) 线性无关，而未知量个数为 \(3\)。非齐次方程有解且为非齐次方程，故导出组解空间维数为 \(2\)。于是通解可以写成
\[
\frac{\eta_2+\eta_3}{2}
+k_1(\eta_3-\eta_1)+k_2(\eta_2-\eta_1).
\]
其中 \((\eta_2+\eta_3)/2\) 是特解，后两项张成齐次解空间。}
\answer{选 C。}
\end{solutionblock}

\begin{solutionblock}
\textbf{4.6 例 10}

\question{已知 \(Ax=\beta\) 的通解为
\[
x=k(1,0,-3,2)^T+(1,-1,0,1)^T.
\]
令 \(B=(\alpha_2+\beta,\alpha_4,\alpha_3,\alpha_2,\alpha_1)\)，其中 \(A=(\alpha_1,\alpha_2,\alpha_3,\alpha_4)\)。求 \(Bx=\beta\) 的通解。}


\analysis{已知
\[
Ax=\beta
\]
的通解为
\[
x=k(1,0,-3,2)^T+(1,-1,0,1)^T.
\]
因此
\[
\beta=A(1,-1,0,1)^T=\alpha_1-\alpha_2+\alpha_4,
\]
且
\[
A(1,0,-3,2)^T=0.
\]
令
\[
B=(\alpha_2+\beta,\alpha_4,\alpha_3,\alpha_2,\alpha_1).
\]
注意
\[
\alpha_2+\beta=\alpha_1+\alpha_4.
\]
设 \(Bx=\beta\)，其中 \(x=(x_1,x_2,x_3,x_4,x_5)^T\)。按 \(A\) 的列向量收集系数，得到
\[
A(x_1+x_5,\ x_4,\ x_3,\ x_1+x_2)^T=\beta.
\]
故存在参数 \(k\)，使
\[
(x_1+x_5,\ x_4,\ x_3,\ x_1+x_2)^T
=(1,-1,0,1)^T+k(1,0,-3,2)^T.
\]
于是
\[
x_4=-1,\quad x_3=-3k,\quad x_1+x_5=1+k,\quad x_1+x_2=1+2k.
\]
令 \(x_1=s,\ k=t\)，得通解。}
\answer{\[
x=(s,\ 1+2t-s,\ -3t,\ -1,\ 1+t-s)^T,\qquad s,t\in\mathbb{R}.
\]}
\end{solutionblock}

\subsection{4.7 向量空间（仅数学一）}

\begin{solutionblock}
\textbf{4.7 例 1}

\question{设
\[
\alpha_1=(1,2,-1,0)^T,\quad
\alpha_2=(1,1,0,2)^T,\quad
\alpha_3=(2,1,1,a)^T.
\]
若 \(\alpha_1,\alpha_2,\alpha_3\) 张成空间维数为 \(2\)，求 \(a\)。}


\analysis{题中
\[
\alpha_1=(1,2,-1,0)^T,\quad
\alpha_2=(1,1,0,2)^T,\quad
\alpha_3=(2,1,1,a)^T.
\]
\(\alpha_1,\alpha_2\) 不成比例，所以二者线性无关。若三向量张成空间维数为 \(2\)，则必须有
\[
\alpha_3=x\alpha_1+y\alpha_2.
\]
由前三个分量：
\[
x+y=2,\quad 2x+y=1,\quad -x=1.
\]
解得
\[
x=-1,\qquad y=3.
\]
再看第四个分量：
\[
a=0\cdot x+2y=6.
\]}
\answer{\(a=6\)}
\end{solutionblock}

\begin{solutionblock}
\textbf{4.7 例 2}

\question{设向量 \(u=(2,0,0)^T\) 在基
\[
\alpha_1=(1,1,0)^T,\quad \alpha_2=(1,0,1)^T,\quad \alpha_3=(0,1,1)^T
\]
下的坐标为 \((x_1,x_2,x_3)^T\)。则
\[
x_1\alpha_1+x_2\alpha_2+x_3\alpha_3=u.
\]
解得
\[
x_1=1,\qquad x_2=1,\qquad x_3=-1.
\]}


\analysis{设向量 \(u=(2,0,0)^T\) 在基
\[
\alpha_1=(1,1,0)^T,\quad \alpha_2=(1,0,1)^T,\quad \alpha_3=(0,1,1)^T
\]
下的坐标为 \((x_1,x_2,x_3)^T\)。则
\[
x_1\alpha_1+x_2\alpha_2+x_3\alpha_3=u.
\]
解得
\[
x_1=1,\qquad x_2=1,\qquad x_3=-1.
\]}
\answer{\[
[u]_{\alpha}=(1,1,-1)^T.
\]}
\end{solutionblock}

\begin{solutionblock}
\textbf{4.7 例 3}

\question{设 \[
P_\alpha=(\alpha_1,\alpha_2)=
\begin{pmatrix}1&1\\0&-1\end{pmatrix},\qquad
P_\beta=(\beta_1,\beta_2)=
\begin{pmatrix}1&1\\1&2\end{pmatrix}.
\]
从基 \(\alpha\) 到基 \(\beta\) 的过渡矩阵 \(C\) 满足
\[
P_\beta=P_\alpha C.
\]}


\analysis{记
\[
P_\alpha=(\alpha_1,\alpha_2)=
\begin{pmatrix}1&1\\0&-1\end{pmatrix},\qquad
P_\beta=(\beta_1,\beta_2)=
\begin{pmatrix}1&1\\1&2\end{pmatrix}.
\]
从基 \(\alpha\) 到基 \(\beta\) 的过渡矩阵 \(C\) 满足
\[
P_\beta=P_\alpha C.
\]
因此
\[
C=P_\alpha^{-1}P_\beta
=\begin{pmatrix}
2&3\\
-1&-2
\end{pmatrix}.
\]}
\answer{\[
\begin{pmatrix}
2&3\\
-1&-2
\end{pmatrix}.
\]}
\end{solutionblock}

\begin{solutionblock}
\textbf{4.7 例 4}

\question{设 \(B\) 为秩为 \(2\) 的 \(5\times4\) 矩阵，且齐次方程 \(Bx=0\) 的解空间由 \(\alpha_1=(1,1,2,3)^T\)、\(\alpha_2=(-1,1,4,-1)^T\) 张成。求 \(Bx=0\) 的一组标准正交基础解系。}


\analysis{因为 \(B\) 是秩为 \(2\) 的 \(5\times4\) 矩阵，所以齐次方程 \(Bx=0\) 的解空间维数为
\[
4-r(B)=2.
\]
题中给出的三个解向量张成该解空间，其中
\[
\alpha_1=(1,1,2,3)^T,\qquad
\alpha_2=(-1,1,4,-1)^T
\]
线性无关，可先取为一组基。作 Gram-Schmidt 正交化：
\[
v_1=\alpha_1,
\]
\[
v_2=\alpha_2-\frac{\alpha_2^T\alpha_1}{\alpha_1^T\alpha_1}\alpha_1
=\alpha_2-\frac13\alpha_1.
\]
将 \(v_2\) 放大为整数向量，可取
\[
v_2=(-2,1,5,-3)^T.
\]
此时
\[
v_1^Tv_2=-2+1+10-9=0.
\]
再单位化：
\[
\|v_1\|=\sqrt{15},\qquad \|v_2\|=\sqrt{39}.
\]}
\answer{一个标准正交基为
\[
e_1=\frac1{\sqrt{15}}(1,1,2,3)^T,\qquad
e_2=\frac1{\sqrt{39}}(-2,1,5,-3)^T.
\]}
\end{solutionblock}

\chapter{矩阵相似理论}

\section{逐题解析}

\subsection{5.1 特征值与特征向量}

\begin{solutionblock}
\textbf{5.1 例 1}

\question{求矩阵 \(A=\begin{pmatrix}1&2&-3\\-1&4&-3\\1&-2&5\end{pmatrix}\) 的全部特征值及其对应的特征向量。}



\analysis{设
\[
A=\begin{pmatrix}
1&2&-3\\
-1&4&-3\\
1&-2&5
\end{pmatrix}.
\]
计算特征多项式：
\[
|\lambda E-A|
=\begin{vmatrix}
\lambda-1&-2&3\\
1&\lambda-4&3\\
-1&2&\lambda-5
\end{vmatrix}
=(\lambda-6)(\lambda-2)^2.
\]
所以特征值为 \(\lambda=2,2,6\)。

当 \(\lambda=2\) 时，
\[
A-2E=\begin{pmatrix}
-1&2&-3\\
-1&2&-3\\
1&-2&3
\end{pmatrix},
\quad -x_1+2x_2-3x_3=0.
\]
令 \(x_2=s,x_3=t\)，得 \(x_1=2s-3t\)，故
\[
\xi=s(2,1,0)^T+t(-3,0,1)^T.
\]
当 \(\lambda=6\) 时，解 \((A-6E)x=0\)，可得
\[
\xi=k(-1,-1,1)^T.
\]}
\answer{\[
\lambda=2:\ k_1(2,1,0)^T+k_2(-3,0,1)^T,\qquad
\lambda=6:\ k(-1,-1,1)^T,
\]
其中对应特征向量要求系数不全为零。}
\examnote{二重特征值要继续求特征子空间维数，不能只写一个特征向量。}
\end{solutionblock}

\begin{solutionblock}
\textbf{5.1 例 2}

\question{分别求下列矩阵的特征值：\((1)\ \begin{pmatrix}1&0&0\\0&2&0\\6&-2&3\end{pmatrix}\)；\((2)\ \begin{pmatrix}1&2&3\\2&4&6\\3&6&9\end{pmatrix}\)。}



\analysis{第 (1) 个矩阵
\[
\begin{pmatrix}
1&0&0\\
0&2&0\\
6&-2&3
\end{pmatrix}
\]
是下三角矩阵，特征值就是主对角线元素：
\[
1,\ 2,\ 3.
\]

第 (2) 个矩阵
\[
\begin{pmatrix}
1&2&3\\
2&4&6\\
3&6&9
\end{pmatrix}
=\alpha\alpha^T,\qquad \alpha=(1,2,3)^T.
\]
秩一矩阵 \(\alpha\alpha^T\) 的非零特征值为
\[
\alpha^T\alpha=1^2+2^2+3^2=14,
\]
其余特征值为 \(0\)。}
\answer{第 (1) 个矩阵的特征值为 \(1,2,3\)；第 (2) 个矩阵的特征值为 \(14,0,0\)。}
\end{solutionblock}

\begin{solutionblock}
\textbf{5.1 例 3}

\question{已知 \(0\) 是矩阵 \(A=\begin{pmatrix}1&0&1\\0&2&0\\1&0&a\end{pmatrix}\) 的一个特征值，求 \(a\)，并求此时 \(A\) 的全部特征值。}



\analysis{\(0\) 是
\[
A=\begin{pmatrix}
1&0&1\\
0&2&0\\
1&0&a
\end{pmatrix}
\]
的特征值，等价于 \(|A|=0\)。计算
\[
|A|=2\begin{vmatrix}1&1\\1&a\end{vmatrix}=2(a-1),
\]
所以 \(a=1\)。此时
\[
A=\begin{pmatrix}
1&0&1\\
0&2&0\\
1&0&1
\end{pmatrix}.
\]
第一、三坐标构成的块为
\[
\begin{pmatrix}1&1\\1&1\end{pmatrix},
\]
特征值为 \(0,2\)，中间坐标给出特征值 \(2\)。}
\answer{\[
a=1,\qquad A\text{ 的全部特征值为 }0,2,2.
\]}
\end{solutionblock}

\begin{solutionblock}
\textbf{5.1 例 4}

\question{设矩阵 \(A\) 可逆，且 \(A^{-1}\) 的特征值为 \(1,-\frac12,\frac13\)，求 \(A\)、\((2A)^*\)、\(A^2+4A+E\) 的特征值。}



\analysis{\(A^{-1}\) 的特征值为 \(1,-\frac12,\frac13\)，所以 \(A\) 的特征值为其倒数：
\[
1,\ -2,\ 3.
\]

对于 \(A\)，若特征值为 \(\lambda_i\)，则 \(A^*\) 的对应特征值为
\[
\frac{|A|}{\lambda_i}.
\]
这里
\[
|A|=1\cdot(-2)\cdot3=-6,
\]
故 \(A^*\) 的特征值为
\[
-6,\ 3,\ -2.
\]
三阶矩阵满足
\[
(2A)^*=2^{3-1}A^*=4A^*,
\]
所以 \((2A)^*\) 的特征值为
\[
-24,\ 12,\ -8.
\]

对多项式矩阵 \(A^2+4A+E\)，特征值由函数
\[
f(x)=x^2+4x+1
\]
作用在 \(A\) 的特征值上得到：
\[
f(1)=6,\qquad f(-2)=-3,\qquad f(3)=22.
\]}
\answer{\[
A:\ 1,-2,3;\qquad
(2A)^*:\ -24,12,-8;\qquad
A^2+4A+E:\ 6,-3,22.
\]}
\end{solutionblock}

\begin{solutionblock}
\textbf{5.1 例 5}

\question{设三阶矩阵 \(A\) 满足 \(|A-E|=|A+2E|=|2A+3E|=0\)，求 \(|2A^*-3E|\)。}



\analysis{由
\[
|A-E|=0,\quad |A+2E|=0,\quad |2A+3E|=0
\]
知 \(A\) 的三个特征值分别为
\[
1,\ -2,\ -\frac32.
\]
故
\[
|A|=1\cdot(-2)\cdot\left(-\frac32\right)=3.
\]
因为 \(A\) 可逆，\(A^*\) 的特征值为 \(|A|/\lambda_i\)，即
\[
3,\ -\frac32,\ -2.
\]
于是 \(2A^*-3E\) 的特征值为
\[
2\cdot3-3=3,\qquad
2\cdot\left(-\frac32\right)-3=-6,\qquad
2\cdot(-2)-3=-7.
\]
所求行列式为三个特征值之积：
\[
|2A^*-3E|=3\cdot(-6)\cdot(-7)=126.
\]}
\answer{\(126\)}
\end{solutionblock}

\begin{solutionblock}
\textbf{5.1 例 6}

\question{设 \(A\) 为 \(n\) 阶可逆矩阵，\(A\) 的每行元素之和均为 \(a\,(a\ne0)\)，且 \(|A|=1\)。求 \(A^*\) 的每行元素之和。}



\analysis{设 \(\mathbf{1}=(1,1,\ldots,1)^T\)。每行元素之和均为 \(a\)，等价于
\[
A\mathbf{1}=a\mathbf{1}.
\]
由于 \(A\) 可逆且 \(a\ne0\)，两边左乘 \(A^{-1}\)，得
\[
A^{-1}\mathbf{1}=\frac1a\mathbf{1}.
\]
又 \(|A|=1\)，所以
\[
A^*=|A|A^{-1}=A^{-1}.
\]
矩阵乘以 \(\mathbf{1}\) 得到的正是各行元素之和，因此 \(A^*\) 的每行元素之和均为 \(\frac1a\)。}
\answer{\(\dfrac1a\)}
\end{solutionblock}

\begin{solutionblock}
\textbf{5.1 例 7}

\question{设
\[
A=\begin{pmatrix}
a&-1&c\\
5&b&3\\
1-c&0&-a
\end{pmatrix},
\qquad |A|=-1.
\]
若 \(A\) 的伴随矩阵 \(A^*\) 有特征值 \(\lambda_0\)，且属于 \(\lambda_0\) 的特征向量为 \(\alpha=(-1,-1,1)^T\)，求 \(a,b,c\) 及 \(\lambda_0\)。}



\analysis{设
\[
A=\begin{pmatrix}
a&-1&c\\
5&b&3\\
1-c&0&-a
\end{pmatrix},\qquad \alpha=(-1,-1,1)^T.
\]
因为 \(\alpha\) 是 \(A^*\) 属于 \(\lambda_0\) 的特征向量，且 \(|A|=-1\ne0\)，所以 \(\alpha\) 也是 \(A\) 的特征向量。设
\[
A\alpha=\mu\alpha.
\]
直接计算
\[
A\alpha=
\begin{pmatrix}
-a+1+c\\
-5-b+3\\
-1+c-a
\end{pmatrix}
=
\begin{pmatrix}
-a+c+1\\
-b-2\\
c-a-1
\end{pmatrix}.
\]
与 \(\mu(-1,-1,1)^T\) 比较，得
\[
-a+c+1=-\mu,\qquad -b-2=-\mu,\qquad c-a-1=\mu.
\]
第一式与第三式给出
\[
-a+c+1=-(c-a-1),
\]
恒成立；再由第二、第三式得
\[
\mu=c-a-1=b+2.
\]
即
\[
b=c-a-3.
\]
再代入 \(|A|=-1\)，可解得
\[
a=2,\qquad b=-3,\qquad c=2,\qquad \mu=-1.
\]
最后利用 \(A^*=|A|A^{-1}\)：若 \(A\alpha=\mu\alpha\)，则
\[
A^*\alpha=\frac{|A|}{\mu}\alpha=\frac{-1}{-1}\alpha=\alpha.
\]
所以 \(\lambda_0=1\)。}
\answer{\[
a=2,\qquad b=-3,\qquad c=2,\qquad \lambda_0=1.
\]}
\examnote{可逆时 \(A\) 与 \(A^*\) 的特征向量相同，特征值按 \(\lambda\mapsto |A|/\lambda\) 转换。}
\end{solutionblock}

\begin{solutionblock}
\textbf{5.1 例 8}

\question{设 \(A\) 为 \(n\) 阶实对称矩阵，\(P\) 为 \(n\) 阶可逆矩阵。已知 \(n\) 维列向量 \(\alpha\) 是矩阵 \(A\) 属于特征值 \(\lambda\) 的特征向量，判断矩阵 \((P^{-1}AP)^T\) 属于特征值 \(\lambda\) 的特征向量：
\[
\text{A. }P^{-1}\alpha\qquad
\text{B. }P^T\alpha\qquad
\text{C. }P\alpha\qquad
\text{D. }(P^{-1})^T\alpha.
\]}



\analysis{设 \(A\alpha=\lambda\alpha\)。因为 \(A\) 为实对称矩阵，故 \(A^T=A\)。记
\[
C=(P^{-1}AP)^T=P^TA^T(P^{-1})^T=P^TA(P^{-1})^T.
\]
考察 \(P^T\alpha\)：
\[
C(P^T\alpha)
=P^TA(P^{-1})^TP^T\alpha
=P^TA\alpha
=\lambda P^T\alpha.
\]
故 \(P^T\alpha\) 是 \(C\) 属于 \(\lambda\) 的特征向量。}
\answer{选 B，\(P^T\alpha\)。}
\end{solutionblock}

\begin{solutionblock}
\textbf{5.1 例 9}

\question{若 \(\alpha_1\) 是矩阵 \(A\) 属于特征值 \(\lambda=1\) 的特征向量，\(\alpha_2,\alpha_3\) 是矩阵 \(A\) 属于特征值 \(\lambda=2\) 的两个线性无关特征向量。设
\[
AP=P\begin{pmatrix}
1&0&0\\
0&2&0\\
0&0&2
\end{pmatrix},
\]
判断矩阵 \(P\) 不可以是哪一个：
\[
\text{A. }(\alpha_1,-2\alpha_2,\alpha_3),\quad
\text{B. }(\alpha_1,\alpha_2+\alpha_3,\alpha_2-\alpha_3),
\]
\[
\text{C. }(\alpha_1,\alpha_3,\alpha_2),\quad
\text{D. }(\alpha_1+\alpha_2,\alpha_1-\alpha_2,\alpha_3).
\]}



\analysis{由
\[
AP=P\begin{pmatrix}
1&0&0\\
0&2&0\\
0&0&2
\end{pmatrix}
\]
可知 \(P\) 的第一列必须是属于 \(\lambda=1\) 的特征向量，第二、三列必须是属于 \(\lambda=2\) 的两个线性无关特征向量。

选项 A、B、C 中第一列仍为 \(\alpha_1\)，第二、三列均是 \(\alpha_2,\alpha_3\) 的线性无关组合，所以可以作为 \(P\)。

选项 D 第一列为 \(\alpha_1+\alpha_2\)。由于
\[
A(\alpha_1+\alpha_2)=\alpha_1+2\alpha_2
\]
不等于 \(1(\alpha_1+\alpha_2)\)，也不等于 \(2(\alpha_1+\alpha_2)\)，故它不是单一特征值对应的特征向量，不能作为第一列。}
\answer{选 D。}
\end{solutionblock}

\subsection{5.2 秩为 1 矩阵专题}

\begin{solutionblock}
\textbf{5.2 例 1}

\question{设 \(\alpha=(1,2,3)^T\)，\(\beta=(1,\frac12,\frac13)^T\)，\(A=\alpha\beta^T\)。求 \(A^n\)。}



\analysis{设
\[
\alpha=(1,2,3)^T,\qquad \beta=\left(1,\frac12,\frac13\right)^T,\qquad A=\alpha\beta^T.
\]
秩一矩阵满足
\[
A^2=(\alpha\beta^T)(\alpha\beta^T)
=\alpha(\beta^T\alpha)\beta^T
=(\beta^T\alpha)A.
\]
这里
\[
\beta^T\alpha=1+1+1=3.
\]
因此
\[
A^n=(\beta^T\alpha)^{n-1}A=3^{n-1}A,\qquad n\ge1.
\]}
\answer{\[
A^n=3^{n-1}\alpha\beta^T=3^{n-1}A.
\]}
\end{solutionblock}

\begin{solutionblock}
\textbf{5.2 例 2}

\question{设 \(A=E+\alpha\beta^T\)，且 \(\alpha^T\beta=2\)。求 \(A^{-1}\)。}



\analysis{令 \(A=E+\alpha\beta^T\)。秩一修正公式为
\[
(E+\alpha\beta^T)^{-1}
=E-\frac{\alpha\beta^T}{1+\beta^T\alpha},
\]
前提是 \(1+\beta^T\alpha\ne0\)。题中
\[
\alpha^T\beta=\beta^T\alpha=2,
\]
故
\[
1+\beta^T\alpha=3\ne0,
\]
于是 \(A\) 可逆，且
\[
A^{-1}=E-\frac13\alpha\beta^T.
\]}
\answer{\[
A^{-1}=E-\frac13\alpha\beta^T.
\]}
\examnote{考研中遇到 \(E+\alpha\beta^T\)，优先想到秩一修正公式；也可设逆为 \(E+k\alpha\beta^T\) 后比较系数求 \(k\)。}
\end{solutionblock}

\begin{solutionblock}
\textbf{5.2 例 3}

\question{设 \(\alpha\beta^T\) 为三阶矩阵，且 \(\alpha\beta^T\) 与
\[
\begin{pmatrix}2&0&0\\0&0&0\\0&0&0\end{pmatrix}
\]
相似。求 \(\beta^T\alpha\)。}



\analysis{\(\alpha\beta^T\) 为三阶秩一矩阵，非零特征值为
\[
\beta^T\alpha.
\]
它与
\[
\begin{pmatrix}
2&0&0\\
0&0&0\\
0&0&0
\end{pmatrix}
\]
相似，故两者特征值相同，为 \(2,0,0\)。因此
\[
\beta^T\alpha=2.
\]}
\answer{\(2\)}
\end{solutionblock}

\begin{solutionblock}
\textbf{5.2 例 4}

\question{设
\[
A=\begin{pmatrix}2&1&3\\4&2&6\\6&3&9\end{pmatrix},
\qquad
B=A+E.
\]
求 \(A\) 与 \(B\) 的特征值和对应特征向量。}



\analysis{设
\[
A=\begin{pmatrix}2&1&3\\4&2&6\\6&3&9\end{pmatrix}
=\alpha\beta^T,\quad
\alpha=(1,2,3)^T,\quad \beta=(2,1,3)^T.
\]
则
\[
\beta^T\alpha=2+2+9=13.
\]
所以 \(A\) 的特征值为 \(13,0,0\)。对应 \(\lambda=13\) 的特征向量可取 \(\alpha=(1,2,3)^T\)；对应 \(\lambda=0\) 的特征向量满足
\[
\beta^Tx=2x_1+x_2+3x_3=0,
\]
即
\[
x=s(-1,2,0)^T+t(-3,0,2)^T.
\]

对
\[
B=\begin{pmatrix}3&1&3\\4&3&6\\6&3&10\end{pmatrix}
=A+E.
\]
若 \(A\) 的特征值为 \(\lambda\)，则 \(A+E\) 的特征值为 \(\lambda+1\)。因此 \(B\) 的特征值为
\[
14,1,1.
\]
特征向量不变：\(\lambda=14\) 对应 \((1,2,3)^T\)；\(\lambda=1\) 对应
\[
2x_1+x_2+3x_3=0.
\]}
\answer{\[
A:\ \lambda=13,\ \xi=k(1,2,3)^T;\quad
\lambda=0,\ \xi=s(-1,2,0)^T+t(-3,0,2)^T.
\]
\[
B:\ \lambda=14,\ \xi=k(1,2,3)^T;\quad
\lambda=1,\ \xi=s(-1,2,0)^T+t(-3,0,2)^T.
\]}
\end{solutionblock}

\subsection{5.3 矩阵相似}

\begin{solutionblock}
\textbf{5.3 例 1}

\question{已知矩阵
\[
A=\begin{pmatrix}-2&-2&1\\2&x&-2\\0&0&-2\end{pmatrix},
\qquad
B=\begin{pmatrix}2&1&0\\0&-1&0\\0&0&y\end{pmatrix}
\]
相似，求 \(x,y\)。}



\analysis{相似矩阵有相同的特征多项式。对
\[
A=\begin{pmatrix}
-2&-2&1\\
2&x&-2\\
0&0&-2
\end{pmatrix},
\quad
B=\begin{pmatrix}
2&1&0\\
0&-1&0\\
0&0&y
\end{pmatrix},
\]
分别计算
\[
|\lambda E-A|
=-(\lambda+2)(-\lambda^2+\lambda x-2\lambda+2x-4),
\]
\[
|\lambda E-B|=(\lambda-2)(\lambda+1)(\lambda-y).
\]
比较系数可得
\[
x=3,\qquad y=-2.
\]
此时两个矩阵的特征值为 \(2,-1,-2\)，互异，故均可对角化并相似于同一个对角矩阵。}
\answer{\[
x=3,\qquad y=-2.
\]}
\end{solutionblock}

\begin{solutionblock}
\textbf{5.3 例 2}

\question{设
\[
B=\begin{pmatrix}0&0&1\\0&1&0\\1&0&0\end{pmatrix},
\qquad A\sim B.
\]
求 \(r(A-2E)+r(A-E)\)。}



\analysis{题设 \(A\sim B\)，其中
\[
B=\begin{pmatrix}
0&0&1\\
0&1&0\\
1&0&0
\end{pmatrix}.
\]
相似变换不改变矩阵 \(A-\lambda E\) 的秩，因为
\[
A=P^{-1}BP
\quad\Longrightarrow\quad
A-\lambda E=P^{-1}(B-\lambda E)P.
\]
矩阵 \(B\) 的特征值为 \(1,1,-1\)。所以 \(2\) 不是特征值，
\[
r(A-2E)=r(B-2E)=3.
\]
而 \(\lambda=1\) 的几何重数为 \(2\)，故
\[
\dim\ker(B-E)=2,\qquad r(B-E)=3-2=1.
\]
于是
\[
r(A-2E)+r(A-E)=3+1=4.
\]}
\answer{选 C，\(4\)。}
\end{solutionblock}

\begin{solutionblock}
\textbf{5.3 例 3}

\question{设 \(A=\alpha\alpha^T\)，其中 \(\alpha=(1,2,-1)^T\)，且 \(A\sim B\)。求 \(|(B+E)^*|\)。}



\analysis{已知
\[
\alpha=(1,2,-1)^T,\qquad A=\alpha\alpha^T.
\]
秩一矩阵 \(A\) 的非零特征值为
\[
\alpha^T\alpha=1+4+1=6,
\]
故 \(A\) 的特征值为 \(6,0,0\)。由于 \(A\sim B\)，\(B\) 的特征值也为 \(6,0,0\)，从而 \(B+E\) 的特征值为
\[
7,1,1.
\]
于是
\[
|B+E|=7.
\]
三阶矩阵 \(C\) 满足
\[
|C^*|=|C|^{3-1}=|C|^2.
\]
取 \(C=B+E\)，得
\[
|(B+E)^*|=7^2=49.
\]}
\answer{\(49\)}
\end{solutionblock}

\begin{solutionblock}
\textbf{5.3 例 4}

\question{若 \(A\) 与 \(B\) 相似，判断有关特征矩阵、可逆性、特征向量和可对角化性的命题中哪一个一定正确。}



\analysis{若 \(A\sim B\)，则 \(A=P^{-1}BP\)。因此
\[
A-\lambda E=P^{-1}(B-\lambda E)P.
\]
取 \(\lambda=0\) 可知 \(A\) 与 \(B\) 同时可逆或同时不可逆，所以选项 B 正确。

选项 A 写成矩阵相等，过强；相似只能推出 \(\lambda E-A\sim \lambda E-B\)，不能推出二者相等。选项 C 错在特征向量依赖具体坐标基，通常不会相同。选项 D 错在相似矩阵不一定可对角化。}
\answer{选 B。}
\end{solutionblock}

\begin{solutionblock}
\textbf{5.3 例 5}

\question{设 \(A,B\) 均可逆且 \(A\sim B\)，判断 \(A^T\sim B^T\)、\(A+A^T\sim B+B^T\)、\(A^{-1}\sim B^{-1}\)、\(A+A^{-1}\sim B+B^{-1}\) 中哪些成立。}



\analysis{设 \(A=P^{-1}BP\)。则
\[
A^T=P^TB^T(P^{-1})^T,
\]
所以 \(A^T\sim B^T\)，选项 A 正确。

若 \(A,B\) 可逆，则
\[
A^{-1}=P^{-1}B^{-1}P,
\]
所以 \(A^{-1}\sim B^{-1}\)，选项 C 正确。进一步
\[
A+A^{-1}=P^{-1}(B+B^{-1})P,
\]
故选项 D 正确。

但 \(A+A^T\) 中同时出现 \(P^{-1}BP\) 与 \(P^TB^T(P^{-1})^T\)，一般不能合并成同一个相似变换，所以不一定与 \(B+B^T\) 相似。}
\answer{选 B。}
\examnote{“多项式保持相似”要求表达式只由 \(A\) 的加减乘幂组成；一旦混入 \(A^T\)，就要格外小心。}
\end{solutionblock}

\begin{solutionblock}
\textbf{5.3 例 6}

\question{判断矩阵 \(\begin{pmatrix}1&1&0\\0&1&1\\0&0&1\end{pmatrix}\) 与给定四个三阶矩阵中的哪一个相似。}



\analysis{目标矩阵
\[
J=\begin{pmatrix}
1&1&0\\
0&1&1\\
0&0&1
\end{pmatrix}
\]
是三阶 Jordan 块。判断相似时，看
\[
N=J-E
\]
的秩结构：
\[
r(N)=2,\qquad r(N^2)=1.
\]
四个选项分别计算 \(M-E\) 与 \((M-E)^2\) 的秩，只有选项 A 满足
\[
r(M-E)=2,\qquad r((M-E)^2)=1.
\]
其余选项的 \(M-E\) 秩为 \(1\) 或平方为零，对应 Jordan 块结构不同。}
\answer{选 A。}
\end{solutionblock}

\begin{solutionblock}
\textbf{5.3 例 7}

\question{判断下列条件中哪一个是 \(A\sim B\) 的充分条件：\(A^*\sim B^*\)、\(A^T\sim B^T\)、\(f(A)\sim f(B)\)、\(AB\sim BA\)。}



\analysis{若 \(A\sim B\)，则必有 \(A^T\sim B^T\)；反过来，若 \(A^T\sim B^T\)，即存在可逆矩阵 \(P\) 使
\[
A^T=P^{-1}B^TP,
\]
两边转置得
\[
A=P^TB(P^{-1})^T,
\]
所以 \(A\sim B\)。因此 \(A^T\sim B^T\) 是 \(A\sim B\) 的充分条件。

选项 A 中伴随矩阵相似不能反推原矩阵相似；选项 C 中 \(f(A)\sim f(B)\) 信息可能丢失；选项 D 中 \(AB\sim BA\) 即使成立，也不能推出 \(A\sim B\)。}
\answer{选 B。}
\end{solutionblock}

\subsection{5.4 矩阵相似对角化}

\begin{solutionblock}
\textbf{5.4 例 1}

\question{设三阶矩阵 \(A\) 的特征值为 \(0,1,2\)，令 \(B=A^3-2A^2\)，求 \(r(B)\)。}



\analysis{设 \(f(x)=x^3-2x^2=x^2(x-2)\)，则
\[
B=A^3-2A^2=f(A).
\]
已知 \(A\) 的特征值为 \(0,1,2\)，且三个特征值互异，所以 \(A\) 可对角化。于是 \(B=f(A)\) 的特征值为
\[
f(0)=0,\qquad f(1)=-1,\qquad f(2)=0.
\]
\(B\) 只有一个非零特征值，故
\[
r(B)=1.
\]}
\answer{\(1\)}
\end{solutionblock}

\begin{solutionblock}
\textbf{5.4 例 2}

\question{设 \(\alpha=(1,1,1)^T\)，\(\beta=(1,0,k)^T\)，且 \(\alpha\beta^T\) 与 \(\operatorname{diag}(4,0,0)\) 相似，求 \(k\)。}



\analysis{秩一矩阵 \(\alpha\beta^T\) 的非零特征值为
\[
\beta^T\alpha.
\]
题设
\[
\alpha=(1,1,1)^T,\qquad \beta=(1,0,k)^T,
\]
所以
\[
\beta^T\alpha=1+0+k=1+k.
\]
又 \(\alpha\beta^T\) 与 \(\operatorname{diag}(4,0,0)\) 相似，故其非零特征值为 \(4\)。于是
\[
1+k=4,\qquad k=3.
\]}
\answer{\(3\)}
\end{solutionblock}

\begin{solutionblock}
\textbf{5.4 例 3}

\question{判断给定四个矩阵中哪一个不能相似对角化。}



\analysis{实对称矩阵一定可以正交相似对角化，所以选项 A 可对角化。

选项 B 为三角矩阵，特征值 \(1,2,3\) 互异，故可对角化。

选项 C 是秩一矩阵，特征值为 \(7,0,0\)，且 \(\lambda=0\) 的特征子空间维数为 \(2\)，也可对角化。

选项 D 的特征多项式为
\[
\lambda^2(\lambda-1),
\]
但
\[
r(D-0E)=2,\qquad \dim\ker D=1,
\]
\(\lambda=0\) 的几何重数小于代数重数 \(2\)，故不能对角化。}
\answer{选 D。}
\end{solutionblock}

\begin{solutionblock}
\textbf{5.4 例 4}

\question{设 \(A=\begin{pmatrix}4&6&0\\-3&-5&0\\-3&-6&1\end{pmatrix}\)，求可逆矩阵 \(P\)，使 \(P^{-1}AP\) 为对角矩阵。}



\analysis{设
\[
A=\begin{pmatrix}
4&6&0\\
-3&-5&0\\
-3&-6&1
\end{pmatrix}.
\]
计算
\[
|\lambda E-A|=(\lambda+2)(\lambda-1)^2.
\]
当 \(\lambda=-2\) 时，特征向量可取
\[
\xi_1=(-1,1,1)^T.
\]
当 \(\lambda=1\) 时，解 \((A-E)x=0\)，可取两个线性无关特征向量
\[
\xi_2=(-2,1,0)^T,\qquad \xi_3=(0,0,1)^T.
\]
令
\[
P=(\xi_1,\xi_2,\xi_3)
=\begin{pmatrix}
-1&-2&0\\
1&1&0\\
1&0&1
\end{pmatrix}.
\]
则 \(|P|\ne0\)，并且
\[
P^{-1}AP=\begin{pmatrix}
-2&0&0\\
0&1&0\\
0&0&1
\end{pmatrix}.
\]}
\answer{可取
\[
P=\begin{pmatrix}
-1&-2&0\\
1&1&0\\
1&0&1
\end{pmatrix},
\quad
P^{-1}AP=\operatorname{diag}(-2,1,1).
\]}
\end{solutionblock}

\begin{solutionblock}
\textbf{5.4 例 5}

\question{设 \(A=\begin{pmatrix}1&-1&1\\x&4&y\\-3&-3&5\end{pmatrix}\)。已知 \(A\) 有三个线性无关特征向量，且 \(\lambda=2\) 为二重特征值，求 \(x,y\)，并求 \(P\) 使 \(P^{-1}AP\) 为对角矩阵。}



\analysis{设
\[
A=\begin{pmatrix}
1&-1&1\\
x&4&y\\
-3&-3&5
\end{pmatrix}.
\]
由题意，\(A\) 有三个线性无关的特征向量，且 \(\lambda=2\) 是二重特征值，因此特征多项式应为
\[
(\lambda-2)^2(\lambda-\lambda_3).
\]
直接计算可得
\[
|\lambda E-A|=(\lambda-2)(\lambda^2-8\lambda+x+3y+16).
\]
若 \(\lambda=2\) 为二重根，则
\[
4-16+x+3y+16=0,\qquad x+3y+4=0. \tag{1}
\]
但仅有二重根还不够，还要求 \(\lambda=2\) 的特征子空间维数为 \(2\)。由
\[
A-2E=\begin{pmatrix}
-1&-1&1\\
x&2&y\\
-3&-3&3
\end{pmatrix}
\]
可知第一行与第三行成比例，要使秩为 \(1\)，第二行也必须与第一行成比例。结合 (1) 得
\[
x=2,\qquad y=-2.
\]
此时另一个特征值由迹确定：
\[
1+4+5=10=2+2+\lambda_3,\qquad \lambda_3=6.
\]
对应 \(\lambda=2\) 的特征向量可取
\[
\xi_1=(-1,1,0)^T,\qquad \xi_2=(1,0,1)^T,
\]
对应 \(\lambda=6\) 的特征向量可取
\[
\xi_3=(1,-2,3)^T.
\]
令
\[
P=(\xi_1,\xi_2,\xi_3)
=\begin{pmatrix}
-1&1&1\\
1&0&-2\\
0&1&3
\end{pmatrix}.
\]
则
\[
P^{-1}AP=\operatorname{diag}(2,2,6).
\]}
\answer{\[
x=2,\qquad y=-2,\qquad
P=\begin{pmatrix}
-1&1&1\\
1&0&-2\\
0&1&3
\end{pmatrix}.
\]}
\examnote{本题的关键陷阱是：二重特征值不代表一定能对角化，还要检查该特征值对应的线性无关特征向量个数。}
\end{solutionblock}

\begin{solutionblock}
\textbf{5.4 例 6}

\question{设 \(A=\begin{pmatrix}1&2&-3\\-1&4&-3\\1&a&5\end{pmatrix}\) 有二重特征值，求 \(a\)，并判断 \(A\) 是否可相似对角化。}



\analysis{设
\[
A=\begin{pmatrix}
1&2&-3\\
-1&4&-3\\
1&a&5
\end{pmatrix}.
\]
计算特征多项式：
\[
|\lambda E-A|=(\lambda-2)(\lambda^2-8\lambda+3a+18).
\]
若有一个二重根，则判别式为零或 \(\lambda=2\) 也是二次因子的根。等价计算可得
\[
a=-2\quad\text{或}\quad a=-\frac23.
\]

当 \(a=-2\) 时，
\[
|\lambda E-A|=(\lambda-6)(\lambda-2)^2,
\]
且 \(\lambda=2\) 的特征子空间维数为 \(2\)，故 \(A\) 可相似对角化。

当 \(a=-\frac23\) 时，
\[
|\lambda E-A|=(\lambda-4)^2(\lambda-2),
\]
但 \(\lambda=4\) 的特征子空间维数为 \(1\)，小于其代数重数 \(2\)，故不能相似对角化。}
\answer{\[
a=-2\ \text{或}\ -\frac23.
\]
其中 \(a=-2\) 时 \(A\) 可相似对角化；\(a=-\frac23\) 时不可相似对角化。}
\end{solutionblock}

\begin{solutionblock}
\textbf{5.4 例 7}

\question{设三阶矩阵 \(A\) 对应线性方程组 \(Ax=b\) 的通解为 \(k_1(-1,1,0)^T+k_2(2,0,1)^T+(1,1,-2)^T\)，其中 \(b=(6,6,-12)^T\)，求 \(A\)。}



\analysis{方程组 \(Ax=b\) 的通解为
\[
x=k_1(-1,1,0)^T+k_2(2,0,1)^T+(1,1,-2)^T.
\]
因此齐次方程组 \(Ax=0\) 的基础解系为
\[
\eta_1=(-1,1,0)^T,\qquad \eta_2=(2,0,1)^T.
\]
又特解 \(x_0=(1,1,-2)^T\) 满足
\[
Ax_0=b=(6,6,-12)^T=6x_0.
\]
所以 \(A\) 在基
\[
P=(\eta_1,\eta_2,x_0)
=\begin{pmatrix}
-1&2&1\\
1&0&1\\
0&1&-2
\end{pmatrix}
\]
下对应对角矩阵
\[
\operatorname{diag}(0,0,6).
\]
即
\[
AP=P\operatorname{diag}(0,0,6),
\quad
A=P\operatorname{diag}(0,0,6)P^{-1}.
\]
计算得
\[
A=\begin{pmatrix}
2&2&-2\\
2&2&-2\\
-4&-4&4
\end{pmatrix}.
\]}
\answer{\[
A=\begin{pmatrix}
2&2&-2\\
2&2&-2\\
-4&-4&4
\end{pmatrix}.
\]}
\end{solutionblock}

\begin{solutionblock}
\textbf{5.4 例 8}

\question{设 \(A=\begin{pmatrix}1&4&-2\\0&-1&0\\1&2&-2\end{pmatrix}\)，求 \(A^{100}\)。}



\analysis{设
\[
A=\begin{pmatrix}
1&4&-2\\
0&-1&0\\
1&2&-2
\end{pmatrix}.
\]
其特征值为 \(0,-1,-1\)，且可对角化。因为 \(100\) 为偶数，
\[
0^{100}=0,\qquad (-1)^{100}=1.
\]
于是 \(A^{100}\) 等于将 \(A\) 的特征值 \(0,-1\) 分别替换为 \(0,1\) 后得到的矩阵。直接由最小多项式也可看出
\[
A^{100}=A^2.
\]
计算
\[
A^2=
\begin{pmatrix}
-1&-4&2\\
0&1&0\\
-1&-2&2
\end{pmatrix}.
\]}
\answer{\[
A^{100}=
\begin{pmatrix}
-1&-4&2\\
0&1&0\\
-1&-2&2
\end{pmatrix}.
\]}
\end{solutionblock}

\begin{solutionblock}
\textbf{5.4 例 9}

\question{设 \(A=\begin{pmatrix}-2&-2&1\\2&x&-2\\0&0&-2\end{pmatrix}\)，\(B=\begin{pmatrix}2&1&0\\0&-1&0\\0&0&y\end{pmatrix}\)。求可逆矩阵 \(P\)，使 \(P^{-1}AP=B\)。}



\analysis{由 5.3 例 1 已知相似时
\[
x=3,\qquad y=-2.
\]
于是
\[
A=\begin{pmatrix}
-2&-2&1\\
2&3&-2\\
0&0&-2
\end{pmatrix},\qquad
B=\begin{pmatrix}
2&1&0\\
0&-1&0\\
0&0&-2
\end{pmatrix}.
\]
要求 \(P^{-1}AP=B\)，等价于
\[
AP=PB.
\]
若 \(P=(p_1,p_2,p_3)\)，则由 \(PB\) 的列关系得
\[
Ap_1=2p_1,\qquad Ap_2=p_1-p_2,\qquad Ap_3=-2p_3.
\]
可取
\[
p_1=(-1,2,0)^T,\qquad p_2=(1,0,0)^T,\qquad p_3=(-1,2,4)^T.
\]
于是
\[
P=\begin{pmatrix}
-1&1&-1\\
2&0&2\\
0&0&4
\end{pmatrix},
\qquad |P|=-8\ne0,
\]
并且直接验证
\[
P^{-1}AP=B.
\]}
\answer{可取
\[
P=\begin{pmatrix}
-1&1&-1\\
2&0&2\\
0&0&4
\end{pmatrix}.
\]}
\end{solutionblock}

\subsection{5.5 实对称矩阵相似对角化}

\begin{solutionblock}
\textbf{5.5 例 1}

\question{设二维向量 \(\alpha,\beta\) 满足 \(\|\alpha+\beta\|=\|\alpha-\beta\|\)，令 \(A=(\alpha,\beta)\)，判断 \(A\) 的列向量关系及 \(A^TA\) 的性质。}



\analysis{由
\[
\|\alpha+\beta\|=\|\alpha-\beta\|
\]
两边平方：
\[
(\alpha+\beta)^T(\alpha+\beta)
=(\alpha-\beta)^T(\alpha-\beta).
\]
展开得
\[
\alpha^T\alpha+2\alpha^T\beta+\beta^T\beta
=\alpha^T\alpha-2\alpha^T\beta+\beta^T\beta,
\]
故
\[
\alpha^T\beta=0.
\]
又 \(\alpha,\beta\) 均为单位列向量，所以矩阵
\[
A=(\alpha,\beta)
\]
的两列是标准正交向量，满足
\[
A^TA=E.
\]
因此 \(A\) 是正交矩阵。}
\answer{选 C。}
\end{solutionblock}

\begin{solutionblock}
\textbf{5.5 例 2}

\question{设 \(A\) 为四阶实对称矩阵，满足 \(A^2=-A\)，且 \(r(A)=3\)，判断 \(A\) 相似于哪个对角矩阵。}



\analysis{\(A\) 为四阶实对称矩阵，故正交相似于对角矩阵。设其特征值为 \(\lambda\)。由
\[
A^2=-A
\]
知
\[
\lambda^2=-\lambda,\qquad \lambda(\lambda+1)=0.
\]
所以特征值只能是 \(0\) 或 \(-1\)。又
\[
r(A)=3,
\]
说明非零特征值有 \(3\) 个，故 \(-1\) 出现三次，\(0\) 出现一次。}
\answer{选 D，即 \(A\sim \operatorname{diag}(-1,-1,-1,0)\)。}
\end{solutionblock}

\begin{solutionblock}
\textbf{5.5 例 3}

\question{设实对称矩阵 \(A\) 满足 \(A^3+A^2+A-3E=0\)，证明 \(A=E\)。}



\analysis{由于 \(A\) 为实对称矩阵，存在正交矩阵 \(Q\)，使
\[
Q^TAQ=\Lambda=\operatorname{diag}(\lambda_1,\lambda_2,\ldots,\lambda_n),
\]
其中 \(\lambda_i\) 均为实数。题设
\[
A^3+A^2+A-3E=0,
\]
所以每个特征值都满足
\[
\lambda^3+\lambda^2+\lambda-3=0.
\]
因式分解：
\[
\lambda^3+\lambda^2+\lambda-3
=(\lambda-1)(\lambda^2+2\lambda+3).
\]
二次因子 \(\lambda^2+2\lambda+3\) 的判别式为
\[
4-12=-8<0,
\]
没有实根。故实特征值只能是
\[
\lambda=1.
\]
于是 \(\Lambda=E\)，从而
\[
A=QEQ^T=E.
\]}
\answer{已证 \(A=E\)。}
\examnote{实对称矩阵的核心优势是“特征值全为实数、可正交对角化”，所以多项式方程可以转化为实数方程。}
\end{solutionblock}

\begin{solutionblock}
\textbf{5.5 例 4}

\question{设三阶实对称矩阵 \(A\) 的特征值为 \(-1,1,2\)。已知 \(\lambda_1=-1\) 的特征向量为 \(\alpha_1=(1,2,1)^T\)，\(\lambda_2=1\) 的特征向量为 \(\alpha_2=(-1,0,1)^T\)，求 \(\lambda_3=2\) 的特征向量。}



\analysis{实对称矩阵不同特征值对应的特征向量正交。已知
\[
\alpha_1=(1,2,1)^T,\qquad \alpha_2=(-1,0,1)^T
\]
分别属于 \(\lambda_1=-1,\lambda_2=1\)，所求 \(\lambda_3=2\) 的特征向量 \(\alpha_3\) 必须同时满足
\[
\alpha_3^T\alpha_1=0,\qquad \alpha_3^T\alpha_2=0.
\]
取叉乘：
\[
\alpha_1\times\alpha_2=(2,-2,2)^T.
\]
故可取
\[
\alpha_3=(1,-1,1)^T.
\]}
\answer{\[
\alpha_3=k(1,-1,1)^T,\qquad k\ne0.
\]}
\end{solutionblock}

\begin{solutionblock}
\textbf{5.5 例 5}

\question{设三阶实对称矩阵 \(A\) 每行元素之和均为 \(3\)，且 \(\lambda_1=\lambda_2=1\) 为二重特征值，求 \(A\) 属于 \(\lambda=1\) 的特征向量空间。}



\analysis{设 \(\mathbf{1}=(1,1,1)^T\)。每行元素之和均为 \(3\)，等价于
\[
A\mathbf{1}=3\mathbf{1}.
\]
因此 \(3\) 是 \(A\) 的一个特征值，对应特征向量为 \((1,1,1)^T\)。

题设还有二重特征值 \(\lambda_1=\lambda_2=1\)。由于 \(A\) 为实对称矩阵，不同特征值对应的特征向量正交，所以属于 \(\lambda=1\) 的特征向量必须与 \((1,1,1)^T\) 正交，即满足
\[
x_1+x_2+x_3=0.
\]
该平面的一组基为
\[
(1,-1,0)^T,\qquad (1,0,-1)^T.
\]}
\answer{\[
\lambda=3:\ \xi=k(1,1,1)^T,\ k\ne0;
\]
\[
\lambda=1:\ \xi=s(1,-1,0)^T+t(1,0,-1)^T,\ (s,t)\ne(0,0).
\]}
\end{solutionblock}

\begin{solutionblock}
\textbf{5.5 例 6}

\question{设 \(A=\begin{pmatrix}1&-2&2\\-2&-2&4\\2&4&-2\end{pmatrix}\)，求正交矩阵 \(Q\)，使 \(Q^{-1}AQ=Q^TAQ=\Lambda\)。}



\analysis{设
\[
A=\begin{pmatrix}
1&-2&2\\
-2&-2&4\\
2&4&-2
\end{pmatrix}.
\]
计算特征多项式：
\[
|\lambda E-A|=(\lambda-2)^2(\lambda+7).
\]
对应 \(\lambda=2\) 的特征子空间可取两个正交向量
\[
\xi_1=(-2,1,0)^T,\qquad \xi_2=(2,4,5)^T.
\]
对应 \(\lambda=-7\) 的特征向量可取
\[
\xi_3=(-1,-2,2)^T.
\]
归一化：
\[
q_1=\frac1{\sqrt5}(-2,1,0)^T,\qquad
q_2=\frac1{3\sqrt5}(2,4,5)^T,\qquad
q_3=\frac13(-1,-2,2)^T.
\]
令
\[
Q=(q_1,q_2,q_3).
\]
则 \(Q\) 为正交矩阵，且
\[
Q^TAQ=\operatorname{diag}(2,2,-7).
\]}
\answer{\[
Q=\begin{pmatrix}
-\frac2{\sqrt5}&\frac2{3\sqrt5}&-\frac13\\
\frac1{\sqrt5}&\frac4{3\sqrt5}&-\frac23\\
0&\frac5{3\sqrt5}&\frac23
\end{pmatrix},
\qquad
\Lambda=\operatorname{diag}(2,2,-7).
\]}
\end{solutionblock}

\begin{solutionblock}
\textbf{5.5 例 7}

\question{设 \(A\) 为三阶实对称矩阵，各行元素之和均为 \(3\)，且 \(\alpha_1=(1,-2,1)^T,\alpha_2=(0,1,-1)^T\) 是 \(Ax=0\) 的基础解系，求 \(A\) 及 \((A-\frac32E)^6\)。}



\analysis{设
\[
\alpha_1=(1,-2,1)^T,\qquad \alpha_2=(0,1,-1)^T.
\]
它们是 \(Ax=0\) 的两个线性无关解，所以
\[
A\alpha_1=0,\qquad A\alpha_2=0.
\]
又 \(A\) 为三阶实对称矩阵且各行元素之和均为 \(3\)，故
\[
A(1,1,1)^T=3(1,1,1)^T.
\]
注意
\[
\alpha_1^T(1,1,1)=0,\qquad \alpha_2^T(1,1,1)=0,
\]
故这三个向量构成一组正交特征向量。于是 \(A\) 在这组特征向量下的特征值为 \(0,0,3\)。这说明 \(A\) 是到 \((1,1,1)^T\) 方向的三倍投影：
\[
A=3\cdot \frac{(1,1,1)^T(1,1,1)}{(1,1,1)(1,1,1)^T}
=3\cdot\frac13
\begin{pmatrix}
1&1&1\\
1&1&1\\
1&1&1
\end{pmatrix}
=\begin{pmatrix}
1&1&1\\
1&1&1\\
1&1&1
\end{pmatrix}.
\]
矩阵 \(A-\frac32E\) 的特征值为
\[
-\frac32,\ -\frac32,\ \frac32.
\]
六次方后均变为
\[
\left(\frac32\right)^6=\frac{729}{64}.
\]
故
\[
\left(A-\frac32E\right)^6=\frac{729}{64}E.
\]}
\answer{\[
A=\begin{pmatrix}
1&1&1\\
1&1&1\\
1&1&1
\end{pmatrix},
\qquad
\left(A-\frac32E\right)^6=\frac{729}{64}E.
\]}
\end{solutionblock}

\chapter{二次型}

\section{逐题解析}

\subsection{6.1 二次型的定义及其表示}

\begin{solutionblock}
\textbf{6.1 例 1}

\question{求二次型 \(f=x_1^2+3x_2^2-x_3^2+2x_1x_2+2x_1x_3+6x_2x_3\) 的矩阵及二次型的秩。}



\analysis{二次型
\[
f=x_1^2+3x_2^2-x_3^2+2x_1x_2+2x_1x_3+6x_2x_3
\]
写成 \(f=x^TAx\) 时，矩阵 \(A\) 必须为对称矩阵。平方项系数直接放在主对角线，交叉项 \(2a_{ij}x_ix_j\) 的系数等于题中交叉项系数。

因此
\[
a_{11}=1,\quad a_{22}=3,\quad a_{33}=-1,
\]
\[
2a_{12}=2,\quad 2a_{13}=2,\quad 2a_{23}=6.
\]
故
\[
A=\begin{pmatrix}
1&1&1\\
1&3&3\\
1&3&-1
\end{pmatrix}.
\]
矩阵的秩为
\[
r(A)=3.
\]}
\answer{\[
A=\begin{pmatrix}
1&1&1\\
1&3&3\\
1&3&-1
\end{pmatrix},\qquad r(f)=r(A)=3.
\]}
\end{solutionblock}

\begin{solutionblock}
\textbf{6.1 例 2}

\question{设 \(f(x_1,x_2,x_3)=\displaystyle\sum_{i=1}^3\sum_{j=1}^3ijx_ix_j\)，求该二次型对应的矩阵 \(A\)。}



\analysis{题中
\[
f(x_1,x_2,x_3)=\sum_{i=1}^3\sum_{j=1}^3 ijx_ix_j.
\]
所以二次型矩阵的第 \(i\) 行第 \(j\) 列元素为
\[
a_{ij}=ij.
\]
列出矩阵：
\[
A=\begin{pmatrix}
1\cdot1&1\cdot2&1\cdot3\\
2\cdot1&2\cdot2&2\cdot3\\
3\cdot1&3\cdot2&3\cdot3
\end{pmatrix}
=\begin{pmatrix}
1&2&3\\
2&4&6\\
3&6&9
\end{pmatrix}.
\]}
\answer{\[
A=\begin{pmatrix}
1&2&3\\
2&4&6\\
3&6&9
\end{pmatrix}.
\]}
\end{solutionblock}

\begin{solutionblock}
\textbf{6.1 例 3}

\question{设 \(f(x_1,x_2,x_3)=x^T\begin{pmatrix}1&4&0\\0&4&0\\0&0&0\end{pmatrix}x\)，求该二次型的秩。}



\analysis{给出的矩阵
\[
M=\begin{pmatrix}
1&4&0\\
0&4&0\\
0&0&0
\end{pmatrix}
\]
不是对称矩阵。二次型
\[
x^TMx
\]
只由矩阵的对称部分决定，即
\[
x^TMx=x^T\left(\frac{M+M^T}{2}\right)x.
\]
因此真正对应的二次型矩阵为
\[
A=\frac{M+M^T}{2}
=\begin{pmatrix}
1&2&0\\
2&4&0\\
0&0&0
\end{pmatrix}.
\]
前两行成比例，且第一行非零，故
\[
r(A)=1.
\]}
\answer{\(1\)}
\examnote{二次型矩阵默认取对称矩阵。若题面给了非对称矩阵，秩要算其对称部分的秩，而不是原矩阵的秩。}
\end{solutionblock}

\begin{solutionblock}
\textbf{6.1 例 4}

\question{设 \(f=(x_1+x_2+x_3)^2+(x_2-x_3)^2+(x_1+2x_2+x_3)^2\)，求二次型矩阵 \(A\)。}



\analysis{令
\[
l_1=x_1+x_2+x_3,\quad l_2=x_2-x_3,\quad l_3=x_1+2x_2+x_3.
\]
则
\[
f=l_1^2+l_2^2+l_3^2.
\]
把线性形式写成矩阵
\[
L=\begin{pmatrix}
1&1&1\\
0&1&-1\\
1&2&1
\end{pmatrix},
\qquad
\begin{pmatrix}l_1\\l_2\\l_3\end{pmatrix}=Lx.
\]
于是
\[
f=(Lx)^T(Lx)=x^T(L^TL)x,
\]
故二次型矩阵为
\[
A=L^TL
=\begin{pmatrix}
2&3&2\\
3&6&2\\
2&2&3
\end{pmatrix}.
\]}
\answer{\[
A=\begin{pmatrix}
2&3&2\\
3&6&2\\
2&2&3
\end{pmatrix}.
\]}
\end{solutionblock}

\subsection{6.2 可逆线性变换与矩阵合同}

\begin{solutionblock}
\textbf{6.2 例 1}

\question{设 \(\alpha_1=(1,2)^T,\alpha_2=(a,1)^T,x=(x_1,x_2)^T\)，令 \(f=\sum_{i=1}^2(\alpha_i,x)^2\)。若 \(f\) 与 \(g=y_1^2+y_2^2+2y_1y_2\) 合同，求 \(a\)。}



\analysis{由
\[
\alpha_1=(1,2)^T,\qquad \alpha_2=(a,1)^T
\]
得
\[
f=(\alpha_1^Tx)^2+(\alpha_2^Tx)^2=x^TAx,
\]
其中
\[
A=\alpha_1\alpha_1^T+\alpha_2\alpha_2^T
=\begin{pmatrix}
1+a^2&2+a\\
2+a&5
\end{pmatrix}.
\]
题中 \(f\) 经可逆线性变换化为
\[
g=y_1^2+y_2^2+2y_1y_2,
\]
其矩阵为
\[
B=\begin{pmatrix}1&1\\1&1\end{pmatrix},
\]
秩为 \(1\)。合同变换保持矩阵秩，因此
\[
r(A)=r(B)=1.
\]
对 \(A\) 计算行列式：
\[
|A|=5(1+a^2)-(2+a)^2=(2a-1)^2.
\]
令 \(|A|=0\)，得
\[
a=\frac12.
\]}
\answer{\(\dfrac12\)}
\end{solutionblock}

\begin{solutionblock}
\textbf{6.2 例 2}

\question{设 \(A=\begin{pmatrix}1&2\\2&1\end{pmatrix}\)，\(B=\begin{pmatrix}1&4\\1&1\end{pmatrix}\)。判断是否存在可逆矩阵 \(C\) 使 \(C^TAC=B\)，以及是否存在正交矩阵 \(Q\) 使 \(Q^TAQ=B\)。}



\analysis{给定
\[
A=\begin{pmatrix}1&2\\2&1\end{pmatrix},\qquad
B=\begin{pmatrix}1&4\\1&1\end{pmatrix}.
\]

第 (1) 个说法：若存在可逆矩阵 \(C\)，使 \(C^TAC=B\)，则 \(B\) 必须是对称矩阵，因为 \(A=A^T\) 时
\[
(C^TAC)^T=C^TA^TC=C^TAC.
\]
但题中 \(B\) 不是对称矩阵，所以第 (1) 个说法错误。

第 (2) 个说法：若存在正交矩阵 \(Q\)，使 \(Q^TAQ=B\)，则 \(B\) 不仅应为对称矩阵，而且与 \(A\) 正交相似，特征值也应相同。题中 \(B\) 非对称，故第 (2) 个说法也错误。}
\answer{两个说法均错误。}
\end{solutionblock}

\subsection{6.3 二次型的标准形和规范形}

\begin{solutionblock}
\textbf{6.3 例 1}

\question{用正交变换 \(x=Qy\) 化二次型 \(f=x_1^2+x_2^2+x_3^2+4x_1x_2+4x_1x_3+4x_2x_3\) 为标准形。}



\analysis{二次型
\[
f=x_1^2+x_2^2+x_3^2+4x_1x_2+4x_1x_3+4x_2x_3
\]
的矩阵为
\[
A=\begin{pmatrix}
1&2&2\\
2&1&2\\
2&2&1
\end{pmatrix}.
\]
该矩阵的特征值为 \(5,-1,-1\)。对应 \(\lambda=5\) 的单位特征向量可取
\[
q_1=\frac1{\sqrt3}(1,1,1)^T.
\]
对应 \(\lambda=-1\) 的特征子空间为 \(x_1+x_2+x_3=0\)，可取一组正交单位向量
\[
q_2=\frac1{\sqrt2}(1,-1,0)^T,\qquad
q_3=\frac1{\sqrt6}(1,1,-2)^T.
\]
令
\[
Q=(q_1,q_2,q_3),
\]
则 \(Q\) 为正交矩阵，且
\[
Q^TAQ=\operatorname{diag}(5,-1,-1).
\]}
\answer{取
\[
Q=\begin{pmatrix}
\frac1{\sqrt3}&\frac1{\sqrt2}&\frac1{\sqrt6}\\
\frac1{\sqrt3}&-\frac1{\sqrt2}&\frac1{\sqrt6}\\
\frac1{\sqrt3}&0&-\frac2{\sqrt6}
\end{pmatrix},
\]
则在正交变换 \(x=Qy\) 下，
\[
f=5y_1^2-y_2^2-y_3^2.
\]}
\end{solutionblock}

\begin{solutionblock}
\textbf{6.3 例 2}

\question{用配方法化二次型 \(f=2x_1^2+3x_2^2+x_3^2+4x_1x_2-4x_1x_3-8x_2x_3\) 为标准形，并写出变换 \(x=Cy\)。}



\analysis{对
\[
f=2x_1^2+3x_2^2+x_3^2+4x_1x_2-4x_1x_3-8x_2x_3
\]
配方。先含 \(x_1\) 的项：
\[
2x_1^2+4x_1x_2-4x_1x_3
=2(x_1+x_2-x_3)^2-2(x_2-x_3)^2.
\]
故
\[
f=2(x_1+x_2-x_3)^2+\bigl[3x_2^2+x_3^2-8x_2x_3-2(x_2-x_3)^2\bigr].
\]
括号内化简为
\[
x_2^2-4x_2x_3-x_3^2=(x_2-2x_3)^2-5x_3^2.
\]
令
\[
y_1=x_1+x_2-x_3,\qquad y_2=x_2-2x_3,\qquad y_3=x_3.
\]
则
\[
f=2y_1^2+y_2^2-5y_3^2.
\]
由上式反解
\[
x_3=y_3,\qquad x_2=y_2+2y_3,\qquad x_1=y_1-y_2-y_3.
\]
所以
\[
x=Cy,\qquad
C=\begin{pmatrix}
1&-1&-1\\
0&1&2\\
0&0&1
\end{pmatrix}.
\]}
\answer{\[
f=2y_1^2+y_2^2-5y_3^2,\qquad
C=\begin{pmatrix}
1&-1&-1\\
0&1&2\\
0&0&1
\end{pmatrix}.
\]}
\end{solutionblock}

\begin{solutionblock}
\textbf{6.3 例 3}

\question{用配方法化二次型 \(f=2x_1x_2+4x_1x_3\) 为标准形，并写出变换 \(x=Cy\)。}



\analysis{题中
\[
f=2x_1x_2+4x_1x_3=2x_1(x_2+2x_3).
\]
令
\[
y_1=x_1+x_2+2x_3,\qquad
y_2=x_1-x_2-2x_3,\qquad
y_3=x_3.
\]
则
\[
y_1^2-y_2^2
=(x_1+x_2+2x_3)^2-(x_1-x_2-2x_3)^2
=4x_1(x_2+2x_3).
\]
因此
\[
f=\frac12y_1^2-\frac12y_2^2.
\]
反解得
\[
x_1=\frac12y_1+\frac12y_2,\qquad
x_2=\frac12y_1-\frac12y_2-2y_3,\qquad
x_3=y_3.
\]
所以
\[
x=Cy,\qquad
C=\begin{pmatrix}
\frac12&\frac12&0\\
\frac12&-\frac12&-2\\
0&0&1
\end{pmatrix}.
\]}
\answer{\[
f=\frac12y_1^2-\frac12y_2^2,\qquad
C=\begin{pmatrix}
\frac12&\frac12&0\\
\frac12&-\frac12&-2\\
0&0&1
\end{pmatrix}.
\]}
\end{solutionblock}

\begin{solutionblock}
\textbf{6.3 例 4}

\question{求二次型 \(f=x_1x_2+x_1x_3+2x_2x_3\) 的标准形，并给出相应可逆线性变换。}



\analysis{二次型
\[
f=x_1x_2+x_1x_3+2x_2x_3
\]
的矩阵为
\[
A=\begin{pmatrix}
0&\frac12&\frac12\\
\frac12&0&1\\
\frac12&1&0
\end{pmatrix}.
\]
作可逆变换
\[
x=Cy,\qquad
C=\begin{pmatrix}
1&1&-\sqrt2\\
1&-1&-\frac1{\sqrt2}\\
0&0&\frac1{\sqrt2}
\end{pmatrix}.
\]
直接计算得
\[
C^TAC=\operatorname{diag}(1,-1,-1).
\]
因此规范形为
\[
z_1^2-z_2^2-z_3^2.
\]
若不缩放第三个变量，也可先得到标准形
\[
y_1^2-y_2^2-2y_3^2.
\]}
\answer{规范形为
\[
z_1^2-z_2^2-z_3^2,
\]
可取上述 \(C\) 作为坐标变换矩阵。}
\end{solutionblock}

\begin{solutionblock}
\textbf{6.3 例 5}

\question{求二次型 \(f=x_1^2+3x_2^2+x_3^2+2x_1x_2+2x_1x_3+2x_2x_3\) 的规范形。}



\analysis{二次型
\[
f=x_1^2+3x_2^2+x_3^2+2x_1x_2+2x_1x_3+2x_2x_3
\]
对应矩阵
\[
A=\begin{pmatrix}
1&1&1\\
1&3&1\\
1&1&1
\end{pmatrix}.
\]
计算其特征值为
\[
4,\ 1,\ 0.
\]
正惯性指数等于正特征值个数，为 \(2\)；负惯性指数等于负特征值个数，为 \(0\)。}
\answer{正惯性指数为 \(2\)，负惯性指数为 \(0\)。}
\end{solutionblock}

\begin{solutionblock}
\textbf{6.3 例 6}

\question{求二次型 \(f=-2x_1x_2-2x_1x_3+6x_2x_3\) 的正惯性指数，并在选项 \(3,2,1,0\) 中作出判断。}



\analysis{二次型
\[
f=-2x_1x_2-2x_1x_3+6x_2x_3
\]
对应矩阵为
\[
A=\begin{pmatrix}
0&-1&-1\\
-1&0&3\\
-1&3&0
\end{pmatrix}.
\]
其特征值为
\[
-3,\quad \frac{3-\sqrt{17}}2,\quad \frac{3+\sqrt{17}}2.
\]
其中
\[
-3<0,\qquad \frac{3-\sqrt{17}}2<0,\qquad \frac{3+\sqrt{17}}2>0.
\]
所以正惯性指数为 \(1\)。}
\answer{选 C，正惯性指数为 \(1\)。}
\end{solutionblock}

\begin{solutionblock}
\textbf{6.3 例 7}

\question{求二次型 \(f=(x_1+x_2)^2+(x_2-x_3)^2+(x_3+x_1)^2\) 的正惯性指数，并在选项 \(3,2,1,0\) 中作出判断。}



\analysis{因为
\[
f=(x_1+x_2)^2+(x_2-x_3)^2+(x_3+x_1)^2
\]
是三个平方和，所以负惯性指数为 \(0\)。展开可得矩阵
\[
A=\begin{pmatrix}
2&1&1\\
1&2&-1\\
1&-1&2
\end{pmatrix}.
\]
其特征值为 \(3,3,0\)，因此正惯性指数为 \(2\)。}
\answer{选 B，正惯性指数为 \(2\)。}
\end{solutionblock}

\subsection{6.4 正定二次型}

\begin{solutionblock}
\textbf{6.4 例 1}

\question{判断二次型 \(f=2x_1^2+4x_2^2+5x_3^2-4x_1x_3\) 是否正定。}



\analysis{二次型
\[
f=2x_1^2+4x_2^2+5x_3^2-4x_1x_3
\]
对应矩阵为
\[
A=\begin{pmatrix}
2&0&-2\\
0&4&0\\
-2&0&5
\end{pmatrix}.
\]
用顺序主子式判别：
\[
\Delta_1=2>0,
\]
\[
\Delta_2=\begin{vmatrix}2&0\\0&4\end{vmatrix}=8>0,
\]
\[
\Delta_3=|A|=24>0.
\]
三个顺序主子式全为正，故 \(A\) 正定，二次型 \(f\) 正定。}
\answer{\(f\) 是正定二次型。}
\end{solutionblock}

\begin{solutionblock}
\textbf{6.4 例 2}

\question{求使二次型 \(f=2x_1^2+x_2^2+x_3^2+2x_1x_2+tx_2x_3\) 正定的参数 \(t\) 的取值范围。}



\analysis{二次型
\[
f=2x_1^2+x_2^2+x_3^2+2x_1x_2+tx_2x_3
\]
对应矩阵为
\[
A=\begin{pmatrix}
2&1&0\\
1&1&\frac t2\\
0&\frac t2&1
\end{pmatrix}.
\]
顺序主子式为
\[
\Delta_1=2>0,
\]
\[
\Delta_2=\begin{vmatrix}2&1\\1&1\end{vmatrix}=1>0,
\]
\[
\Delta_3=|A|=1-\frac{t^2}{2}.
\]
正定要求 \(\Delta_3>0\)，所以
\[
1-\frac{t^2}{2}>0
\quad\Longleftrightarrow\quad
t^2<2.
\]}
\answer{\[
-\sqrt2<t<\sqrt2.
\]}
\end{solutionblock}

\begin{solutionblock}
\textbf{6.4 例 3}

\question{设 \(A=\begin{pmatrix}1&0&1\\0&2&0\\1&0&1\end{pmatrix}\)，\(B=(A+kE)^2\)。求使 \(B\) 正定的 \(k\) 的取值范围。}



\analysis{矩阵
\[
A=\begin{pmatrix}
1&0&1\\
0&2&0\\
1&0&1
\end{pmatrix}
\]
为实对称矩阵，特征值为
\[
0,\ 2,\ 2.
\]
矩阵
\[
B=(A+kE)^2
\]
的特征值为
\[
(0+k)^2,\quad (2+k)^2,\quad (2+k)^2.
\]
\(B\) 正定的充要条件是这三个特征值均大于 \(0\)，即
\[
k\ne0,\qquad k\ne-2.
\]}
\answer{\[
k\in\mathbb{R}\setminus\{0,-2\}.
\]}
\end{solutionblock}

\begin{solutionblock}
\textbf{6.4 例 4}

\question{设 \(f=(ax_1+2x_2-3x_3)^2+(x_2-2x_3)^2+(x_1+ax_2-x_3)^2\)，判断 \(f\) 正定时 \(a\) 应满足的条件。}



\analysis{题中
\[
f=(ax_1+2x_2-3x_3)^2+(x_2-2x_3)^2+(x_1+ax_2-x_3)^2.
\]
这是三个线性形式平方和。设
\[
R=\begin{pmatrix}
a&2&-3\\
0&1&-2\\
1&a&-1
\end{pmatrix}.
\]
则
\[
f=\|Rx\|^2=x^TR^TRx.
\]
该二次型正定的充要条件是 \(Rx=0\) 只有零解，即 \(R\) 可逆。计算
\[
|R|=(a-1)(2a+1).
\]
故正定要求
\[
(a-1)(2a+1)\ne0,
\]
即
\[
a\ne1,\qquad a\ne-\frac12.
\]}
\answer{选 C，\(a\ne1\) 且 \(a\ne-\dfrac12\)。}
\end{solutionblock}

\begin{solutionblock}
\textbf{6.4 例 5}

\question{设 \(f=x^TAx\)，\(A^T=A\)。判断 \(f\) 正定的充要条件。}



\analysis{对称矩阵 \(A\) 对应二次型 \(f=x^TAx\)。若 \(f\) 正定，则 \(A\) 正定，必可逆，因此 \(A^{-1}\) 也正定；反之若 \(A^{-1}\) 正定，则其逆矩阵 \(A\) 也正定，故 \(f\) 正定。因此 D 是充要条件。

选项 A 只有 \(|A|>0\)，不足以保证所有特征值均为正。选项 B “负惯性指数为 \(0\)”只说明没有负平方项，但可能存在零平方项，属于半正定。选项 C “秩为 \(n\)”只说明非退化，也可能是不定型。}
\answer{选 D。}
\end{solutionblock}

\begin{solutionblock}
\textbf{6.4 例 6}

\question{设 \(A\) 为 \(m\times n\) 矩阵，\(E\) 为 \(n\) 阶单位矩阵，\(B=\lambda E+A^TA\)。证明当 \(\lambda>0\) 时 \(B\) 正定。}



\analysis{设 \(x\ne0\)。由
\[
B=\lambda E+A^TA
\]
得
\[
x^TBx=\lambda x^Tx+x^TA^TAx
=\lambda\|x\|^2+\|Ax\|^2.
\]
当 \(\lambda>0\) 时，
\[
\lambda\|x\|^2>0,\qquad \|Ax\|^2\ge0,
\]
所以
\[
x^TBx>0.
\]
故 \(B\) 为正定矩阵。}
\answer{已证：当 \(\lambda>0\) 时，\(B\) 为正定矩阵。}
\end{solutionblock}

\begin{solutionblock}
\textbf{6.4 例 7}

\question{设 \(A\) 为 \(n\) 阶正定矩阵，证明存在 \(n\) 阶可逆矩阵 \(D\) 使 \(A=D^TD\)，并证明存在 \(n\) 阶正定矩阵 \(B\) 使 \(A=B^2\)。}



\analysis{第 (1) 问：若 \(A\) 正定，则存在正交矩阵 \(Q\)，使
\[
Q^TAQ=\Lambda=\operatorname{diag}(\lambda_1,\ldots,\lambda_n),
\]
且
\[
\lambda_i>0.
\]
令
\[
\Lambda^{1/2}=\operatorname{diag}(\sqrt{\lambda_1},\ldots,\sqrt{\lambda_n}),
\qquad
D=\Lambda^{1/2}Q^T.
\]
则 \(D\) 可逆，且
\[
D^TD=Q\Lambda^{1/2}\Lambda^{1/2}Q^T=Q\Lambda Q^T=A.
\]

第 (2) 问：令
\[
B=Q\Lambda^{1/2}Q^T.
\]
则 \(B\) 为实对称矩阵，且特征值均为 \(\sqrt{\lambda_i}>0\)，所以 \(B\) 正定。并且
\[
B^2=Q\Lambda^{1/2}Q^TQ\Lambda^{1/2}Q^T
=Q\Lambda Q^T=A.
\]}
\answer{两个结论均成立。}
\examnote{这是正定矩阵的平方根分解：\(A=D^TD\) 是 Cholesky 型分解，\(A=B^2\) 是对称正定平方根。}
\end{solutionblock}

\subsection{6.5 矩阵等价、相似、合同}

\begin{solutionblock}
\textbf{6.5 例 1}

\question{设 \(A=\begin{pmatrix}2&-1&-1\\-1&2&-1\\-1&-1&2\end{pmatrix}\)，\(B=\begin{pmatrix}1&0&0\\0&1&0\\0&0&0\end{pmatrix}\)，判断 \(A\) 与 \(B\) 的等价、相似、合同关系。}



\analysis{设
\[
A=\begin{pmatrix}
2&-1&-1\\
-1&2&-1\\
-1&-1&2
\end{pmatrix},\qquad
B=\begin{pmatrix}
1&0&0\\
0&1&0\\
0&0&0
\end{pmatrix}.
\]
矩阵 \(A\) 的特征值为
\[
3,\ 3,\ 0,
\]
因此 \(r(A)=2\)。而 \(r(B)=2\)，所以 \(A\) 与 \(B\) 等价。

相似要求特征值完全相同。\(B\) 的特征值为
\[
1,\ 1,\ 0,
\]
与 \(A\) 的特征值不同，所以二者不相似。

合同分类看惯性指数。\(A\) 的正惯性指数为 \(2\)、负惯性指数为 \(0\)；\(B\) 的正惯性指数为 \(2\)、负惯性指数为 \(0\)。因此 \(A\) 与 \(B\) 合同。}
\answer{选 B，合同但不相似。}
\end{solutionblock}

\begin{solutionblock}
\textbf{6.5 例 2}

\question{设 \(A\) 为 \(n\) 阶实对称矩阵，将 \(A\) 的第 \(i\) 行与第 \(j\) 行交换，再将第 \(i\) 列与第 \(j\) 列交换得到 \(B\)，判断 \(A,B\) 的等价、相似、合同关系。}



\analysis{交换 \(A\) 的第 \(i\) 列和第 \(j\) 列，可写成右乘一个置换矩阵 \(P\)；再交换第 \(i\) 行和第 \(j\) 行，可写成左乘同一个置换矩阵 \(P\)。因此
\[
B=PAP.
\]
置换矩阵满足
\[
P^T=P,\qquad P^{-1}=P.
\]
于是
\[
B=P^TAP,
\]
所以 \(A\) 与 \(B\) 合同；同时
\[
B=P^{-1}AP,
\]
所以 \(A\) 与 \(B\) 相似。由于初等行列变换不改变等价关系，且相似或合同都蕴含等价矩阵同秩，因此 \(A\) 与 \(B\) 也等价。}
\answer{选 D，等价、相似、合同。}
\end{solutionblock}


\end{document}
